\documentclass[11pt]{article}

\usepackage[margin=1.1in]{geometry}

\usepackage{amssymb,mathtools,natbib}

% Anchored formal environments for causalsmith papers.
% First mandatory argument = obj_id (machine anchor joining the paper to the
% Lean crosswalk; invisible in the rendered PDF). Optional argument = name.
% Each environment also defines \label{obj:<obj_id>} so prose can use
% \cref{obj:T-1}; cleveref derives the printed kind and number from the target.
\usepackage{amsmath}
\usepackage{mathrsfs}
\usepackage{amsthm}
\usepackage{booktabs}
% Long displays may break across pages; let TeX stretch a little harder before an
% overfull \hbox so wide formulas/tables are less likely to run off the page.
\allowdisplaybreaks
\setlength{\emergencystretch}{3em}
\newtheorem{theorem}{Theorem}
\newtheorem{lemma}{Lemma}
\newtheorem{assumption}{Assumption}
\newtheorem{definition}{Definition}
\newtheorem{proposition}{Proposition}
\newtheorem{remark}{Remark}
\newtheorem{citedresult}{Cited result}
% The amsthm counter is `algorithmbox` (not `algorithm`) so a later `\usepackage{algorithm}` cannot clash.
\newcounter{algorithmbox}
\renewcommand{\thealgorithmbox}{\arabic{algorithmbox}}
\NewDocumentEnvironment{theoremv}{m o s}{\begin{theorem}\label{obj:#1}{\normalfont[\texttt{\detokenize{#1}}]\IfValueT{#2}{\ (#2)}\IfBooleanT{#3}{\verificationfootnotemark}\ }}{\end{theorem}}
\NewDocumentEnvironment{lemmav}{m o s}{\begin{lemma}\label{obj:#1}{\normalfont[\texttt{\detokenize{#1}}]\IfValueT{#2}{\ (#2)}\IfBooleanT{#3}{\verificationfootnotemark}\ }}{\end{lemma}}
\NewDocumentEnvironment{assumptionv}{m o s}{\begin{assumption}\label{obj:#1}{\normalfont[\texttt{\detokenize{#1}}]\IfValueT{#2}{\ (#2)}\IfBooleanT{#3}{\verificationfootnotemark}\ }}{\end{assumption}}
\NewDocumentEnvironment{definitionv}{m o s}{\begin{definition}\label{obj:#1}{\normalfont[\texttt{\detokenize{#1}}]\IfValueT{#2}{\ (#2)}\IfBooleanT{#3}{\verificationfootnotemark}\ }}{\end{definition}}
% A `propositionv` is a proved result tiered as a Proposition (theorem-grade, machine-verified).
\NewDocumentEnvironment{propositionv}{m o s}{\begin{proposition}\label{obj:#1}{\normalfont[\texttt{\detokenize{#1}}]\IfValueT{#2}{\ (#2)}\IfBooleanT{#3}{\verificationfootnotemark}\ }}{\end{proposition}}
% A `remarkv` holds NON-load-bearing interpretive / open-question content (e.g. a causalsmith `oeq:`
% open question). It is neither proved nor assumed; no theorem may depend on it.
\NewDocumentEnvironment{remarkv}{m o s}{\begin{remark}\label{obj:#1}{\normalfont[\texttt{\detokenize{#1}}]\IfValueT{#2}{\ (#2)}\IfBooleanT{#3}{\verificationfootnotemark}\ }}{\end{remark}}
% An `algorithmv` is a DEFINITION-kind object presented as a boxed, numbered-step procedure (an
% estimator / selection rule built in stages). Same anchor contract as `definitionv`; its body is
% ordinary LaTeX whose steps are `\begin{enumerate}` items, so tex2html needs no extra machinery.
% Presented in the CONVENTIONAL algorithm style readers expect from `algorithm`+`algorithmic`:
% a rule above, an "Algorithm N (Title)" caption line, a rule under the caption, the numbered
% steps, and a closing rule. It is NOT a float — the anchor contract requires the object to stay
% where the outline places it, and floats drift. The obj-id tag is suppressed here (it is
% bookkeeping, and a caption line carrying it reads as debug output in a finished paper).
\NewDocumentEnvironment{algorithmv}{m o s}%
  {\par\addvspace{\medskipamount}\noindent\rule{\linewidth}{0.8pt}\par\nobreak\vspace{2pt}%
   \refstepcounter{algorithmbox}\label{obj:#1}%
   \noindent\textbf{Algorithm \thealgorithmbox}\IfValueT{#2}{ \textnormal{(#2)}}%
   \IfBooleanT{#3}{\verificationfootnotemark}\par\nobreak\vspace{2pt}%
   \noindent\rule{\linewidth}{0.4pt}\par\nobreak\vspace{2pt}\normalfont}%
  {\par\nobreak\vspace{2pt}\noindent\rule{\linewidth}{0.8pt}\par\addvspace{\medskipamount}}
% A `citedv` is an IMPORTED external result the paper relies on but does NOT prove
% (a causalsmith `gate_class:"cited"` node). It is machine-checked against the cited
% source at F2.5, never proved here; the fixed provenance note makes that explicit
% so the environment can never be read as a contribution of this paper. The \citep
% attribution is supplied by the surrounding prose (P2), keyed to references.bib.
\NewDocumentEnvironment{citedv}{m o s}{\begin{citedresult}\label{obj:#1}{\normalfont[\texttt{\detokenize{#1}}]\IfValueT{#2}{\ (#2)}\IfBooleanT{#3}{\verificationfootnotemark}\ \textit{(imported result; machine-checked against the cited source, not proved here.)}\ }}{\end{citedresult}}
% \leanref{obj_id}{display text}: an inline first-mention link to a from-note object's
% verified Lean code. In the PDF it renders the display text plainly (the Lean link is a
% web-only affordance); tex2html turns it into a drawer-opening span.
\newcommand{\leanref}[2]{#2}
% For a theorem/lemma whose Lean declaration accepts a source-matched published proposition as an
% input, the starred anchored environment puts the mark in its heading; the generated text remains
% outside the frozen mathematical environment. Keep the combined macro for older paper bundles.
\newcommand{\verificationfootnotemark}{\footnotemark}
\newcommand{\verificationfootnotetext}[1]{\footnotetext{#1}}
\newcommand{\verificationfootnote}[1]{\verificationfootnotemark\verificationfootnotetext{#1}}
% xcolor before hyperref (hyperref would otherwise pull in plain `color` for colorlinks, and a
% later xcolor load clashes). The page-1 identity stamp at the end of this file needs a grey.
\usepackage{xcolor}
\usepackage{graphicx}
\usepackage{eso-pic}
% hyperref follows natbib and the environment macros; cleveref must follow hyperref.
\usepackage[colorlinks=true,linkcolor=blue,citecolor=blue,urlcolor=blue]{hyperref}
\usepackage[capitalise,nameinlink,noabbrev]{cleveref}
% Pin the names explicitly: labels are placed inside the underlying amsthm environments, and these
% declarations make the PDF contract independent of cleveref's theorem-name inference. House style:
% reference names always print with a capital first letter ("Theorem 1", "Definition 2"), in any
% sentence position — hence `capitalise` above and capitalized \crefname forms below.
\crefname{theorem}{Theorem}{Theorems}
\Crefname{theorem}{Theorem}{Theorems}
\crefname{lemma}{Lemma}{Lemmas}
\Crefname{lemma}{Lemma}{Lemmas}
\crefname{assumption}{Assumption}{Assumptions}
\Crefname{assumption}{Assumption}{Assumptions}
\crefname{definition}{Definition}{Definitions}
\Crefname{definition}{Definition}{Definitions}
\crefname{proposition}{Proposition}{Propositions}
\Crefname{proposition}{Proposition}{Propositions}
\crefname{remark}{Remark}{Remarks}
\Crefname{remark}{Remark}{Remarks}
\crefname{citedresult}{Cited result}{Cited results}
\Crefname{citedresult}{Cited result}{Cited results}
\crefname{algorithmbox}{Algorithm}{Algorithms}
\Crefname{algorithmbox}{Algorithm}{Algorithms}
% ---------------------------------------------------------------------------
% Page-1 identity stamp (arXiv style) and the pinned title-block date.
%
% Split of responsibility: THIS file owns the FORMAT (placement, font, colour, punctuation),
% the generated `paper_stamp.tex` owns only the VALUES (number+version, area, revised date,
% DOI, link target). P4 refreshes this template on every emit, so a layout fix reaches every
% bundle from a recompile; and because `paper_stamp.tex` is not one of the P2 assembly
% digest's authored inputs, a DOI that arrives after publication needs only that one small
% file rewritten plus a recompile — never a redraft, never a reassembly.
%
% Defaults first, so a bundle WITHOUT a stamp file compiles exactly as it always did:
% `\cspaperdate` falls back to `\today` and no stamp is drawn.
\providecommand*{\cspaperdate}{\today}  % the \date{} target
\providecommand*{\csstampid}{}          % e.g. CSWP-2026-008v3 (empty ⇒ no stamp at all)
\providecommand*{\csstamparea}{}        % e.g. Stat
\providecommand*{\csstampdate}{}        % e.g. 27 Aug 2026
\providecommand*{\csstampdoi}{}         % e.g. 10.5281/zenodo.123 (empty ⇒ DOI part omitted)
\providecommand*{\csstamplink}{}        % href target (DOI url, else the paper's page; may be empty)
\IfFileExists{paper_stamp.tex}{% GENERATED by CausalSmith P4 — paper identity stamp values. Do not hand-edit.
%
% Field order on the stamp (owned by paper_macros.tex):
%   <number>v<version> · doi:<doi> · [<Area>] <date>   — the DOI is never the last token, so a
%   text extractor cannot run it into whatever follows. Without a DOI that part is omitted.
% Format lives in paper_macros.tex; only the values are here, and this file is NOT one of
% the P2 assembly digest's authored inputs. A DOI arriving after publication therefore needs
% this file rewritten plus a `latexmk` recompile — never a redraft or a reassembly.
\renewcommand*{\csstampid}{CSWP-2026-033v4}
\renewcommand*{\csstamparea}{Stat}
\renewcommand*{\csstampdate}{2 Oct 2026}
\renewcommand*{\csstampdoi}{}
\renewcommand*{\csstamplink}{https://causalsmith.org/papers/stat_optvalue_vanishingoverlap_rate_v1}
\renewcommand*{\cspaperdate}{2 October 2026}
}{}
\makeatletter
% Field order is deliberate: the DOI is NEVER the last token on the line. A trailing identifier
% runs into whatever a text extractor emits next, which makes it unsafe to match mechanically
% (the Zenodo stamp matcher reads exactly this line); the date is a safe terminator.
%   <number>v<version> · doi:<doi> · [<Area>] <date>      — with a DOI
%   <number>v<version> · [<Area>] <date>                  — without
\newcommand{\cs@stampsep}{\,\textperiodcentered\,}
\newcommand{\cs@stamptext}{%
  \csstampid
  \ifx\csstampdoi\@empty\else\cs@stampsep doi:\csstampdoi\fi
  \ifx\csstamparea\@empty\else\cs@stampsep[\csstamparea]\fi
  \ifx\csstampdate\@empty\else\quad\csstampdate\fi
}
\ifx\csstampid\@empty\else
  \definecolor{csstampgrey}{gray}{0.45}
  % Background shipout material: it occupies no space, so the text block and the \thanks
  % footnote are byte-identical to an unstamped compile. The starred form applies to the
  % NEXT page shipped only — i.e. page 1. Placed in the left margin (the text block starts
  % at 1.1in), reading bottom-to-top, in a Times-like face at 8pt.
  \AddToShipoutPictureBG*{%
    \AtPageLowerLeft{%
      \put(\LenToUnit{0.42in},\LenToUnit{1.1in}){%
        \rotatebox{90}{%
          \fontfamily{ptm}\fontsize{8}{10}\selectfont
          \ifx\csstamplink\@empty
            \textcolor{csstampgrey}{\cs@stamptext}%
          \else
            \expandafter\href\expandafter{\csstamplink}{\textcolor{csstampgrey}{\cs@stamptext}}%
          \fi
        }%
      }%
    }%
  }
\fi
\makeatother


\title{Estimating optimal treatment value with many covariate cells and shrinking overlap}

\author{CausalSmith\thanks{Machine-generated by the CausalSmith research pipeline.}}

\date{\cspaperdate}

\begin{document}

\maketitle

\begin{abstract}
We characterize the minimax squared-error risk for estimating the observed optimal treatment value from independent and identically distributed observational data with categorical covariates, binary treatments, and binary outcomes. For sample size \(n\), number of covariate cells \(d\ge2\), and public treatment-overlap floor \(0<\epsilon\le1/2\), the risk is of order
\[
\min\left\{1,\frac{d}{n\epsilon\log(ed)}\right\},
\]
with universal comparison constants. The characterization is uniform over unknown cell masses, propensities satisfying the public floor, and conditional outcome means, including treatment-effect ties and null cells. We construct a deterministic observed-data estimator that attains this rate through armwise localization, Jackson polynomial approximation, factorial statistics, and conditional averaging of an auxiliary capped-Poisson construction. The matching lower bound uses weighted approximation and moment-matched probability priors with prescribed common mean, together with a comparison of the full observed likelihood. Consequently, uniform mean-square consistency along growing-cell and shrinking-overlap sequences is possible exactly when \(d/[n\epsilon\log(ed)]\to0\). Under consistency and armwise conditional exchangeability, the same risk characterization applies to the causal oracle value.
\end{abstract}

\section{Introduction}

The population value of individualized treatment choice averages the best conditional outcome attainable in each covariate cell. We establish its minimax squared-error risk in a finite observational experiment with independent and identically distributed observations, categorical covariates, binary treatments, and binary outcomes. With \(n\) the sample size, \(d\ge2\) the number of covariate cells, and \(0<\epsilon\le1/2\) the public treatment-overlap floor, \cref{obj:thm:matched-frontier} gives the rate
\[
\min\left\{1,\frac{d}{n\epsilon\log(ed)}\right\}
\]
with comparison constants universal in all three parameters. The result characterizes the joint cost of an expanding covariate alphabet and weakening overlap for estimation of a scalar optimized population value.

This target follows the optimal-regime and welfare-based treatment-choice formulations of \citet{Murphy2003,Manski2004}. In each cell, it selects the larger of the two conditional outcome means and weights that choice by the population cell probability, as specified in \cref{obj:def:observed-value}. Under consistency, armwise conditional exchangeability, and the public overlap restriction, \cref{obj:prop:identification} identifies this observed functional with the causal oracle value. It also establishes that every admissible observed law has a causal completion. These properties connect the observed-law minimax problem to estimation of the causal oracle value over the corresponding full-data class.

The uniformity domain is central to the result. The class in \cref{obj:def:observed-class} permits arbitrary unknown cell probabilities and conditional outcome means, together with heterogeneous unknown treatment propensities satisfying the public floor in every occupied cell. It includes null cells, boundary outcome means, exact treatment-effect ties, and arbitrarily small positive treatment contrasts. The risk comparison holds for every admissible sample size, alphabet size, and overlap floor. In particular, its constants remain universal as overlap shrinks.

The rate yields a necessary and sufficient condition for uniform mean-square consistency. For sequences \(d_n\), the cell counts, and \(\epsilon_n\), the public overlap floors, satisfying \(d_n\ge2\) and \(0<\epsilon_n\le1/2\), \cref{obj:thm:consistency-threshold} establishes that uniform consistency is possible exactly when
\[
\frac{d_n}{n\epsilon_n\log(e d_n)}\longrightarrow0.
\]
The quantity \(n\epsilon_n\) is the rare-arm sample-size scale allowed by the overlap restriction. The criterion therefore compares that scale with the number of cells divided by its logarithm. For a fixed alphabet, the risk has order \(\min\{1,1/(n\epsilon)\}\), with constants depending on the alphabet; for fixed overlap, it has order \(\min\{1,d/[n\log(ed)]\}\), with constants depending on the overlap floor.

The upper bound is attained by the deterministic armwise factorial estimator in \cref{obj:def:armwise-estimator}. Its active branch uses pilot counts to localize the four observed atom probabilities in each cell, approximates the cell contribution with a tensor Jackson polynomial, and estimates polynomial terms through factorial count statistics. Armwise clipping accounts for the sensitivity of the unknown treatment denominators. Conditional averaging over an auxiliary capped-Poisson count and its allocation marks produces a deterministic function of the observed sample. The parameterized definition also includes empirical estimation for bounded alphabets and a midpoint statistic in the saturated regime. \Cref{obj:thm:observable-upper} guarantees a universal bounded-alphabet cutoff \(D_0\ge2\); the attaining estimator uses the branches specified in \cref{obj:def:armwise-estimator}.

The converse combines weighted polynomial approximation with constrained moment priors. For \(K\), the polynomial degree, and \(M\), the scalar intensity support endpoint, \cref{obj:thm:weighted-separation-and-frontier} establishes that weighted approximation error and target separation between moment-matched probability priors of common mean one have order \(\sqrt M/K\), uniformly for \(K\ge2\), \(2\le M\le K^2\), and \(0\le\epsilon\le1/2\). A localized Jackson packet and common-atom completion construct the required priors. The statistical reduction retains normalization of the sparse alternatives and the information carried by both treatment arms, yielding the lower bound for every measurable estimator in \cref{obj:thm:all-estimator-lower}.

The closest estimand-specific literature develops inference for optimal treatment values under differentiability, convergence, variance, or smoothing conditions \citep[Theorems~1--2]{LuedtkeVanderLaan2016OptimalValue}; \citep[Theorems~2--3]{Shi2020}; \citep[Theorem~3.3 and Corollary~3.4]{WhitehouseChenAusternSyrgkanis2025}. \citet[Theorem~3]{CausalSmithFixedOverlap2026} prove the same-target fixed-overlap rate \(\min\{1,d/[n\log(ed)]\}\) for every fixed \(\epsilon\in(0,1/2)\), with constants allowed to depend on that floor. Our result is uniform for \(0<\epsilon\le1/2\) as the floor shrinks: it identifies the extra \(\epsilon^{-1}\) factor with universal comparison constants. The proof of this extension uses armwise sensitivity control, weighted common-mean priors, and comparison of the full boundary-propensity likelihood. The discrete causal analysis of \citet[Theorems~1--2]{ZengBalakrishnanHanKennedy2024Discrete} provides a nearby comparison for linear arm means and treatment effects. Here the cellwise maximum introduces a nonlinear target. Methodologically, the construction adapts polynomial approximation and unbiased polynomial estimation from discrete-functional estimation \citep{JiaoVenkatHanWeissman2015Functionals,JiaoHanWeissman2018L1} to maxima of unknown observational arm ratios.

\Cref{sec:appendix-lower-bound} develops the statistical lower bound containing \cref{obj:thm:all-estimator-lower}.

The paper proceeds through related work, the observational model and causal interpretation, the minimax rates and consistency criterion, the attaining estimator, and discussion of implications and open questions. The appendices develop identification and the upper bound, weighted approximation and constrained moment priors, the statistical lower bound and verification scope, and proofs of the main results.


\section{Related work}

The comparison with existing work turns on three distinctions: the population quantity being estimated, the information supplied by the statistical experiment, and the class over which performance is uniform. Our target is the scalar optimized value in \cref{obj:def:observed-value}, with its causal interpretation established in \cref{obj:prop:identification}. The \leanref{S-1}{finite observational experiment} has categorical covariates, binary treatments, and binary outcomes. Its sample size \leanref{sym:n}{\ensuremath{n}}, number of covariate cells \leanref{sym:d}{\ensuremath{d}}, and public overlap floor \leanref{sym:\epsilon}{\ensuremath{\epsilon}} jointly index the risk comparison. Under \cref{obj:ass:iid-sampling,obj:ass:shrinking-overlap}, the uniformity domain is the observed-law class in \cref{obj:def:observed-class}. Within that class, \cref{obj:thm:matched-frontier,obj:thm:consistency-threshold} characterize squared-error risk and the conditions for uniform consistency as the number of cells grows and overlap shrinks.

The closest estimand-specific literature studies differentiability and inference for the mean outcome under an optimal treatment rule. \citet[Theorem~1 and equation~(3)]{LuedtkeVanderLaan2016OptimalValue} characterize pathwise differentiability under strictly positive propensities, bounded outcomes, and finite variance of the candidate gradient. Their condition permits treatment-effect ties when both conditional outcome variances vanish on the tie set. Their online estimator also supports asymptotic inference at exceptional laws under conditions controlling standardized increments, variance estimation, the nuisance remainder, and the value loss of estimated rules \citep[Theorem~2, conditions~(C1)--(C5)]{LuedtkeVanderLaan2016OptimalValue}. The efficiency conclusion in their Corollary~3 adds fixed positivity, boundedness, convergence, and variance conditions. These results describe inference under specified asymptotic conditions at the underlying law; the present comparison describes worst-case squared-error risk over an entire finite-cell class indexed by sample size and overlap.

Subsample aggregation provides another approach to inference for the same optimized-value estimand. \citet[Theorems~2--3]{Shi2020} combine treatment-rule estimation on subsamples with evaluation on complementary observations. Their analysis assumes fixed positivity and an expected value-loss rate with exponent greater than one half, together with compatible subsample growth \citep[conditions~(A3)--(A4)]{Shi2020}. For a rule obtained from an estimated treatment contrast, their margin and contrast-estimation conditions provide sufficient conditions for this value-loss requirement \citep[conditions~(A5)--(A6)]{Shi2020}. The feasible inference result additionally imposes variance and nuisance-estimation conditions \citep[Theorem~3 and equations~(9)--(11)]{Shi2020}. This approach explains how aggregation can accommodate nonunique optimal rules while supporting asymptotic inference. Our finite-cell risk characterization instead takes the uniformity domain directly from the observed-law and overlap restrictions.

Softmax smoothing connects optimal-value inference to estimation of a differentiable surrogate. \citet[Theorem~3.3 and Corollary~3.4]{WhitehouseChenAusternSyrgkanis2025} establish asymptotic linearity and studentized inference under strong positivity, controlled density of positive suboptimality gaps near zero, suitable growth of the inverse smoothing scale, and nuisance-estimation conditions \citep[Assumptions~1--2]{WhitehouseChenAusternSyrgkanis2025}. Their density condition permits an atom at zero, thereby accommodating ties. The theorem balances smoothing bias against a smoothing-dependent mean-square rate for the outcome regression and a product-rate condition for the regression and its representer; boundedness and positive limiting variance also enter the inference guarantee. The corresponding distinction here is the uniformity domain: \cref{obj:thm:matched-frontier} evaluates estimation risk across the finite-cell class, including laws with ties and arbitrarily small positive treatment contrasts, as the public overlap floor varies.

The nearest comparison in the discrete causal experiment is \citet{ZengBalakrishnanHanKennedy2024Discrete}. Their June 2026 revision studies arm means and average treatment effects with many discrete covariate categories. For a treated-arm mean, their Theorem~1 gives a uniform mean-squared-error upper bound of the form
\[
\frac{(1-\epsilon)^2}{\epsilon^2}\frac{d^2}{n^2}
+\frac{C}{n\epsilon},
\]
where \(C\) is an absolute constant \citep[Theorem~1]{ZengBalakrishnanHanKennedy2024Discrete}. Their minimax lower bound has order
\[
\left(\frac{d}{n\log n}\right)^2+\frac{1}{n}
\]
in the regime \(d\lesssim n\log n\), with a comparison constant depending on overlap \citep[Theorem~2]{ZengBalakrishnanHanKennedy2024Discrete}. These results expose the interaction between sparse cell counts and propensity denominators for a target that is linear in the arm-specific outcome means. The optimized value adds a maximum within each cell. Our comparison retains the overlap dependence in both sides of the risk characterization for this nonlinear target.

The approximation architecture draws on the literature on estimating functionals of discrete distributions. \citet{JiaoVenkatHanWeissman2015Functionals} combine polynomial approximation in nonsmooth regions with unbiased estimation of polynomial terms and bias correction in smooth regions. Their entropy characterization applies when the sample size is at least a constant multiple of alphabet size divided by its logarithm, and separates approximation bias from a variance contribution \citep[Theorem~1]{JiaoVenkatHanWeissman2015Functionals}. Thus the logarithmic improvement has an estimand-specific risk expression and a stated sampling regime. In our experiment, the localized approximation is applied to the armwise extension in \cref{obj:def:armwise-extension}, and the attaining construction estimates its polynomial contributions through factorial statistics \cref{obj:def:armwise-estimator}. The propensity denominators make the overlap floor part of the approximation and statistical analysis.

The \(L_1\)-distance results of \citet{JiaoHanWeissman2018L1} provide a particularly close comparison because absolute-value nonsmoothness is closely related to the maximum of two arm means. They treat both estimation relative to a known reference distribution and estimation from samples from two unknown distributions. For alphabet size \(d\), their worst-case squared-error rate \(d/(n\log n)\) applies in the regime \(n\gtrsim d/\log d\) and \(\log n\lesssim\log d\), with the precise distribution-specific qualifications stated for the known-reference problem \citep[Theorems~2--4 and~6, Remark~1]{JiaoHanWeissman2018L1}. These regime restrictions matter when comparing logarithmic improvements across estimands. Here, weighted approximation and priors with prescribed common mean and matching moments connect the scalar approximation problem to the causal experiment \cref{obj:def:weighted-error,obj:def:constrained-prior-separation,obj:thm:weighted-separation-and-frontier}.

Policy learning evaluates a different statistical decision. \citet[Assumption~2.1 and Theorems~2.1--2.2]{Kitagawa2018} establish uniform welfare-regret bounds for empirical welfare maximization over a policy class, with class complexity and overlap entering the comparison. \citet[Theorem~1]{Athey2021} use estimated doubly robust scores to obtain asymptotic regret guarantees under nuisance-estimation, policy-class, and outcome conditions. \citet{JinRenYangWang2025Pessimism} construct pessimistic policy choices using known behavior propensities, with guarantees governed by coverage of an optimal policy and the complexity of the candidate class, including settings with adaptive data collection. In these problems, the statistical output is a treatment rule, and loss measures its welfare shortfall relative to a policy-class optimum. The risk in \cref{obj:def:minimax-risk} instead measures squared error in a scalar population value optimized over all cellwise treatment choices.

Specified-policy evaluation offers a further useful distinction. \citet{Ma2022} characterize minimax performance for evaluating a given target policy in a multi-armed bandit experiment, separating known, unknown, and partially known behavior policies. Their partial-knowledge analysis also uses polynomial approximation when a public lower bound on behavior probabilities is available \citep[Section~3.3]{Ma2022}. This supplies a related treatment of limited sampling probabilities, while the present target incorporates optimization over the conditional outcome means.

The closest same-target predecessor is the fixed-overlap result of \citet[Theorem~3]{CausalSmithFixedOverlap2026}. For every fixed \(\epsilon\in(0,1/2)\), it proves the rate \(\min\{1,d/[n\log(ed)]\}\), with comparison constants that may depend on \(\epsilon\), using the same Jackson--factorial architecture. The present result identifies the joint rate for \(0<\epsilon\le1/2\) when that floor may shrink: \(\min\{1,d/[n\epsilon\log(ed)]\}\) with universal comparison constants, and the corresponding consistency threshold in \cref{obj:thm:consistency-threshold}. The additional argument controls the maximum through separate armwise sensitivities, constructs weighted moment-matched priors with prescribed common mean, and compares the full observed likelihood at boundary propensities. Thus the fixed-overlap specialization is recovered while the dependence on overlap is made uniform and explicit.


\section{Observational model and causal interpretation}

Throughout, \(c\) and \(C\) denote positive universal constants whose values may change between occurrences. The notation \(\|v\|_1\) denotes the sum of the absolute coordinates of a finite vector. We use \(\lesssim\), \(\gtrsim\), and \(\asymp\) for comparisons up to positive universal constants, and all logarithms are natural. Sample indices run from \(1\) to \(n\), treatment and outcome indices take values in \(\{0,1\}\), and covariate indices use the zero-based alphabet defined below. The statistical analysis concerns \(n\ge1\) observed units, \(d\ge2\) covariate cells, and a public overlap floor \(0<\epsilon\le1/2\); the cell count and overlap floor may vary with sample size.

\subsection{Observed experiment and loss}

We first state the sampling condition shared by every observed law in the experiment.

\begin{assumptionv}{ass:iid-sampling}[Independent and identically distributed observations]
For each \(\mathbb P\in\mathcal D_{d,\epsilon}^{\mathrm{obs}}\), the observed-law class, the observed units \(O_1,\ldots,O_n\) are independent and have common law \(\mathbb P\).
\end{assumptionv}

Independent and identically distributed sampling controls dependence across units within each experiment; allowing \(d\) and \(\epsilon\) to vary with \(n\) specifies a sequence of such experiments \citep{ZengBalakrishnanHanKennedy2024Discrete}.

The next definition fixes the covariate alphabet \leanref{sym:[d]}{\ensuremath{[d]}} and observed atom index set \leanref{sym:\mathcal J_d}{\ensuremath{\mathcal J_d}}. We write \leanref{sym:x}{\ensuremath{x}} for a cell index, \leanref{sym:a}{\ensuremath{a}} for a treatment index, and \leanref{sym:y}{\ensuremath{y}} for an outcome index.

\begin{definitionv}{synth_1}[Covariate alphabet and atoms]
The finite covariate alphabet \([d]\) and observed atom index set \(\mathcal J_d\) are defined, for each nonnegative integer \(d\), by
\[
[d] := \{x\in\mathbb Z : 0\le x<d\},
\qquad
\mathcal J_d := [d]\times\{0,1\}\times\{0,1\}.
\]
\end{definitionv}

Cell masses determine how much each covariate category contributes to the population target. Propensities allocate that mass between treatment arms, while outcome means summarize the observed response within each arm. The model restricts treatment allocation through the following public floor.

The restriction in \cref{obj:ass:shrinking-overlap} is the triangular-array version of standard strong positivity: both arms have positive conditional probability in every occupied cell, with a floor that may shrink along the sequence of experiments \citep{ZengBalakrishnanHanKennedy2024Discrete}. A public floor supplies a common bound for the statistical analysis and estimator tuning, while permitting heterogeneous propensities across cells. Null cells follow the zero-ratio convention in \cref{obj:def:observed-class}. These quantities specify the observed-law class \leanref{sym:\mathcal D_{d,\epsilon}^{\mathrm{obs}}}{\ensuremath{\mathcal D_{d,\epsilon}^{\mathrm{obs}}}} of \cref{obj:def:observed-class}.

\begin{assumptionv}{ass:shrinking-overlap}[Public treatment overlap restriction]
Let \(p_x\) denote the probability of covariate cell \(x\), and let \(\pi_x\) denote its conditional treatment-one probability. For every \(x\in[d]\), the finite covariate alphabet, with \(p_x>0\), the overlap parameter \(\epsilon\) satisfies
\[
\epsilon\le\pi_x\le1-\epsilon.
\]
\end{assumptionv}

The following definition collects the atom probabilities, cell masses, propensities, outcome means, and their probability constraints into the observed-law class.

\begin{definitionv}{def:observed-class}[Observed laws under overlap]
The observed class \(\mathcal D_{d,\epsilon}^{\mathrm{obs}}\) consists of probability laws \(\mathbb P\) on \([d]\times\{0,1\}^2\) satisfying \cref{obj:ass:shrinking-overlap} on every cell with \(p_x>0\). Its parameters are \(d\in\mathbb N\) and \(\epsilon\in\mathbb R\), and the covariate alphabet is
\[
[d]:=\{x\in\mathbb Z:0\le x<d\}.
\]
For the coordinate variables \(X\), \(A\), and \(Y\), representing the covariate, treatment, and observed outcome, respectively, the observed atom probabilities, cell probabilities, treatment propensities, and arm-specific outcome means are defined by
\[
q_{ay,x}:=\mathbb P(X=x,A=a,Y=y),
\qquad
p_x:=\sum_{a\in\{0,1\}}\sum_{y\in\{0,1\}}q_{ay,x},
\]
\[
\pi_x:=\frac{q_{10,x}+q_{11,x}}{p_x},
\qquad
\mu_{ax}:=\frac{q_{a1,x}}{q_{a0,x}+q_{a1,x}},
\]
with each ratio defined as zero when its denominator is zero. These quantities satisfy, for every \(x\in[d]\) and \(a\in\{0,1\}\),
\[
q_{a1,x}=p_x\pi_x^a(1-\pi_x)^{1-a}\mu_{ax},
\qquad
q_{a0,x}=p_x\pi_x^a(1-\pi_x)^{1-a}(1-\mu_{ax}),
\]
\[
p_x\ge0,
\qquad
\sum_{x\in[d]}p_x=1,
\qquad
\mu_{ax}\in[0,1].
\]
\end{definitionv}

The class allows arbitrary cell masses and outcome means subject to the displayed probability constraints, and arbitrary propensities satisfying the public floor.

The scalar target is the observed optimal treatment value \leanref{sym:\Psi(\mathbb P)}{\ensuremath{\Psi(\mathbb P)}} of \cref{obj:def:observed-value}. It chooses the larger conditional outcome mean in each cell and averages those choices using the population cell masses.

\begin{definitionv}{def:observed-value}[Observed optimal treatment value]
The observed value of a probability law \(\mathbb P\) on \([d]\times\{0,1\}\times\{0,1\}\), for a nonnegative integer \(d\) and the zero-based covariate alphabet \([d]\) of \cref{obj:synth_1}, is
\[
\Psi(\mathbb P)
=\sum_{x\in[d]}p_x\max_{a\in\{0,1\}}\mu_{ax}.
\]
Here the observed atom probabilities, covariate cell probabilities, and arm-specific outcome means are defined by
\[
q_{ay,x}=\mathbb P(X=x,A=a,Y=y),
\qquad
p_x=\sum_{a\in\{0,1\}}\sum_{y\in\{0,1\}}q_{ay,x},
\qquad
\mu_{ax}=\frac{q_{a1,x}}{q_{a0,x}+q_{a1,x}},
\]
with a zero denominator giving \(\mu_{ax}=0\). Each null-cell summand is therefore zero.
\end{definitionv}

The maximum introduces the cellwise optimization into the functional, while the zero-denominator convention gives every null cell a zero contribution. We assess estimation of this scalar under squared error, uniformly over the observed-law class, through the minimax risk \leanref{sym:\mathfrak R_{n,d,\epsilon}}{\ensuremath{\mathfrak R_{n,d,\epsilon}}} defined next.

\begin{definitionv}{def:minimax-risk}[Minimax risk \(\mathfrak R_{n,d,\epsilon}\)]
The minimax squared-error risk is
\[
\mathfrak R_{n,d,\epsilon}
:=
\inf_{\widehat V}
\sup_{\mathbb P\in\mathcal D_{d,\epsilon}^{\mathrm{obs}}}
E_{\mathbb P}
\left[
\left\{
\widehat V(O_1,\ldots,O_n)-\Psi(\mathbb P)
\right\}^{2}
\right],
\]
where \(O_1,\ldots,O_n\) denotes the observed sample, \(\Psi(\mathbb P)\) denotes the observed value, and the infimum ranges over all measurable real-valued estimators.
\end{definitionv}

Under \cref{obj:ass:iid-sampling}, the expectation is taken over the product sampling law. The supremum varies the cell masses, propensities, and outcome means within \cref{obj:def:observed-class}, so the risk measures uniform performance when those quantities are unknown. The subsequent rate comparisons use the logarithmic alphabet scale \leanref{sym:L_d}{\ensuremath{L_d}}, defined by
\[
L_d=\log(ed).
\]
This convention fixes the logarithmic factor in the risk bounds and consistency criterion.

\subsection{Causal interpretation and cellwise representation}

The observed functional acquires a causal interpretation through a \leanref{S-2}{binary potential-outcome overlay}. We introduce the full-data notation together so that the causal restrictions and completion class refer to the same experiment.



The oracle value chooses a treatment separately for each covariate cell using its conditional potential-outcome means. Relating those means to observed outcomes requires consistency and armwise conditional exchangeability. The first restriction connects the realized outcome to the potential outcome under the treatment received.

\begin{assumptionv}{ass:consistency}[Consistency of observed outcomes]
The observed outcome \(Y\) equals the potential outcome \(Y(A)\) under the realized treatment \(A\) almost surely:
\[
Y=Y(A)\quad\text{almost surely}.
\]
\end{assumptionv}

Consistency is the standard condition that gives the observed outcome its treatment-specific potential-outcome interpretation \citep{ZengBalakrishnanHanKennedy2024Discrete}. The second restriction addresses treatment selection within covariate cells.

\begin{assumptionv}{ass:conditional-exchangeability}[Conditional exchangeability of treatment]
For each treatment \(a\in\{0,1\}\), the potential outcome \(Y(a)\) is conditionally independent of the treatment \(A\) given the categorical covariate \(X\):
\[
Y(a)\perp A\mid X.
\]
\end{assumptionv}

Armwise conditional exchangeability is the standard condition under which conditioning on the covariate accounts for treatment selection relevant to each potential outcome \citep{ZengBalakrishnanHanKennedy2024Discrete}. Its substantive interpretation is that the recorded covariate provides sufficient adjustment for that selection. Together with consistency and positive overlap in occupied cells, it supplies the conditions for the value identification in \cref{obj:prop:identification}. We collect these restrictions in the causal completion class \leanref{sym:\mathcal P_{d,\epsilon}}{\ensuremath{\mathcal P_{d,\epsilon}}}.

\begin{definitionv}{def:causal-class}[Causal-law class \(\mathcal P_{d,\epsilon}\)]
Define
\[
\mathcal P_{d,\epsilon}
:=
\left\{
P:
P\text{ is a law of }(X,A,Y,Y(0),Y(1))
\text{ satisfying the conditions below}
\right\}.
\]
Here \(\operatorname{Obs}(P)\) denotes the observed marginal law of \((X,A,Y)\).
\begin{itemize}
\item \textbf{(Consistency.)} \(P\) satisfies \cref{obj:ass:consistency}.
\item \textbf{(Exchangeability.)} \(P\) satisfies \cref{obj:ass:conditional-exchangeability}.
\item \textbf{(Observed marginal.)}
\[
\operatorname{Obs}(P)\in\mathcal D_{d,\epsilon}^{\mathrm{obs}},
\]
whose overlap restriction is given in \cref{obj:ass:shrinking-overlap}.
\end{itemize}
This definition applies for every \(n\ge1\), \(d\ge2\), and \(0<\epsilon\le1/2\).
\end{definitionv}

The statistical problem is defined over observed laws, while this class specifies their causal interpretation. To connect value identification to the estimator construction, we also introduce a representation of each cell's contribution in terms of its four observed atom probabilities.



The armwise extension \leanref{sym:F_\epsilon(u)}{\ensuremath{F_\epsilon(u)}} of \cref{obj:def:armwise-extension} represents the contribution of an admissible cell and supplies a function on arbitrary nonnegative four-vectors for the estimator. For a generic nonnegative four-vector \leanref{sym:u}{\ensuremath{u}}, its arm mass \leanref{sym:s_a(u)}{\ensuremath{s_a(u)}}, total mass \leanref{sym:s(u)}{\ensuremath{s(u)}}, and overlap-truncated arm denominator \leanref{sym:D_a(u)}{\ensuremath{D_a(u)}} are defined together with the extension.

\begin{definitionv}{def:armwise-extension}[Armwise extension \(F_\epsilon\)]
The armwise extension is
\[
F_\epsilon(u)
=
\begin{cases}
\displaystyle
\max_{a\in\{0,1\}}\frac{s(u)u_{a1}}{D_a(u)},&u\ne0,\\[6pt]
0,&u=0,
\end{cases}
\qquad u\in\mathbb R_+^4,
\]
where the coordinates of \(u\) are indexed by \(a,y\in\{0,1\}\), and
\[
s_a(u)=u_{a0}+u_{a1},
\qquad
s(u)=s_0(u)+s_1(u),
\qquad
D_a(u)=\max\{s_a(u),\epsilon s(u)\}.
\]
\end{definitionv}

The denominator uses the larger of the arm mass and the public floor times total mass, and the zero branch retains the null-cell convention. The following result connects this representation to the observed target and connects that target to the causal oracle value.

\begin{propositionv}{prop:identification}[Completion and value identification]
Let \(d\) be the number of covariate cells and \(\epsilon\) the overlap floor, with
\begin{itemize}
\item \textbf{(Cell count.)} \(d\in\mathbb N\) and \(d\ge 2\).
\item \textbf{(Positive overlap.)} \(0<\epsilon\).
\item \textbf{(Overlap range.)} \(\epsilon\le 1/2\).
\end{itemize}
The following statements hold.

Every observed law \(\mathbb P\in\mathcal D_{d,\epsilon}^{\mathrm{obs}}\) of \cref{obj:def:observed-class} admits a full-data completion \(P\in\mathcal P_{d,\epsilon}\) of \cref{obj:def:causal-class} such that
\[
\operatorname{Obs}(P)=\mathbb P,
\]
where \(\operatorname{Obs}(P)\) is the marginal law of the categorical covariate \(X\), binary treatment \(A\), and binary observed outcome \(Y\).

Every \(P\in\mathcal P_{d,\epsilon}\) satisfies
\[
V^\star(P)
=
\sum_{x\in[d]} P(X=x)
\max_{a\in\{0,1\}}
\mathbb E_P[Y(a)\mid X=x]
=
\Psi\{\operatorname{Obs}(P)\},
\]
where \([d]\) is the finite covariate alphabet, \(Y(a)\) is the potential outcome under arm \(a\), \(V^\star(P)\) is the causal oracle value defined by the displayed sum, and \(\Psi\) is the observed-value functional of \cref{obj:def:observed-value}. Conditional potential-outcome means in cells of zero probability are assigned zero.

For every \(\mathbb P\in\mathcal D_{d,\epsilon}^{\mathrm{obs}}\) and every \(x\in[d]\), define the four-vector of observed atom probabilities and its cell mass and armwise outcome means by
\[
q_{ay,x}=\mathbb P(X=x,A=a,Y=y),
\qquad
q_x=(q_{00,x},q_{01,x},q_{10,x},q_{11,x}),
\]
\[
p_x=\sum_{a=0}^{1}\sum_{y=0}^{1}q_{ay,x},
\qquad
\mu_{ax}=\frac{q_{a1,x}}{q_{a0,x}+q_{a1,x}},
\]
with a zero denominator giving \(\mu_{ax}=0\). The armwise extension of \cref{obj:def:armwise-extension} then satisfies
\[
F_\epsilon(q_x)=p_x\max_{a\in\{0,1\}}\mu_{ax}.
\]
In particular, both sides equal zero when \(p_x=0\).
\end{propositionv}

\Cref{obj:prop:identification} gives the observed estimation problem its causal scope. Every law in the observed class has an admissible causal completion, and every such completion has the same oracle value determined by its observed marginal. Consequently, the minimax risk in \cref{obj:def:minimax-risk} also describes estimation of the causal oracle value over \cref{obj:def:causal-class} using observed samples. The cellwise identity supplies the representation
\[
\Psi(\mathbb P)=\sum_{x\in[d]}F_\epsilon(q_x).
\]
This representation preserves the target on the overlap-constrained observed class while extending each cell contribution to the nonnegative vectors used in the estimator construction. Identification therefore connects the causal target, the observed loss, and the cellwise estimation problem under the stated sampling and causal restrictions.


\section{Minimax rates and consistency}

The minimax squared-error risk \leanref{sym:\mathfrak R_{n,d,\epsilon}}{\ensuremath{\mathfrak R_{n,d,\epsilon}}}, defined in \cref{obj:def:minimax-risk}, measures the precision attainable uniformly over the observed class. Its rate depends jointly on the number of covariate cells and the sample-size scale associated with the public overlap floor. The following result characterizes this dependence with universal comparison constants and identifies a deterministic estimator attaining the rate.

\begin{theoremv}{thm:matched-frontier}[Universal matched risk frontier]
There exist universal real constants \(c_\star,C_\star,H_0,\kappa\) and an integer \(D_0\) satisfying
\[
0<c_\star\le C_\star,\qquad H_0>0,\qquad \kappa>0,\qquad D_0\ge2,
\]
such that the following bounds hold for every:
\begin{itemize}
\item \textbf{(Sample size.)} Integer \(n\ge1\).
\item \textbf{(Alphabet size.)} Integer \(d\ge2\).
\item \textbf{(Overlap.)} Real overlap parameter \(0<\epsilon\le1/2\).
\end{itemize}
Write \(L_d=\log(ed)\) for the logarithmic alphabet scale, and let
\(\widehat V_{n,d,\epsilon}^{\mathrm{AF}}\) be the prescribed deterministic armwise estimator with calibration constants \(H_0,\kappa,D_0\). The observed minimax squared-error risk,
\[
\mathfrak R_{n,d,\epsilon}
=
\inf_{\widehat V}
\sup_{\mathbb P\in\mathcal D_{d,\epsilon}^{\mathrm{obs}}}
E_{\mathbb P}\!\left[
\bigl(\widehat V-\Psi(\mathbb P)\bigr)^2
\right],
\]
where the infimum ranges over measurable real-valued estimators based on \(n\) observed units, satisfies
\[
c_\star\min\left\{1,\frac{d}{n\epsilon L_d}\right\}
\le
\mathfrak R_{n,d,\epsilon}
\le
\sup_{\mathbb P\in\mathcal D_{d,\epsilon}^{\mathrm{obs}}}
E_{\mathbb P}\!\left[
\bigl(\widehat V_{n,d,\epsilon}^{\mathrm{AF}}-\Psi(\mathbb P)\bigr)^2
\right]
\le
C_\star\min\left\{1,\frac{d}{n\epsilon L_d}\right\}.
\]
For the causal class, the same estimator and constants satisfy
\[
\begin{aligned}
c_\star\min\left\{1,\frac{d}{n\epsilon L_d}\right\}
&\le
\inf_{\widehat V}
\sup_{P\in\mathcal P_{d,\epsilon}}
E_{\operatorname{Obs}(P)}\!\left[
\bigl(\widehat V-V^\star(P)\bigr)^2
\right]
\\
&\le
\sup_{P\in\mathcal P_{d,\epsilon}}
E_{\operatorname{Obs}(P)}\!\left[
\bigl(\widehat V_{n,d,\epsilon}^{\mathrm{AF}}-V^\star(P)\bigr)^2
\right]
\\
&\le
C_\star\min\left\{1,\frac{d}{n\epsilon L_d}\right\},
\end{aligned}
\]
where the infimum again ranges over measurable real-valued estimators based on \(n\) observed units, \(V^\star(P)\) is the causal oracle value, and the expectation is under the observed marginal law \(\operatorname{Obs}(P)\).
\end{theoremv}

The ratio \(d/(n\epsilon L_d)\) summarizes the joint statistical cost of many cells and weak overlap. Here \(n\epsilon\) is the rare-arm sample-size scale permitted by the overlap floor: when an arm's treatment probability equals that floor throughout the population, its expected sample count is \(n\epsilon\). Holding \(n\) and \(\epsilon\) fixed, the characterized comparison scale increases with \(d\) through \(d/\log(ed)\). In the saturation regime \(n\epsilon\le d/L_d\), the minimax risk is bounded above and below by positive universal constants. When \(n\epsilon>d/L_d\), the squared-error rate is \(d/(n\epsilon L_d)\). The causal interpretation of these bounds follows from \cref{obj:prop:identification}, which connects the observed target and class to the causal oracle value and its admissible completions.

Along sequences in which the cell count and overlap floor vary with sample size, the same bounds give an exact criterion for uniform mean-square consistency.

\begin{theoremv}{thm:consistency-threshold}[Consistency threshold and fixed-parameter rates]
Let \(\mathfrak R_{n,d,\epsilon}\) be the minimax squared-error risk in \cref{obj:def:minimax-risk}. Consider sequences \((d_n)_{n\in\mathbb N}\) of alphabet sizes and \((\epsilon_n)_{n\in\mathbb N}\) of overlap floors satisfying:
\begin{itemize}
\item \textbf{(Integer alphabet.)} \(d_n\in\mathbb N\) for every \(n\).
\item \textbf{(Alphabet size.)} \(d_n\ge2\) for every \(n\).
\item \textbf{(Positive overlap.)} \(\epsilon_n>0\) for every \(n\).
\item \textbf{(Overlap range.)} \(\epsilon_n\le1/2\) for every \(n\).
\end{itemize}
There exists a sequence of measurable real-valued estimators \(\widehat V_n\), each based on \(n\) independent observed units, such that
\[
\sup_{\mathbb P\in\mathcal D_{d_n,\epsilon_n}^{\mathrm{obs}}}
\mathbb E_{\mathbb P}
\bigl[(\widehat V_n-\Psi(\mathbb P))^2\bigr]
\longrightarrow0
\]
if and only if
\[
\frac{d_n}{n\epsilon_n\log(e d_n)}
\longrightarrow0
\qquad (n\to\infty).
\]

For each fixed integer \(d\ge2\), there exist real constants \(c,C\), depending on \(d\), with \(0<c\le C\), such that for every integer \(n\ge1\) and every overlap floor \(0<\epsilon\le1/2\),
\[
c\min\left\{1,\frac{1}{n\epsilon}\right\}
\le \mathfrak R_{n,d,\epsilon}
\le C\min\left\{1,\frac{1}{n\epsilon}\right\}.
\]

For each fixed overlap floor \(0<\epsilon\le1/2\), there exist real constants \(c,C\), depending on \(\epsilon\), with \(0<c\le C\), such that for every pair of integers \(n\ge1\) and \(d\ge2\),
\[
c\min\left\{1,\frac{d}{n\log(e d)}\right\}
\le \mathfrak R_{n,d,\epsilon}
\le C\min\left\{1,\frac{d}{n\log(e d)}\right\}.
\]
\end{theoremv}

The criterion requires the rare-arm sample-size scale to grow faster than the cell count divided by its logarithmic scale. It therefore allows growing alphabets and shrinking overlap simultaneously, provided their joint burden vanishes relative to the available sample. For a fixed alphabet, the cell-count factor is absorbed into constants depending on that alphabet, yielding the rate \(\min\{1,1/(n\epsilon)\}\); consistency then requires \(n\epsilon_n\to\infty\). For fixed overlap, constants may depend on the overlap floor, and the rate becomes \(\min\{1,d/[n\log(ed)]\}\). These specializations make the joint criterion directly interpretable for designs with either a stable covariate alphabet or a stable overlap floor.

The attaining construction is specified in \cref{obj:def:armwise-estimator}, and its uniform risk bound is established in \cref{obj:thm:observable-upper}. The converse in \cref{obj:thm:all-estimator-lower} applies to every measurable estimator. Its central ingredient is the weighted approximation and constrained-prior separation result in \cref{obj:thm:weighted-separation-and-frontier}, which supplies the target separation used in the statistical lower bound. \Cref{sec:appendix-weighted,sec:appendix-lower-bound} develop the approximation results and the associated experiment construction.


\section{An attaining estimator}

The armwise factorial estimator \leanref{sym:\widehat V_{n,d,\epsilon}^{\mathrm{AF}}}{\ensuremath{\widehat V_{n,d,\epsilon}^{\mathrm{AF}}}} attains the rate in \cref{obj:thm:matched-frontier} using the observed sample and the public overlap floor. Its construction separates three regimes. The bounded-alphabet branch uses empirical cell frequencies and arm-specific outcome means. For larger alphabets, the saturated branch returns the midpoint of the value range when \(n\epsilon<d/L_d\), where \(L_d\) is the logarithmic alphabet scale. The active branch, used when \(n\epsilon\ge d/L_d\), combines pilot localization with polynomial approximation and factorial estimation of the cell contributions. The following algorithm specifies all three branches and the averaging that produces a deterministic statistic.

\begin{algorithmv}{def:armwise-estimator}[Armwise factorial estimator]
The estimator \(\widehat V_{n,d,\epsilon}^{\mathrm{AF}}\) is a deterministic function of the observed sample \(O_i=(X_i,A_i,Y_i)\in[d]\times\{0,1\}\times\{0,1\}\), \(i=1,\ldots,n\), using the zero-based alphabet \([d]=\{0,\ldots,d-1\}\) of \cref{obj:synth_1}. Its parameters satisfy:
\begin{itemize}
\item \textbf{(Tuning.)} The pilot-radius constant \(H_0\) and degree constant \(\kappa\) satisfy \(H_0>0\) and \(\kappa>0\); the integer alphabet cutoff \(D_0\) satisfies \(D_0\ge2\).
\item \textbf{(Sample and alphabet.)} The integers \(n\) and \(d\) satisfy \(n\ge1\) and \(d\ge2\).
\item \textbf{(Overlap parameter.)} The parameter \(\epsilon\) satisfies \(0<\epsilon\le1/2\).
\end{itemize}
Define the logarithmic scale, evaluation intensity, and approximation degree by
\[
L=L_d=\log(ed),\qquad
m=\frac n8,\qquad
K_n=\max\{2,\lfloor\kappa L\rfloor\}.
\]
Write \(\Pi_{[l,u]}(z)=\min\{u,\max\{l,z\}\}\) for projection onto an interval \([l,u]\).

For \(d<D_0\), define the full-sample counts
\[
N_x=\sum_{i=1}^{n}\mathbf 1\{X_i=x\},\qquad
N_{ax}=\sum_{i=1}^{n}\mathbf 1\{X_i=x,A_i=a\},\qquad
S_{ax}=\sum_{i=1}^{n}Y_i\mathbf 1\{X_i=x,A_i=a\},
\qquad x\in[d],\quad a\in\{0,1\},
\]
and set
\[
\widehat V_{n,d,\epsilon}^{\mathrm{AF}}
=
\Pi_{[0,1]}
\left(
\sum_{x\in[d]}\frac{N_x}{n}
\max_{a\in\{0,1\}}
\left\{\frac{S_{ax}}{N_{ax}}\right\}
\right),
\]
where each ratio \(S_{ax}/N_{ax}\) is defined to be zero when \(N_{ax}=0\).

For \(d\ge D_0\) and \(n\epsilon<d/L\), set
\[
\widehat V_{n,d,\epsilon}^{\mathrm{AF}}=\frac12.
\]

For \(d\ge D_0\) and \(n\epsilon\ge d/L\), define the estimator by the following auxiliary construction. Independently of the observations, draw
\[
\mathsf M\sim\operatorname{Pois}(n/4).
\]
On \(\{\mathsf M\le n\}\), draw independent fair marks, independently of the observations and of \(\mathsf M\),
\[
\mathsf B_i\sim\operatorname{Bernoulli}(1/2),
\qquad i=1,\ldots,\mathsf M.
\]
Define the pilot and evaluation atom counts by
\[
N'_{ay,x}
=
\sum_{i=1}^{\mathsf M}
\mathbf 1\{X_i=x,A_i=a,Y_i=y,\mathsf B_i=1\},
\]
\[
N_{ay,x}
=
\sum_{i=1}^{\mathsf M}
\mathbf 1\{X_i=x,A_i=a,Y_i=y,\mathsf B_i=0\},
\qquad a,y\in\{0,1\},\quad x\in[d].
\]
All subsequent cellwise quantities are defined on \(\{\mathsf M\le n\}\). Set
\[
c_{ay,x}=\frac{N'_{ay,x}}m,\qquad
h_{ay,x}
=
H_0\left\{\sqrt{\frac{c_{ay,x}L}{m}}+\frac Lm\right\},
\]
\[
\ell_{ay,x}=\max\{c_{ay,x}-h_{ay,x},0\},\qquad
u_{ay,x}=c_{ay,x}+h_{ay,x},
\]
\[
b_{ay,x}=\frac{\ell_{ay,x}+u_{ay,x}}2,\qquad
r_{ay,x}=\frac{u_{ay,x}-\ell_{ay,x}}2,
\]
\[
b_x=(b_{ay,x})_{a,y\in\{0,1\}},\qquad
r_x=(r_{ay,x})_{a,y\in\{0,1\}},\qquad
Q_x=\prod_{a,y\in\{0,1\}}[\ell_{ay,x},u_{ay,x}].
\]
Use \(F_\epsilon\), \(s\), and \(D_a\) from \cref{obj:def:armwise-extension}. Define the normalized Jackson kernel by
\[
J_{K_n}(t)
=
c_{K_n}^{\mathrm J}
\left\{\frac{\sin(K_nt/2)}{\sin(t/2)}\right\}^{4},
\qquad
\int_{-\pi}^{\pi}J_{K_n}(t)\,dt=1,
\]
where removable values are filled continuously and \(c_{K_n}^{\mathrm J}>0\) is the unique normalizing constant. With \(\odot\) denoting componentwise multiplication and cosine applied componentwise, define the algebraic polynomial \(P_{x,K_n}\) on \(Q_x\) by
\[
P_{x,K_n}(b_x+r_x\odot\cos\theta)
=
\int_{[-\pi,\pi]^4}
F_\epsilon\!\left(b_x+r_x\odot\cos(\theta+\tau)\right)
\prod_{a,y\in\{0,1\}}J_{K_n}(\tau_{ay})\,d\tau,
\qquad \theta\in\mathbb R^4.
\]
Its degree in each coordinate is at most
\[
D=2(K_n-1).
\]
Define the coefficients \(\beta_{x,\boldsymbol\alpha}\) by the unique expansion
\[
P_{x,K_n}(v)-F_\epsilon(b_x)
=
\sum_{\boldsymbol\alpha\in\{0,\ldots,D\}^4}
\beta_{x,\boldsymbol\alpha}
\prod_{a,y\in\{0,1\}}
\left(\frac{v_{ay}-b_{ay,x}}{r_{ay,x}}\right)^{\alpha_{ay}},
\qquad v\in Q_x.
\]
For nonnegative integers \(N\) and \(t\), define the falling factorial by
\[
(N)_t=\prod_{j=0}^{t-1}(N-j),\qquad (N)_0=1.
\]
For an integer \(k\ge0\) and a real center \(\zeta\), define
\[
U_k(N;\zeta)
=
\sum_{t=0}^{k}\binom kt(-\zeta)^{k-t}\frac{(N)_t}{m^t},
\]
and set
\[
Z_x
=
F_\epsilon(b_x)
+
\sum_{\boldsymbol\alpha\in\{0,\ldots,D\}^4}
\beta_{x,\boldsymbol\alpha}
\prod_{a,y\in\{0,1\}}
\frac{U_{\alpha_{ay}}(N_{ay,x};b_{ay,x})}
{r_{ay,x}^{\alpha_{ay}}}.
\]
Define the armwise clipping scales and clipped cell contributions by
\[
R_{a,x}=r_{a0,x}+r_{a1,x},\qquad
w_{a,x}=1+\frac{s(b_x)}{D_a(b_x)},\qquad
S_x^{\mathrm{aw}}=2\sum_{a=0}^{1}w_{a,x}R_{a,x},
\]
\[
T_x
=
\Pi_{[F_\epsilon(b_x)-d^{1/4}S_x^{\mathrm{aw}},\,
       F_\epsilon(b_x)+d^{1/4}S_x^{\mathrm{aw}}]}(Z_x).
\]
Finally, set
\[
\widetilde V
=
\begin{cases}
\displaystyle\Pi_{[0,1]}\left(\sum_{x\in[d]}T_x\right),
&\mathsf M\le n,\\[4pt]
0,&\mathsf M>n,
\end{cases}
\qquad
\widehat V_{n,d,\epsilon}^{\mathrm{AF}}
=
E[\widetilde V\mid O_1,\ldots,O_n].
\]
The conditional expectation integrates the auxiliary count and marks. Equivalently, on this branch,
\[
\widehat V_{n,d,\epsilon}^{\mathrm{AF}}
=
\sum_{k=0}^{n}
e^{-n/4}\frac{(n/4)^k}{k!}
\sum_{(\mathsf B_1,\ldots,\mathsf B_k)\in\{0,1\}^k}
2^{-k}\,
\Pi_{[0,1]}\left(\sum_{x\in[d]}T_x\right),
\]
where each summand uses the counts and cell contributions constructed from the first \(k\) observations and the displayed mark vector. This defines a deterministic measurable real-valued function of the observed sample.
\end{algorithmv}

The calibration proved in \cref{sec:deferred-proofs} takes
\[
H_0=1025,
\qquad
\kappa=\frac1{17952},
\qquad
D_0=2.
\]
With this choice, the admissible range \(d\ge2\) uses the saturated and active branches. The parameterized definition in \cref{obj:def:armwise-estimator} also records an empirical branch for a general cutoff.

In the bounded-alphabet branch, the full-sample cell count \leanref{sym:N_x}{\ensuremath{N_x}} supplies the empirical cell weight, while the arm count \leanref{sym:N_{ax}}{\ensuremath{N_{ax}}} and success count \leanref{sym:S_{ax}}{\ensuremath{S_{ax}}} supply the empirical outcome mean. The zero-count convention completes the definition for empty arms. In the saturated branch, the midpoint \(1/2\) has squared error at most \(1/4\) because the observed value lies in the unit interval. These branches address the regimes in which a fixed alphabet or a constant risk bound suffices.

The active branch first localizes each cell's four observed atom masses. The auxiliary count \leanref{sym:\mathsf M}{\ensuremath{\mathsf M}} selects an initial portion of the sample, and the fair marks \leanref{sym:\mathsf B_i}{\ensuremath{\mathsf B_i}} allocate its observations to pilot and evaluation counts. The pilot counts \leanref{sym:N'_{ay,x}}{\ensuremath{N'_{ay,x}}} determine a localization rectangle \leanref{sym:Q_x}{\ensuremath{Q_x}} with center \leanref{sym:b_x}{\ensuremath{b_x}}. Its radii combine a term that varies with the estimated atom mass and an additive term that remains positive when the pilot count is zero. Thus the polynomial construction has positive coordinate radii even in cells with empty pilot atoms.

Within this rectangle, the tensor Jackson polynomial \leanref{sym:P_{x,K_n}}{\ensuremath{P_{x,K_n}}} approximates the extension \(F_\epsilon\) defined in \cref{obj:def:armwise-extension}. The approximation degree \leanref{sym:K_n}{\ensuremath{K_n}} grows with the logarithmic alphabet scale. The evaluation counts \leanref{sym:N_{ay,x}}{\ensuremath{N_{ay,x}}} then enter through the centered factorial statistic \leanref{sym:U_k(N;\zeta)}{\ensuremath{U_k(N;\zeta)}}, which replaces centered powers in the polynomial expansion by falling-factorial count terms. This produces the cell statistic \leanref{sym:Z_x}{\ensuremath{Z_x}}. The separation of localization and evaluation is motivated by Poisson splitting: in the auxiliary uncapped Poisson experiment, pilot and evaluation atom counts are independent, and falling factorials estimate powers of atom masses. The displayed cap incorporates this construction into a sample of fixed size.

Clipping adapts to the sensitivity of the cell contribution in each treatment arm. The armwise weights \leanref{sym:w_{a,x}}{\ensuremath{w_{a,x}}} compare total cell mass with the overlap-truncated arm mass at the pilot center. They enter the clipping scale \leanref{sym:S_x^{\mathrm{aw}}}{\ensuremath{S_x^{\mathrm{aw}}}} together with the arm-specific localization widths. Each clipped contribution \leanref{sym:T_x}{\ensuremath{T_x}} consequently remains within a prescribed interval around the center's contribution, and projection of their sum yields the bounded auxiliary statistic \leanref{sym:\widetilde V}{\ensuremath{\widetilde V}}.

Finally, conditional averaging integrates the auxiliary count and every mark allocation with their specified probabilities. The finite sum in \cref{obj:def:armwise-estimator} makes the resulting dependence explicit: for every \(k\le n\), it enumerates \(2^k\) mark vectors, and each summand also requires the four-dimensional Jackson integrals or their polynomial coefficients. The formula defines a deterministic theoretical statistic through finite averaging over auxiliary counts and mark allocations. For any fixed target value, conditional averaging cannot increase squared-error risk, by conditional Jensen's inequality. The uniform guarantee for this statistic is as follows.

\begin{theoremv}{thm:observable-upper}[Uniform observable risk bound]
There exist universal constants \(H_0>0\), \(\kappa>0\), \(0<C_\star<\infty\), and an integer \(D_0\ge2\) such that, for every choice of:
\begin{itemize}
\item \textbf{(Sample size.)} An integer \(n\ge1\).
\item \textbf{(Alphabet size.)} An integer \(d\ge2\).
\item \textbf{(Overlap.)} A public overlap parameter \(0<\epsilon\le1/2\).
\end{itemize}
the estimator \(\widehat V_{n,d,\epsilon}^{\mathrm{AF}}\) specified in \cref{obj:def:armwise-estimator}, with these tuning constants, is a measurable, deterministic function of the observed sample. For every family of sampling laws satisfying \cref{obj:ass:iid-sampling} for each \(\mathbb P\in\mathcal D_{d,\epsilon}^{\mathrm{obs}}\), the class of admissible observed laws with public overlap floor \(\epsilon\), its squared-error risk for the observed value \(\Psi(\mathbb P)\) defined in \cref{obj:def:observed-value} satisfies
\[
\sup_{\mathbb P\in\mathcal D_{d,\epsilon}^{\mathrm{obs}}}
E_{\mathbb P}\!\left[
\left\{\widehat V_{n,d,\epsilon}^{\mathrm{AF}}-\Psi(\mathbb P)\right\}^{2}
\right]
\le C_\star\min\left\{1,\frac{d}{n\epsilon L_d}\right\},
\qquad L_d=\log(ed).
\]
\end{theoremv}

The armwise sensitivity bound in \cref{obj:lem:armwise-modulus} explains the scaling of this guarantee. Changes in the four atom masses affect the cell contribution through the mass available in each arm, so localization and clipping retain separate armwise sensitivity weights. Polynomial approximation addresses the maximum over arm-specific means within the pilot rectangle, and factorial estimation translates that approximation into observed count statistics. Together, these ingredients yield the logarithmic improvement in the risk bound while accounting for the public overlap floor. \Cref{sec:appendix-upper} organizes the analytic ingredients, and \cref{sec:deferred-proofs} gives the detailed proof and calibration. Combined with \cref{obj:thm:matched-frontier}, the guarantee establishes attainment of the minimax rate by the deterministic statistic in \cref{obj:def:armwise-estimator}.


\section{Discussion and open questions}

The consistency criterion in \cref{obj:thm:consistency-threshold} links sample size, categorical support, and overlap through a single joint requirement. With \(n\) observed units, \(d\) covariate cells, overlap floor \(\epsilon\), and logarithmic alphabet scale \(L_d\), consistency requires
\[
\frac{d}{n\epsilon L_d}\longrightarrow 0.
\]
Thus, sample growth must compensate jointly for an expanding covariate alphabet and declining overlap. The criterion makes this tradeoff explicit: the relevant comparison is between the number of cells and sample size multiplied by the overlap floor and logarithmic alphabet scale.

For categorical treatment-choice problems, the minimax comparison in \cref{obj:thm:matched-frontier} provides a benchmark for assessing the precision of optimal-value estimation. In the nonsaturated regime, the squared-error scale is \(d/(n\epsilon L_d)\), up to universal comparison constants. This scale supports assessments of how changes in sample size, covariate resolution, and overlap affect the smallest achievable worst-case risk. The guarantee in \cref{obj:thm:observable-upper} supplies an observed-data estimator attaining that scale, linking the precision benchmark to the construction in \cref{obj:def:armwise-estimator}.

The benchmark also accommodates heterogeneity in treatment assignment and indifference between treatment arms. Within the class in \cref{obj:def:observed-class}, propensities may differ across occupied cells while satisfying the public overlap bounds. The value in \cref{obj:def:observed-value} uses the larger conditional outcome mean in each cell, including cells where the two means coincide. Such ties leave the value well defined even when either treatment attains the same cellwise optimum. Under the consistency and conditional exchangeability conditions in \cref{obj:ass:consistency,obj:ass:conditional-exchangeability}, \cref{obj:prop:identification} gives this observed-law benchmark its causal interpretation.

\subsection{Limitations and future work}

The attaining estimator is explicit as a measurable finite-sample statistic, but its displayed conditional average enumerates \(2^k\) mark assignments for each auxiliary count \(k\le n\), in addition to evaluating four-dimensional Jackson integrals or precomputed polynomial coefficients. The result therefore establishes statistical attainment rather than an efficient numerical algorithm. Finding a polynomial-time approximation with the same uniform risk guarantee, and quantifying the numerical error permitted without changing the rate, remain open computational questions.

The proved rate comparison determines the order of minimax risk. Leading-constant convergence is a separate question, particularly when the covariate alphabet grows and overlap vanishes together. The following question isolates that issue within the nonsaturated regime.

\begin{remarkv}{oeq:sharp-path-constant}[Open questions on path constants]
Consider admissible nonsaturated sequences satisfying
\begin{itemize}
\item \textbf{(Growing alphabet.)} \(d\to\infty\).
\item \textbf{(Vanishing overlap.)} \(\epsilon\to0\).
\item \textbf{(Nonsaturation.)} \(d/(n\epsilon L_d)\to0\), where \(L_d\) is the logarithmic alphabet scale.
\end{itemize}
For \(\mathfrak R_{n,d,\epsilon}\), the minimax squared-error risk, consider the normalized sequence
\[
\frac{n\epsilon L_d}{d}\,\mathfrak R_{n,d,\epsilon}.
\]
It remains open whether there exists a single finite positive constant to which the entire normalized-risk sequence converges along every admissible path. Alternatively, do there exist two admissible overlap paths along which the entire normalized-risk sequences converge to distinct finite positive constants? Failure of convergence along an admissible path is a further possibility.

Future work could investigate a cone-constrained approximation--moment saddle for the boundary-propensity likelihood and an observable no-loss transfer.
\end{remarkv}

The nonsaturation condition places these sequences in the region where the minimax risk tends to zero. Along them, \cref{obj:thm:matched-frontier} bounds the normalized risk above and below by universal positive constants. Convergence to a common constant would sharpen that comparison into a path-independent calibration of leading risk. Distinct limits along two admissible paths would instead reveal dependence on how sample size, support, and overlap evolve jointly. The alternatives in \cref{obj:oeq:sharp-path-constant} concern convergence of entire sequences; convergence along selected subsequences would leave the stated question unresolved.

The proposed approximation--moment direction seeks a sharp characterization compatible with the boundary-propensity likelihood and the constrained prior pairs in \cref{obj:def:constrained-prior-separation}. An observable no-loss transfer would aim to carry the resulting leading constant into the observable statistical problem without deterioration. Developing these two steps could connect sharp approximation and moment calculations to a leading-risk calibration beyond the established rate comparison.


\appendix

\section*{Appendices}

\section{Identification and upper-bound proofs}\label{sec:appendix-upper}

This appendix collects the analytic ingredients for \cref{obj:prop:identification,obj:thm:observable-upper}. The causal argument connects the observed-law class to admissible full-data completions. The upper-bound analysis uses armwise sensitivity to organize localization, polynomial approximation, factorial estimation, and clipping. Throughout, the estimator and its sampling conventions are those of \cref{obj:def:armwise-estimator}; the detailed arguments appear in \cref{sec:deferred-proofs}.

\subsection{Causal completion and value identification}

The two causal clauses of \cref{obj:prop:identification} have distinct mathematical roles. Existence of a completion ensures that every law in \cref{obj:def:observed-class} has a causal interpretation satisfying \cref{obj:ass:consistency,obj:ass:conditional-exchangeability}. Value identification then establishes that all completions in \cref{obj:def:causal-class} yield the observed value of their marginal law. Together, these clauses connect the observed statistical problem to the causal oracle value under the stated overlap restriction.

Positive overlap supplies both conditional outcome means in every occupied cell. Consistency and conditional exchangeability relate these observed means to the corresponding conditional potential-outcome means. The oracle can choose an arm separately in each cell, so its value takes the cell-mass-weighted maximum appearing in \cref{obj:def:observed-value}. Cells of zero probability contribute zero under the conventions of \cref{obj:prop:identification}.

The final clause of \cref{obj:prop:identification} supplies the cellwise representation used in estimation. At an admissible vector of observed atom probabilities, the overlap-truncated denominators in \cref{obj:def:armwise-extension} agree with the arm masses. The extension therefore represents the target contribution of that cell, including its zero-mass case. Its domain also accommodates the nonnegative vectors generated by pilot localization.

\subsection{Armwise sensitivity}

The local analysis retains the sensitivity of each treatment arm separately. The following bounds control perturbations of the extension and relate the resulting armwise weights to the observed atom masses.

\begin{lemmav}{lem:armwise-modulus}[Armwise perturbation and mass bounds]
Let \(\epsilon\) be the overlap parameter, with \(0<\epsilon\leq 1/2\). For every pair of four-vectors \(u,v\) satisfying
\begin{itemize}
\item \textbf{(Nonnegativity.)} \(u,v\in\mathbb R_+^4\).
\item \textbf{(Nonzero comparison.)} \(v\neq 0\).
\end{itemize}
define the armwise perturbations by
\[
\Delta_a=\sum_{y=0}^1|u_{ay}-v_{ay}|,\qquad a\in\{0,1\}.
\]
With \(F_\epsilon\), \(s\), and \(D_a\) as in \cref{obj:def:armwise-extension},
\[
|F_\epsilon(u)-F_\epsilon(v)|
\leq
2\sum_{a=0}^1\Delta_a
+
2\sum_{a=0}^1\frac{s(v)}{D_a(v)}\Delta_a.
\]

For every observed law \(\mathbb P\) on
\(\{0,\ldots,d-1\}\times\{0,1\}\times\{0,1\}\) and every cell \(x\in\{0,\ldots,d-1\}\), let \(v=q_x\) be its vector of observed atom probabilities and \(p_x\) its cell probability:
\[
v_{ay}=\mathbb P\bigl(\{(x,a,y)\}\bigr),
\qquad
p_x=\sum_{a=0}^1\sum_{y=0}^1v_{ay}.
\]
If
\begin{itemize}
\item \textbf{(Dimension.)} \(d\) is an integer with \(d\geq2\).
\item \textbf{(Observed class.)} \(\mathbb P\in\mathcal D_{d,\epsilon}^{\mathrm{obs}}\), as in \cref{obj:def:observed-class}.
\item \textbf{(Positive cell mass.)} \(p_x>0\).
\end{itemize}
then, with \(s_a\) as in \cref{obj:def:armwise-extension},
\[
\sum_{a=0}^1\frac{p_x}{s_a(v)}
\sum_{y=0}^1\sqrt{v_{ay}}
\leq
2\sqrt{2}\sqrt{\frac{p_x}{\epsilon}}.
\]
\end{lemmav}

\begin{proof}[Proof of \cref{obj:lem:armwise-modulus}]
\emph{1. Mass and denominator perturbations.}
Use the quantities in \cref{obj:def:armwise-extension}, and define
\[
\delta:=\sum_{a=0}^1\Delta_a.
\]
Each \(\Delta_a\) is nonnegative, so \(\delta\geq0\). The triangle inequality gives, for each \(a\in\{0,1\}\),
\[
\begin{aligned}
|s_a(u)-s_a(v)|
&=\bigl|(u_{a0}-v_{a0})+(u_{a1}-v_{a1})\bigr|
\leq\Delta_a,\\
|s(u)-s(v)|
&\leq |s_0(u)-s_0(v)|+|s_1(u)-s_1(v)|
\leq\delta.
\end{aligned}
\]
Because the coordinates of \(v\) are nonnegative, \(s(v)\geq0\); equality would force every coordinate to vanish. Thus \(v\neq0\) implies
\[
s(v)>0,\qquad
D_a(v)\geq\epsilon s(v)>0,\qquad
\frac{s(v)}{D_a(v)}\geq0.
\]
Comparing the two entries of each maximum yields
\[
\begin{aligned}
|D_a(u)-D_a(v)|
&\leq
\max\bigl\{|s_a(u)-s_a(v)|,\,
          |\epsilon s(u)-\epsilon s(v)|\bigr\}\\
&\leq\max\{\Delta_a,\epsilon\delta\}
\leq\Delta_a+\epsilon\delta.
\end{aligned}
\]
Here \(\epsilon>0\) gives
\(
|\epsilon s(u)-\epsilon s(v)|
=\epsilon|s(u)-s(v)|\leq\epsilon\delta
\).
% lean: CausalSmith.Stat.OptvalueVanishingoverlapRate.armwise_modulus_extension

\emph{2. The zero-vector case.}
Suppose \(u=0\). Nonnegativity and the definition of the denominator imply
\[
0\leq v_{a1}\leq s_a(v)\leq D_a(v),
\qquad
0\leq\frac{s(v)v_{a1}}{D_a(v)}\leq s(v).
\]
Taking the maximum over the two arms gives
\[
0\leq F_\epsilon(v)\leq s(v).
\]
Also,
\[
s(v)=|s(u)-s(v)|\leq\delta,
\qquad
\sum_{a=0}^1\frac{s(v)}{D_a(v)}\Delta_a\geq0.
\]
Since \(F_\epsilon(u)=0\), it follows that
\[
|F_\epsilon(u)-F_\epsilon(v)|
=F_\epsilon(v)
\leq\delta
\leq2\delta+
2\sum_{a=0}^1\frac{s(v)}{D_a(v)}\Delta_a.
\]
% lean: CausalSmith.Stat.OptvalueVanishingoverlapRate.armwiseExtension_nonneg_le_totalMassVec

\emph{3. A coordinate-ratio estimate.}
Now suppose \(u\neq0\). The same nonnegative-coordinate argument gives \(s(u)>0\), and hence, for each arm,
\[
D_a(u)\geq\epsilon s(u)>0,\qquad
0\leq u_{a1}\leq D_a(u),\qquad
0\leq v_{a1}\leq D_a(v).
\]
We first establish
\[
\left|\frac{u_{a1}}{D_a(u)}-\frac{v_{a1}}{D_a(v)}\right|
\leq
\frac{|u_{a1}-v_{a1}|+|D_a(u)-D_a(v)|}{D_a(v)}.
\]
If \(D_a(u)\leq D_a(v)\), the identity
\[
D_a(v)u_{a1}-D_a(u)v_{a1}
=
D_a(u)(u_{a1}-v_{a1})
+(D_a(v)-D_a(u))u_{a1}
\]
and \(0\leq u_{a1}\leq D_a(u)\) give
\[
\begin{aligned}
|D_a(v)u_{a1}-D_a(u)v_{a1}|
&\leq D_a(u)|u_{a1}-v_{a1}|
 +(D_a(v)-D_a(u))u_{a1}\\
&\leq D_a(u)
 \bigl(|u_{a1}-v_{a1}|+|D_a(u)-D_a(v)|\bigr).
\end{aligned}
\]
If \(D_a(v)\leq D_a(u)\), use instead
\[
D_a(v)u_{a1}-D_a(u)v_{a1}
=
D_a(v)(u_{a1}-v_{a1})
-(D_a(u)-D_a(v))v_{a1}.
\]
Since \(0\leq v_{a1}\leq D_a(v)\leq D_a(u)\),
\[
\begin{aligned}
|D_a(v)u_{a1}-D_a(u)v_{a1}|
&\leq D_a(v)|u_{a1}-v_{a1}|
 +(D_a(u)-D_a(v))v_{a1}\\
&\leq D_a(u)
 \bigl(|u_{a1}-v_{a1}|+|D_a(u)-D_a(v)|\bigr).
\end{aligned}
\]
In both cases, divide by the positive product \(D_a(u)D_a(v)\), using
\[
\frac{u_{a1}}{D_a(u)}-\frac{v_{a1}}{D_a(v)}
=
\frac{D_a(v)u_{a1}-D_a(u)v_{a1}}{D_a(u)D_a(v)},
\]
to obtain the coordinate-ratio estimate.
% lean: CausalSmith.Stat.OptvalueVanishingoverlapRate.armwise_ratio_bound

\emph{4. Scaling the ratios and comparing the maxima.}
For the two nonzero vectors under consideration, define
\[
g_a(\xi):=\frac{s(\xi)\xi_{a1}}{D_a(\xi)},
\qquad \xi\in\{u,v\},\quad a\in\{0,1\}.
\]
The identity
\[
g_a(u)-g_a(v)
=
(s(u)-s(v))\frac{u_{a1}}{D_a(u)}
+s(v)\left(\frac{u_{a1}}{D_a(u)}
           -\frac{v_{a1}}{D_a(v)}\right)
\]
gives
\[
\begin{aligned}
|g_a(u)-g_a(v)|
&\leq |s(u)-s(v)|\frac{u_{a1}}{D_a(u)}
+s(v)\left|\frac{u_{a1}}{D_a(u)}
           -\frac{v_{a1}}{D_a(v)}\right|\\
&\leq\delta+
\frac{s(v)}{D_a(v)}
\bigl(\Delta_a+|D_a(u)-D_a(v)|\bigr).
\end{aligned}
\]
Indeed, \(0\leq u_{a1}/D_a(u)\leq1\), and
\(
|u_{a1}-v_{a1}|\leq\Delta_a
\).
% lean: CausalSmith.Stat.OptvalueVanishingoverlapRate.armwise_coordinate_ratio_bound
Moreover,
\[
\epsilon\frac{s(v)}{D_a(v)}
=\frac{\epsilon s(v)}{D_a(v)}\leq1.
\]
Substituting the denominator perturbation bound therefore yields
\[
\begin{aligned}
|g_a(u)-g_a(v)|
&\leq\delta+
\frac{s(v)}{D_a(v)}(2\Delta_a+\epsilon\delta)\\
&\leq2\delta+2\frac{s(v)}{D_a(v)}\Delta_a.
\end{aligned}
\]
Finally, comparison of the two entries of the maximum gives
\[
\begin{aligned}
|F_\epsilon(u)-F_\epsilon(v)|
&=|\max\{g_0(u),g_1(u)\}-\max\{g_0(v),g_1(v)\}|\\
&\leq\max_{a\in\{0,1\}}|g_a(u)-g_a(v)|\\
&\leq\max_{a\in\{0,1\}}
 \left\{2\delta+2\frac{s(v)}{D_a(v)}\Delta_a\right\}\\
&\leq2\delta+
2\sum_{a=0}^1\frac{s(v)}{D_a(v)}\Delta_a.
\end{aligned}
\]
The last inequality uses the nonnegativity of both weighted perturbations. This proves the first assertion.
% lean: CausalSmith.Stat.OptvalueVanishingoverlapRate.armwise_modulus_extension

\emph{5. Observed arm-mass floors.}
For the second assertion, fix the stated law and cell, and take \(v=q_x\). The atom probabilities are nonnegative. The definitions in \cref{obj:def:observed-class,obj:def:armwise-extension}, together with \(p_x>0\), give
\[
p_x=s_0(v)+s_1(v),\qquad
p_x\pi_x=s_1(v),\qquad
p_x(1-\pi_x)=s_0(v).
\]
The last identity follows by subtracting the second from the first. Membership in the observed class supplies \cref{obj:ass:shrinking-overlap}, so
\[
\epsilon\leq\pi_x,\qquad \epsilon\leq1-\pi_x.
\]
Multiplication by \(p_x>0\) gives, separately for the two arms,
\[
\epsilon p_x\leq p_x(1-\pi_x)=s_0(v),
\qquad
\epsilon p_x\leq p_x\pi_x=s_1(v).
\]
Consequently,
\[
s_a(v)\geq\epsilon p_x>0,\qquad a\in\{0,1\}.
\]
% lean: CausalSmith.Stat.OptvalueVanishingoverlapRate.observed_armwise_sqrt_bound

\emph{6. The weighted square-root estimate.}
For each arm, nonnegativity of its two coordinates gives
\[
\begin{aligned}
2s_a(v)-(\sqrt{v_{a0}}+\sqrt{v_{a1}})^2
&=(\sqrt{v_{a0}}-\sqrt{v_{a1}})^2\geq0,\\
(\sqrt{v_{a0}}+\sqrt{v_{a1}})^2
&\leq2s_a(v).
\end{aligned}
\]
Both sides of the resulting square-root comparison are nonnegative, hence
\[
\sqrt{v_{a0}}+\sqrt{v_{a1}}\leq\sqrt{2s_a(v)}.
\]
% lean: CausalSmith.Stat.OptvalueVanishingoverlapRate.sqrt_pair_le_sqrt_twice_sum
Multiplying the squared bound by \(p_x^2\), and then using the arm-mass floor, yields
\[
\begin{aligned}
\bigl(p_x(\sqrt{v_{a0}}+\sqrt{v_{a1}})\bigr)^2
&\leq2p_x^2s_a(v)\\
&=\frac{2p_xs_a(v)}{\epsilon}(\epsilon p_x)\\
&\leq\frac{2p_xs_a(v)}{\epsilon}s_a(v)\\
&=\frac{2p_x}{\epsilon}s_a(v)^2\\
&=\left(\sqrt{\frac{2p_x}{\epsilon}}\,s_a(v)\right)^2.
\end{aligned}
\]
The compared quantities are nonnegative. Taking square roots and dividing by \(s_a(v)>0\) proves
\[
\frac{p_x}{s_a(v)}
\bigl(\sqrt{v_{a0}}+\sqrt{v_{a1}}\bigr)
\leq\sqrt{\frac{2p_x}{\epsilon}},
\qquad a\in\{0,1\}.
\]
% lean: CausalSmith.Stat.OptvalueVanishingoverlapRate.weighted_sqrt_pair_bound

\emph{7. Summing the two arms.}
Adding the two armwise bounds and factoring the square root gives
\[
\begin{aligned}
\sum_{a=0}^1\frac{p_x}{s_a(v)}
\sum_{y=0}^1\sqrt{v_{ay}}
&\leq2\sqrt{\frac{2p_x}{\epsilon}}\\
&=2\sqrt{2}\sqrt{\frac{p_x}{\epsilon}}.
\end{aligned}
\]
This proves the second assertion and completes the proof.
% lean: CausalSmith.Stat.OptvalueVanishingoverlapRate.armwise_modulus
\end{proof}

The armwise perturbation \leanref{sym:\Delta_a}{\ensuremath{\Delta_a}} measures the total coordinate change within one treatment arm. Nonnegativity preserves the mass interpretation of the extension, while a nonzero comparison vector makes each overlap-truncated denominator positive. The first inequality applies even when the perturbed vector is zero, which is useful for localized estimation near empty cells.

The second inequality concerns occupied cells in the observed-law class. It combines the armwise sensitivity weights with square roots of atom masses, the quantities governing count fluctuations. Its dependence on overlap is proportional to the inverse square root of the overlap floor. This mass-sensitive control is the analytic ingredient supporting the single inverse power of overlap in the squared-risk guarantee of \cref{obj:thm:observable-upper}.

\subsection{Factorial statistics and sampling conventions}

Polynomial estimation requires statistics corresponding to powers of atom probabilities. A deterministic split gives a multinomial construction, which we record before discussing its relationship to the active estimator's Poisson construction.

\begin{definitionv}{def:fixed-sample-factorial}[Centered factorial statistic \(U_{\boldsymbol h}\)]
The centered multinomial factorial statistic is
\[
U_{\boldsymbol h}(N^{\mathrm e};b)
=
\sum_{\boldsymbol t\le\boldsymbol h}
\left\{
\prod_{j\in\mathcal J_d}
\binom{h_j}{t_j}(-b_j)^{h_j-t_j}
\right\}
\frac{\prod_{j\in\mathcal J_d}(N_j^{\mathrm e})_{t_j}}
{(n_{\mathrm e})_{|\boldsymbol t|}},
\qquad
|\boldsymbol h|\le n_{\mathrm e}.
\]
Here \(\mathcal J_d\) is the observed atom index set, \(b=(b_j)_{j\in\mathcal J_d}\) is the centering vector, and \(\boldsymbol h,\boldsymbol t\) are nonnegative integer multi-indices, with componentwise ordering and total orders
\[
|\boldsymbol h|=\sum_{j\in\mathcal J_d}h_j,
\qquad
|\boldsymbol t|=\sum_{j\in\mathcal J_d}t_j.
\]
The deterministic sample split and its atom counts are
\[
n_{\mathrm p}=\lfloor n/2\rfloor,
\qquad
n_{\mathrm e}=n-n_{\mathrm p},
\qquad
I_{\mathrm p}=\{1,\ldots,n_{\mathrm p}\},
\qquad
I_{\mathrm e}=\{n_{\mathrm p}+1,\ldots,n\},
\]
\[
N_j^{\mathrm p}
=\sum_{i\in I_{\mathrm p}}\mathbf 1\{O_i=j\},
\qquad
N_j^{\mathrm e}
=\sum_{i\in I_{\mathrm e}}\mathbf 1\{O_i=j\},
\qquad
N^{\mathrm e}=(N_j^{\mathrm e})_{j\in\mathcal J_d}.
\]
The falling factorials are
\[
(m)_k=m(m-1)\cdots(m-k+1),
\qquad
(m)_0=1.
\]
\end{definitionv}

The centered multinomial statistic \leanref{sym:U_{\boldsymbol h}(N^{\mathrm e};b)}{\ensuremath{U_{\boldsymbol h}(N^{\mathrm e};b)}} uses the joint factorial moments of a fixed-size evaluation histogram. Its denominator depends on the total order of each term, reflecting the common multinomial sample size. The order restriction ensures that these normalizations are defined. Under \cref{obj:ass:iid-sampling}, the two deterministic subsamples are independent, allowing the centering vector to be selected from the pilot sample.

The active branch of \cref{obj:def:armwise-estimator} uses a separate sampling convention: an auxiliary Poisson count and fair marks generate pilot and evaluation counts. In the uncapped Poisson experiment, the evaluation atom counts are independent Poisson variables. The coordinatewise statistic defined there consequently normalizes each falling factorial by a power of the evaluation intensity. Both constructions translate centered polynomial terms into count statistics, with their respective multinomial and Poisson normalizations. The risk guarantee in \cref{obj:thm:observable-upper} concerns the capped Poisson construction specified in \cref{obj:def:armwise-estimator}.

\subsection{Risk analysis by estimator branch}

The branches of \cref{obj:def:armwise-estimator} separate the roles of empirical estimation, bounded loss, and local polynomial estimation. For a bounded alphabet, the relevant quantities are empirical cell weights and arm-specific outcome means, including the prescribed zero-count convention. Since the alphabet cutoff is universal, the logarithmic alphabet scale remains bounded in this branch. Its contribution to \cref{obj:thm:observable-upper} is therefore assessed within a fixed range of alphabet sizes.

For larger alphabets in the saturated regime, the estimator returns the midpoint of the unit interval. The target lies in that interval by \cref{obj:def:observed-value}, giving a squared-error bound of one quarter. This branch supplies the constant-risk part of \cref{obj:thm:observable-upper}.

The active branch assigns separate mathematical roles to localization, approximation, factorial variance, and clipping. Pilot localization determines a rectangle for each cell's atom probabilities. Its radii include both a mass-dependent fluctuation term and a positive additive term, so the polynomial coordinates remain defined when pilot counts vanish. The perturbation bound in \cref{obj:lem:armwise-modulus} measures variation over these rectangles through the armwise widths and sensitivities.

Polynomial bias concerns the discrepancy between the extension and its tensor Jackson approximation on the pilot rectangle. The approximation degree grows with the logarithmic alphabet scale, while localization determines the region on which the approximation is used. Factorial estimation then represents the polynomial through evaluation counts. Independence in the uncapped Poisson experiment separates the pilot choice of polynomial from evaluation noise; factorial moments govern the variance cost of estimating its centered terms. The tuning constants in \cref{obj:thm:observable-upper} calibrate these approximation and variance roles jointly.

Clipping uses the same armwise geometry. The interval around the center's contribution is scaled by the armwise localization widths and sensitivity weights specified in \cref{obj:def:armwise-estimator}. It bounds each polynomial statistic before aggregation. Projection of the aggregate onto the unit interval aligns its range with that of the target and cannot increase squared error for a target in that interval.

Finally, the Poisson cap connects the auxiliary experiment to the available sample of exactly \(n\) observations. On the cap event, the statistic uses the selected initial observations and their marks; overflow receives the prescribed value zero. The bounded range permits the overflow probability to enter the risk accounting as a controlled discrepancy. Conditional averaging over the auxiliary count and marks then produces the deterministic observed-data estimator. With the proved calibration \(D_0=2\) and the admissible range \(d\ge2\), \cref{obj:thm:observable-upper} supplies the uniform risk bound through the saturated and active branches. The empirical branch belongs to the general parameterized definition and is inactive under this calibration.


\section{Weighted approximation and constrained moment priors}\label{sec:appendix-weighted}

The observable upper bound in \cref{obj:thm:observable-upper} combines polynomial approximation with factorial estimation. The complementary lower-bound argument requires probability priors whose moments agree while their target values remain separated. This appendix presents the \leanref{S-4}{scalar approximation problem} underlying that construction, the localization bounds used to obtain the priors, and the resulting comparison between approximation error and constrained prior separation.

\subsection{The scalar functional and probability constraints}

The scalar target combines a change in slope at intensity one with the overlap-dependent weighting. We introduce its formula before defining the approximation problem.



The weighted scalar positive-part functional \leanref{sym:\phi_\epsilon(z)}{\ensuremath{\phi_\epsilon(z)}} is
\[
\phi_\epsilon(z)
=\frac{\{(1-\epsilon)+\epsilon z\}\max\{z-1,0\}}{1+z}.
\]
It is continuous and vanishes below intensity one. Its change in slope supplies the local approximation difficulty. The polynomial degree \leanref{sym:K}{\ensuremath{K}} and support endpoint \leanref{sym:M}{\ensuremath{M}} determine the approximation scale. Weighting the uniform error by \(1+z\), where \leanref{sym:z}{\ensuremath{z}} is the scalar intensity, also aligns the approximation criterion with the common-mean restriction on the priors.

\begin{definitionv}{def:weighted-error}[Weighted approximation error \(\mathcal E_{K,M,\epsilon}^{w}\)]
The weighted best uniform polynomial approximation error is
\[
\mathcal E_{K,M,\epsilon}^{w}
:=
\inf_{P\in\mathbb R_K[z]}
\sup_{0\le z\le M}
\frac{|\phi_\epsilon(z)-P(z)|}{1+z},
\]
where \(\phi_\epsilon\) is the weighted scalar positive-part functional and \(\mathbb R_K[z]\) is the set of real polynomials in \(z\) of degree at most \(K\).
\end{definitionv}

The weighted approximation error \leanref{sym:\mathcal E_{K,M,\epsilon}^{w}}{\ensuremath{\mathcal E_{K,M,\epsilon}^{w}}} measures the discrepancy after allowing the pointwise error to grow with intensity. For probability measures of mean one, the weight \(1+z\) has integral two. This normalization gives approximation error a direct role in bounding target differences between moment-matched priors. The admissible probability pairs are defined next.

\begin{definitionv}{def:constrained-prior-separation}[Constrained prior separation \(\Delta_{K,M,\epsilon}\)]
The normalized constrained prior separation is
\[
\Delta_{K,M,\epsilon}
=
\sup_{(\nu_0,\nu_1)\in\mathfrak N_{K,M}}
\left|
\int_0^M\phi_\epsilon(z)\,d(\nu_1-\nu_0)(z)
\right|,
\]
where \(\phi_\epsilon\) is the weighted scalar positive-part functional and the parameters satisfy
\begin{itemize}
\item \textbf{(Degree.)} \(K\) is an integer with \(K\geq2\).
\item \textbf{(Support endpoint.)} \(M\) is real with \(2\leq M\leq K^2\).
\item \textbf{(Overlap.)} \(0\leq\epsilon\leq1/2\).
\end{itemize}
The set \(\mathfrak N_{K,M}\) consists of pairs \((\nu_0,\nu_1)\) of probability measures on \([0,M]\) satisfying
\[
\int_0^M z\,d\nu_0(z)
=
\int_0^M z\,d\nu_1(z)
=1
\]
and
\[
\int_0^M z^j\,d\nu_0(z)
=
\int_0^M z^j\,d\nu_1(z),
\qquad j=0,1,\ldots,K.
\]
\end{definitionv}

The constrained prior separation \leanref{sym:\Delta_{K,M,\epsilon}}{\ensuremath{\Delta_{K,M,\epsilon}}} therefore incorporates both moment agreement and an absolute normalization of the first moment. Moment agreement cancels polynomial target components through the prescribed degree, while the common mean fixes the average intensity. The parameter range in \cref{obj:def:constrained-prior-separation} is the range covered by the comparison below, including zero overlap for this scalar approximation problem.

A signed measure supplies the target contrast; a common atom completes its two nonnegative parts into probability priors. We define that completion as a map on measures. The support and normalization properties needed for the constructed packet inputs are established in \cref{obj:thm:weighted-separation-and-frontier}.

\begin{definitionv}{def:weighted-handle}[Common-atom completion map \(\mathscr H^w\)]
The common-atom completion map is
\[
\mathscr H^w(\sigma_+,\sigma_-)
=
\left(
\sigma_-+(1-t)_+\delta_{z_0},
\ \sigma_++(1-t)_+\delta_{z_0}
\right),
\]
with
\[
t=\sigma_+(\mathbb R),
\qquad
u=\int_{\mathbb R}z\,d\sigma_+(z),
\qquad
z_0=\frac{1-u}{1-t}.
\]
Its inputs satisfy the following conditions:
\begin{itemize}
\item \textbf{(Measures.)} \(\sigma_+\) and \(\sigma_-\) are finite nonnegative measures on \(\mathbb R\).
\item \textbf{(First moments.)} Each measure has a finite first moment.
\item \textbf{(Mass.)} \(t<1\).
\end{itemize}
Here \((r)_+=\max\{r,0\}\), and \(\delta_{z_0}\) is the unit point mass at \(z_0\).
\end{definitionv}

The common-atom completion map \leanref{sym:\mathscr H^{w}}{\ensuremath{\mathscr H^w}} adds the same measure to both inputs, preserving their signed difference. Its location uses the positive input's remaining mass and first moment. For the packet inputs in \cref{obj:thm:weighted-separation-and-frontier}, the stated bounds place this atom inside the intensity support and certify that both completed measures have total mass and mean one.

\subsection{Fourier coefficients and packet localization}

A modulated Jackson function provides a localized signed contrast with a prescribed frequency band. Its twice antiderivative separates the response at the change in slope from the contribution spread across the interval. The following grouped notation fixes the periodic normalization, coefficient conventions, and endpoint extension before stating the localization bounds.



The shifted squared frequencies in \cref{obj:lem:jackson-packet-localization} encode two integrations of the modulated function. The endpoint calculation concerns the coefficients after extension by zero: forward differences crossing either end of the coefficient band enter the same sum as interior differences. This convention gives a bound for the complete Fourier expression. The localization statement records that coefficient bound together with frequency cancellation, the value at the periodic origin, and spatial decay.

\begin{lemmav}{lem:jackson-packet-localization}[Jackson packet localization bounds]
There exist real constants \(0<c\leq C\) such that the following holds for every integer \(N\geq2\). Define the normalized periodic Jackson function by
\[
\mathcal J_N(u)
=c_N\left\{\frac{\sin(Nu/2)}{\sin(u/2)}\right\}^{4},
\qquad
c_N=
\left(
\int_{-\pi}^{\pi}
\left\{\frac{\sin(Nu/2)}{\sin(u/2)}\right\}^{4}
\frac{du}{2\pi}
\right)^{-1},
\]
with removable values filled continuously. Set
\[
L=4N,\qquad
k_N(u)=\mathcal J_N(u)\cos(Lu),
\qquad
\widehat h(j)=\int_{-\pi}^{\pi}h(u)e^{-iju}\,\frac{du}{2\pi}.
\]
For \(j\in\mathbb Z\), define the zero-extended coefficient sequences and their forward differences by
\[
a_j^\pm=
\begin{cases}
\widehat{\mathcal J_N}(j)/(L\pm j)^2,& |j|\leq2N-2,\\
0,& |j|>2N-2,
\end{cases}
\qquad
(\delta a)_j=a_{j+1}-a_j,
\qquad
(\delta^2a)_j=a_{j+2}-2a_{j+1}+a_j.
\]
Define also
\[
G_N(u)
=-\frac12\sum_{j=-(2N-2)}^{2N-2}
\left\{
a_j^+\cos((L+j)u)+a_j^-\cos((L-j)u)
\right\},
\qquad
\|u\|_{\mathbb T}
=\inf_{j\in\mathbb Z}|u-2\pi j|.
\]
Then
\[
\int_{-\pi}^{\pi}\mathcal J_N(u)\,\frac{du}{2\pi}=1.
\]
For each sign, the second forward differences are absolutely summable and satisfy
\[
\sum_{j\in\mathbb Z}|\delta^2a_j^\pm|
\leq\frac{C}{N^3}.
\]
The differences include the endpoints of the displayed zero extension. For every \(r\in\mathbb Z\) satisfying
\[
|r|<2N+2\quad\text{or}\quad |r|>6N-2,
\]
the trigonometric integrals vanish:
\[
\int_{-\pi}^{\pi}k_N(u)\cos(ru)\,du=0,
\qquad
\int_{-\pi}^{\pi}k_N(u)\sin(ru)\,du=0.
\]
The function \(G_N\) satisfies
\[
G_N''(u)=k_N(u)\quad\text{for every }u\in\mathbb R,
\qquad
\int_{-\pi}^{\pi}G_N(u)\,du=0,
\]
and the bounds
\[
\frac{c}{N}\leq |G_N(0)|\leq\frac{C}{N},
\qquad
|G_N(u)|\leq
\frac{C}{N\{1+N\|u\|_{\mathbb T}\}^{2}}
\quad\text{for every }u\in\mathbb R,
\qquad
\int_{-\pi}^{\pi}|G_N(u)|\,\frac{du}{2\pi}
\leq\frac{C}{N^2}.
\]
\end{lemmav}

\begin{proof}[Proof of \cref{obj:lem:jackson-packet-localization}]
\begin{enumerate}
\item We first establish the coefficient representation needed for the center estimate. For integers \(N\geq1\) and real \(u\), write
\[
\mathcal R_N(u)=
\begin{cases}
\left\{\dfrac{\sin(Nu/2)}{\sin(u/2)}\right\}^{4},
& \sin(u/2)\neq0,\\[1ex]
N^4,&\sin(u/2)=0.
\end{cases}
\]
The sine addition formulas give
\[
\sin^2((N+1)u/2)-\sin^2(Nu/2)
=\sin((2N+1)u/2)\sin(u/2),
\]
and, for every positive integer \(j\),
\[
\sin((2j+1)u/2)-\sin((2j-1)u/2)
=2\cos(ju)\sin(u/2).
\]
Summing the latter identity, then using the former, yields
\[
\sqrt{\mathcal R_{N+1}(u)}-\sqrt{\mathcal R_N(u)}
=1+2\sum_{j=1}^{N}\cos(ju)
\qquad\text{when }\sin(u/2)\neq0.
\]
When \(\sin(u/2)=0\), one has \(\cos(u/2)=1\) or \(-1\), and the same identity reads
\((N+1)^2-N^2=1+2N\). Starting from \(\mathcal R_1(u)=1\), induction therefore gives
\[
\sqrt{\mathcal R_N(u)}
=N+2\sum_{j=1}^{N-1}(N-j)\cos(ju).
\]

For \(N\geq2\) and \(j\in\mathbb Z\), define
\[
\tau_{N,j}=\max\{N-|j|,0\},\qquad
\gamma_{N,j}=\sum_{j_1\in\mathbb Z}
\tau_{N,j_1}\tau_{N,j-j_1},\qquad
\mathcal S_N=\{j\in\mathbb Z:|j|\leq2N-2\}.
\]
These sums are finite. The sequence \(\tau_{N,j}\) is even and vanishes for
\(|j|\geq N\), so \(\gamma_{N,j}\) is even, nonnegative, and vanishes outside
\(\mathcal S_N\). The preceding expansion becomes
\[
\sqrt{\mathcal R_N(u)}
=\sum_{j\in\mathbb Z}\tau_{N,j}\cos(ju).
\]
To square it, use
\[
\cos(j_1u)\cos(j_2u)
=\frac{\cos((j_1+j_2)u)+\cos((j_1-j_2)u)}2.
\]
Evenness makes the two resulting double sums equal after replacing
\(j_2\) by \(-j_2\). Reindexing their common value by \(j=j_1+j_2\) gives
\[
\mathcal R_N(u)
=\sum_{j_1,j_2\in\mathbb Z}
\tau_{N,j_1}\tau_{N,j_2}\cos((j_1+j_2)u)
=\sum_{j\in\mathcal S_N}\gamma_{N,j}\cos(ju).
\]
At frequency zero,
\[
\gamma_{N,0}
=N^2+2\sum_{j=1}^{N-1}(N-j)^2
=\frac{2N^3+N}{3}>0,
\]
where the last equality follows from the finite sum of squares. Integrating the cosine expansion consequently gives
\[
\int_{-\pi}^{\pi}\mathcal R_N(u)\,\frac{du}{2\pi}
=\gamma_{N,0},
\qquad
c_N=\gamma_{N,0}^{-1}.
\]
Define, for every \(j\in\mathbb Z\),
\[
f_{N,j}=\frac{\gamma_{N,j}}{\gamma_{N,0}}.
\]
Thus
\[
\mathcal J_N(u)=\sum_{j\in\mathcal S_N}f_{N,j}\cos(ju),
\qquad
f_{N,j}\geq0,\qquad
f_{N,-j}=f_{N,j},\qquad
f_{N,j}=0\quad(j\notin\mathcal S_N).
\]

For completeness, cosine orthogonality reads, for \(j_1,j\in\mathbb Z\),
\[
\int_{-\pi}^{\pi}\cos(j_1u)\cos(ju)\,\frac{du}{2\pi}
=
\begin{cases}
1,&j_1=j=0,\\
\frac12,&|j_1|=|j|\neq0,\\
0,&\text{otherwise}.
\end{cases}
\]
This follows by the cosine product identity and integration of the integer modes. The sine part of the Fourier integral vanishes because \(\mathcal J_N\) is even. Hence
\[
\widehat{\mathcal J_N}(j)
=
\begin{cases}
f_{N,0},&j=0,\\
\dfrac{f_{N,j}+f_{N,-j}}2,&j\neq0,
\end{cases}
=f_{N,j}.
\]
In particular, the zero-extended sequences in the statement satisfy
\[
a_j^\pm=
\begin{cases}
f_{N,j}(L\pm j)^{-2},&j\in\mathcal S_N,\\
0,&j\notin\mathcal S_N.
\end{cases}
\]
% lean: CausalSmith.Stat.OptvalueVanishingoverlapRate.kernelCoeff_eq_staged

\item We now obtain the center bounds. Summing the triangular sequence gives
\[
\sum_{j\in\mathbb Z}\tau_{N,j}
=N+2\sum_{j=1}^{N-1}(N-j)=N^2.
\]
Reindexing the finite convolution sum yields
\[
\sum_{j\in\mathbb Z}\gamma_{N,j}=N^4,
\qquad
\sum_{j\in\mathbb Z}|\gamma_{N,j}|
\leq
\sum_{j,j_1\in\mathbb Z}
\tau_{N,j_1}\tau_{N,j-j_1}
=N^4.
\]
Using the value of \(\gamma_{N,0}\), we obtain
\[
\frac N2
\leq
\sum_{j\in\mathcal S_N}f_{N,j}
=\frac{N^4}{\gamma_{N,0}}
\leq2N.
\]
Indeed, the lower inequality is equivalent to \(4N^2\geq1\), which holds for \(N\geq2\), while
\(N^4/\gamma_{N,0}\leq2N\) follows directly from
\(\gamma_{N,0}=(2N^3+N)/3\).

For \(j\in\mathcal S_N\), both shifted frequencies obey
\[
2N\leq L+j\leq6N,\qquad
2N\leq L-j\leq6N.
\]
Their reciprocal squares therefore satisfy
\[
\frac1{36N^2}\leq\frac1{(L\pm j)^2}\leq\frac1{4N^2},
\]
and nonnegativity of \(f_{N,j}\) gives
\[
\frac{f_{N,j}}{18N^2}
\leq a_j^++a_j^-
\leq\frac{f_{N,j}}{2N^2}.
\]
At \(u=0\), every cosine equals one and every coefficient is nonnegative, so
\[
|G_N(0)|=\frac12\sum_{j\in\mathcal S_N}(a_j^++a_j^-).
\]
Consequently,
\[
\frac1{72N}
=
\frac12\,\frac{N/2}{18N^2}
\leq |G_N(0)|
\leq
\frac12\,\frac{2N}{2N^2}
=\frac1{2N}.
\]
% lean: Causalean.Mathlib.Analysis.Approximation.Chebyshev.Jackson.packet_center_bounds

\item To prepare the pointwise estimate, we bound the full-line second differences of the weighted coefficients. Directly from the triangular sequence,
\[
(\delta\tau_N)_j=
\begin{cases}
1,&-N\leq j\leq-1,\\
-1,&0\leq j\leq N-1,\\
0,&\text{otherwise},
\end{cases}
\qquad
\sum_{j\in\mathbb Z}|(\delta\tau_N)_j|=2N.
\]
Shifting the finite convolution sums gives
\[
(\delta\gamma_N)_j
=\sum_{j_1\in\mathbb Z}
(\delta\tau_N)_{j_1}\tau_{N,j-j_1},
\qquad
(\delta^2\gamma_N)_j
=\sum_{j_1\in\mathbb Z}
(\delta\tau_N)_{j_1}(\delta\tau_N)_{j-j_1}.
\]
The triangle inequality, followed by reindexing \(j-j_1\), now yields
\[
\begin{aligned}
\sum_j|(\delta\gamma_N)_j|
&\leq
\sum_{j,j_1}|(\delta\tau_N)_{j_1}|\tau_{N,j-j_1}
=2N^3,\\
\sum_j|(\delta^2\gamma_N)_j|
&\leq
\sum_{j,j_1}|(\delta\tau_N)_{j_1}|
                 |(\delta\tau_N)_{j-j_1}|
=4N^2.
\end{aligned}
\]
Division by the positive constant \(\gamma_{N,0}\) commutes with differences. Thus
\[
\sum_j|f_{N,j}|\leq2N,\qquad
\sum_j|(\delta f_N)_j|
\leq\frac{2N^3}{\gamma_{N,0}}\leq4,\qquad
\sum_j|(\delta^2f_N)_j|
\leq\frac{4N^2}{\gamma_{N,0}}\leq\frac6N.
\]

For \(\varsigma\in\{-1,1\}\) and \(j\in\mathbb Z\), define
\[
\omega_j^{\varsigma}=L+\varsigma j,\qquad
\rho_j^{\varsigma}=
\begin{cases}
(\omega_j^{\varsigma})^{-2},&\omega_j^{\varsigma}\neq0,\\
0,&\omega_j^{\varsigma}=0,
\end{cases}
\qquad
a_j^{(\varsigma)}=
\begin{cases}
a_j^+,&\varsigma=1,\\
a_j^-,&\varsigma=-1.
\end{cases}
\]
The convention at a zero denominator makes \(\rho^{\varsigma}\) a sequence on all of \(\mathbb Z\). Since \(f_N\) vanishes outside \(\mathcal S_N\) and the denominators are positive on \(\mathcal S_N\),
\[
a_j^{(\varsigma)}=f_{N,j}\rho_j^{\varsigma}
\qquad(j\in\mathbb Z).
\]
If \(|j|\leq2N\), then \(N\geq2\) ensures
\[
\omega_j^{\varsigma}\geq N,\qquad
\omega_{j+1}^{\varsigma}\geq N,\qquad
\omega_{j+2}^{\varsigma}\geq N,
\qquad
\omega_{j+1}^{\varsigma}-\omega_j^{\varsigma}=\varsigma.
\]
Factoring the first difference gives
\[
\begin{aligned}
|(\delta\rho^{\varsigma})_j|
&=
\frac1{\omega_j^{\varsigma}(\omega_{j+1}^{\varsigma})^2}
+\frac1{(\omega_j^{\varsigma})^2\omega_{j+1}^{\varsigma}}\\
&\leq\frac2{N^3}\leq\frac{16}{N^3}.
\end{aligned}
\]
For the second difference, the unit spacing of the three denominators gives the exact identity
\[
(\delta^2\rho^{\varsigma})_j
=
\frac{2\{3(\omega_{j+1}^{\varsigma})^2-1\}}
 {(\omega_j^{\varsigma})^2
  (\omega_{j+1}^{\varsigma})^2
  (\omega_{j+2}^{\varsigma})^2}.
\]
Its numerator is nonnegative, and hence
\[
|(\delta^2\rho^{\varsigma})_j|
\leq
\frac6{(\omega_j^{\varsigma})^2(\omega_{j+2}^{\varsigma})^2}
\leq\frac6{N^4}\leq\frac{64}{N^4},
\qquad
|\rho_{j+1}^{\varsigma}|\leq\frac1{N^2}.
\]

The forward-difference product identity is
\[
\begin{aligned}
(\delta^2a^{(\varsigma)})_j
={}&f_{N,j+2}(\delta^2\rho^{\varsigma})_j\\
&+\{(\delta f_N)_{j+1}+(\delta f_N)_j\}
                  (\delta\rho^{\varsigma})_j
+\rho_{j+1}^{\varsigma}(\delta^2f_N)_j.
\end{aligned}
\]
For \(|j|\leq2N\), the preceding estimates imply
\[
\begin{aligned}
|(\delta^2a^{(\varsigma)})_j|
\leq{}&
\frac{64}{N^4}|f_{N,j+2}|
+\frac{16}{N^3}|(\delta f_N)_{j+1}|\\
&+\frac{16}{N^3}|(\delta f_N)_j|
+\frac1{N^2}|(\delta^2f_N)_j|.
\end{aligned}
\]
For \(|j|>2N\), all three values \(f_{N,j}\), \(f_{N,j+1}\), and
\(f_{N,j+2}\) vanish, so this inequality also holds there. In particular,
\(\delta^2a^{(\varsigma)}\) has finite support. Summing and using invariance of full-line sums under integer shifts gives
\[
\begin{aligned}
\sum_j|(\delta^2a^{(\varsigma)})_j|
&\leq
\frac{64}{N^4}\sum_j|f_{N,j}|
+\frac{32}{N^3}\sum_j|(\delta f_N)_j|
+\frac1{N^2}\sum_j|(\delta^2f_N)_j|\\
&\leq\frac{262}{N^3}
\leq\frac{512}{N^3}.
\end{aligned}
\]
% lean: Causalean.Mathlib.Analysis.Approximation.Chebyshev.Jackson.weightedCoeff_delta2_l1

\item We next derive the pointwise localization bound. The reciprocal-square bound on the coefficient support gives, for either sign,
\[
\sum_{j\in\mathbb Z}|a_j^{(\varsigma)}|
\leq\frac1{N^2}\sum_j|f_{N,j}|
\leq\frac2N.
\]
For \(\varsigma\in\{-1,1\}\) and \(u\in\mathbb R\), define the finite Fourier sums
\[
\mathcal F_N^{\varsigma}(u)
=\sum_{j\in\mathbb Z}a_j^{(\varsigma)}e^{iju}.
\]
The definition of \(G_N\) gives
\[
G_N(u)
=-\frac12\operatorname{Re}
\left[
e^{iLu}\mathcal F_N^1(u)
+e^{iLu}\mathcal F_N^{-1}(-u)
\right].
\]
Since the modulation has modulus one,
\[
|G_N(u)|
\leq
\frac{|\mathcal F_N^1(u)|+|\mathcal F_N^{-1}(-u)|}{2}
\leq\frac2N.
\]

For real \(u\), introduce
\[
\eta(u)=\|u\|_{\mathbb T},\qquad
\jmath(u)=\left\lfloor\frac{u}{2\pi}+\frac12\right\rfloor\in\mathbb Z,
\qquad
\upsilon(u)=u-2\pi\jmath(u).
\]
The defining property of the nearest integer gives
\[
|\upsilon(u)|\leq\pi,
\qquad
0\leq\eta(u)\leq|\upsilon(u)|\leq\pi.
\]
Also, replacing the integer in the defining infimum by its negative shows
\[
\eta(-u)=\eta(u).
\]

If \(N\eta(u)\leq1\), then
\[
(1+N\eta(u))^2\leq4,
\qquad
|G_N(u)|\leq\frac2N
\leq\frac{8192}{N(1+N\eta(u))^2}.
\]
This case includes \(\eta(u)=0\).

Suppose instead that \(N\eta(u)>1\), so \(\eta(u)>0\). For every integer
\(j_0\), shifting the finite Fourier sum gives
\[
\sum_j a_{j+j_0}^{(\varsigma)}e^{iju}
=e^{-ij_0u}\mathcal F_N^{\varsigma}(u).
\]
Applying this with shifts one and two yields
\[
\sum_j(\delta^2a^{(\varsigma)})_je^{iju}
=(e^{-iu}-1)^2\mathcal F_N^{\varsigma}(u).
\]
Therefore,
\[
|e^{-iu}-1|^2\,|\mathcal F_N^{\varsigma}(u)|
\leq\sum_j|(\delta^2a^{(\varsigma)})_j|
\leq\frac{512}{N^3}.
\]
Periodicity and the elementary sine bound on \([-\pi/2,\pi/2]\) give
\[
|e^{-iu}-1|
=2|\sin(\upsilon(u)/2)|
\geq\frac2\pi|\upsilon(u)|
\geq\frac2\pi\eta(u)
\geq\frac{\eta(u)}2,
\]
where the last inequality uses \(\pi\leq4\). It follows that
\[
|\mathcal F_N^{\varsigma}(u)|
\leq\frac{2048}{N^3\eta(u)^2}.
\]
Applying this estimate at \(u\) and \(-u\) in the representation of \(G_N\) gives
\[
|G_N(u)|\leq\frac{2048}{N^3\eta(u)^2}.
\]
Finally, \(N\eta(u)>1\) implies
\[
(1+N\eta(u))^2\leq4(N\eta(u))^2,
\]
and hence
\[
|G_N(u)|
\leq\frac{2048}{N^3\eta(u)^2}
\leq\frac{8192}{N(1+N\eta(u))^2}.
\]
Together, the two cases establish
\[
|G_N(u)|
\leq\frac{8192}{N(1+N\|u\|_{\mathbb T})^2}
\qquad(u\in\mathbb R).
\]
% lean: Causalean.Mathlib.Analysis.Approximation.Chebyshev.Jackson.packet_pointwise_localization

\item We obtain the normalized integral bound by integrating a rational envelope. First, for \(|u|\leq\pi\), the integer zero in the defining infimum gives
\(\eta(u)\leq|u|\). For every \(j\leq-1\),
\[
|u-2\pi j|\geq u+2\pi\geq\pi\geq|u|,
\]
and for every \(j\geq1\),
\[
|u-2\pi j|\geq2\pi-u\geq\pi\geq|u|.
\]
The remaining case \(j=0\) gives equality. Thus
\[
\eta(u)=|u|\qquad(-\pi\leq u\leq\pi).
\]
Since
\[
1+(Nu)^2\leq(1+N|u|)^2,
\]
the pointwise estimate implies
\[
\frac{|G_N(u)|}{2\pi}
\leq
\frac{8192}{2\pi N}\,
\frac1{1+(Nu)^2}
\qquad(-\pi\leq u\leq\pi).
\]
The rational envelope has integral
\[
\begin{aligned}
\int_{-\pi}^{\pi}\frac{du}{1+(Nu)^2}
&=\frac{\arctan(N\pi)-\arctan(-N\pi)}N\\
&\leq\frac\pi N,
\end{aligned}
\]
because the arctangent takes values between \(-\pi/2\) and \(\pi/2\).
Both integrands in the comparison are continuous on the compact interval.
Consequently,
\[
\int_{-\pi}^{\pi}|G_N(u)|\,\frac{du}{2\pi}
\leq
\frac{8192}{2\pi N}\,\frac\pi N
=\frac{4096}{N^2}.
\]
% lean: Causalean.Mathlib.Analysis.Approximation.Chebyshev.Jackson.packet_l1_localization

\item All these estimates hold uniformly for \(N\geq2\). Choose
\[
c=\frac1{72},\qquad
C_{\mathrm{center}}=\frac12,\qquad
C_{\mathrm{point}}=8192,\qquad
C_{\mathrm{integral}}=4096,
\]
and set
\[
C=\max\left\{
512,\,
\max\left\{
C_{\mathrm{center}},
\max\{C_{\mathrm{point}},C_{\mathrm{integral}}\}
\right\}
\right\}
=8192.
\]
Then
\[
0<c\leq C,\qquad
512\leq C,\qquad
C_{\mathrm{center}}\leq C,\qquad
C_{\mathrm{point}}\leq C,\qquad
C_{\mathrm{integral}}\leq C.
\]
We now fix any integer \(N\geq2\) and verify the assertions with these common constants.
% lean: CausalSmith.Stat.OptvalueVanishingoverlapRate.jackson_packet_localization

\item The positive normalizing denominator computed above gives
\[
\int_{-\pi}^{\pi}\mathcal J_N(u)\,\frac{du}{2\pi}
=c_N\int_{-\pi}^{\pi}\mathcal R_N(u)\,\frac{du}{2\pi}
=\gamma_{N,0}^{-1}\gamma_{N,0}
=1.
\]
% lean: CausalSmith.Stat.OptvalueVanishingoverlapRate.packetKernel_normalized_integral

\item For either sign, the zero extension implies
\[
(\delta^2a^\pm)_j=0\qquad(|j|>2N).
\]
Hence the absolute second differences are summable, and their full-line sums satisfy
\[
\sum_{j\in\mathbb Z}|(\delta^2a^\pm)_j|
=
\sum_{j=-2N}^{2N}|(\delta^2a^\pm)_j|
\leq\frac{512}{N^3}
\leq\frac C{N^3}.
\]
Because the differences are taken after extending the coefficients by zero, these sums include the differences crossing both endpoints of \(\mathcal S_N\).
% lean: CausalSmith.Stat.OptvalueVanishingoverlapRate.packetCoeffPlus_secondDiff_summable
% lean: CausalSmith.Stat.OptvalueVanishingoverlapRate.packetCoeffMinus_secondDiff_summable
% lean: Causalean.Mathlib.Analysis.Approximation.Chebyshev.Jackson.weightedCoeff_delta2_l1

\item Multiplying the cosine expansion of \(\mathcal J_N\) by \(\cos(Lu)\) gives
\[
k_N(u)
=\frac12\sum_{j\in\mathcal S_N}f_{N,j}
\left\{\cos((L+j)u)+\cos((L-j)u)\right\}.
\]
Every frequency appearing here satisfies
\[
2N+2\leq L\pm j\leq6N-2.
\]
Let \(r\in\mathbb Z\) satisfy the stated alternative. If \(|r|<2N+2\), both shifted frequencies exceed \(|r|\); if \(|r|>6N-2\), both are smaller than \(|r|\). Thus
\[
|L+j|\neq|r|,\qquad |L-j|\neq|r|
\qquad(j\in\mathcal S_N).
\]
Cosine orthogonality consequently gives
\[
\int_{-\pi}^{\pi}\cos((L\pm j)u)\cos(ru)\,du=0.
\]
For the sine integrals, the product identity gives
\[
\begin{aligned}
&\int_{-\pi}^{\pi}\cos((L\pm j)u)\sin(ru)\,du\\
&\qquad=
\frac12\int_{-\pi}^{\pi}
\left\{\sin((r+L\pm j)u)+\sin((r-(L\pm j))u)\right\}\,du
=0.
\end{aligned}
\]
Here every integer sine mode has integral zero; the zero mode is identically zero, and the other modes have equal cosine values at the two endpoints in their antiderivatives. Linearity in the finite expansion therefore yields
\[
\int_{-\pi}^{\pi}k_N(u)\cos(ru)\,du=0,
\qquad
\int_{-\pi}^{\pi}k_N(u)\sin(ru)\,du=0.
\]
% lean: Causalean.Mathlib.Analysis.Approximation.Chebyshev.Jackson.packet_shifted_support

\item Termwise differentiation of the finite cosine polynomial gives, for every real \(u\),
\[
G_N'(u)
=\frac12\sum_{j\in\mathcal S_N}
\left\{
a_j^+(L+j)\sin((L+j)u)
+a_j^-(L-j)\sin((L-j)u)
\right\},
\]
and then
\[
G_N''(u)
=\frac12\sum_{j\in\mathcal S_N}
\left\{
a_j^+(L+j)^2\cos((L+j)u)
+a_j^-(L-j)^2\cos((L-j)u)
\right\}.
\]
The shifted frequencies are positive on \(\mathcal S_N\), so their squares cancel the denominators:
\[
a_j^+(L+j)^2=f_{N,j},
\qquad
a_j^-(L-j)^2=f_{N,j}.
\]
Using
\[
\cos((L+j)u)+\cos((L-j)u)
=2\cos(ju)\cos(Lu),
\]
we conclude that
\[
G_N''(u)
=
\left(\sum_{j\in\mathcal S_N}f_{N,j}\cos(ju)\right)\cos(Lu)
=\mathcal J_N(u)\cos(Lu)
=k_N(u).
\]
% lean: Causalean.Mathlib.Analysis.Approximation.Chebyshev.Jackson.packet_twice_deriv

\item Every frequency \(L\pm j\), for \(j\in\mathcal S_N\), is a nonzero integer. Its cosine integral vanishes by direct integration. Thus
\[
\begin{aligned}
\int_{-\pi}^{\pi}G_N(u)\,du
=-\frac12\sum_{j\in\mathcal S_N}
\left\{
a_j^+\int_{-\pi}^{\pi}\cos((L+j)u)\,du
+a_j^-\int_{-\pi}^{\pi}\cos((L-j)u)\,du
\right\}
=0.
\end{aligned}
\]
% lean: Causalean.Mathlib.Analysis.Approximation.Chebyshev.Jackson.packet_zero_mean

\item Finally, enlarging the three upper constants to the common constant \(C\) gives
\[
\frac cN\leq|G_N(0)|
\leq\frac{C_{\mathrm{center}}}N
\leq\frac CN,
\]
\[
|G_N(u)|
\leq
\frac{C_{\mathrm{point}}}
 {N(1+N\|u\|_{\mathbb T})^2}
\leq
\frac C{N(1+N\|u\|_{\mathbb T})^2}
\qquad(u\in\mathbb R),
\]
and
\[
\int_{-\pi}^{\pi}|G_N(u)|\,\frac{du}{2\pi}
\leq\frac{C_{\mathrm{integral}}}{N^2}
\leq\frac C{N^2}.
\]
These bounds, together with the preceding identities and summability statements, establish every assertion.
% lean: CausalSmith.Stat.OptvalueVanishingoverlapRate.jackson_packet_localization
\end{enumerate}
\end{proof}

Frequency cancellation and spatial localization play complementary roles in \cref{obj:lem:jackson-packet-localization}. The frequency band supplies the cancellation required for matching low-order moments after the packet is transferred to the intensity interval. The twice antiderivative has an origin value of order \(N^{-1}\) and an integrated absolute magnitude of order \(N^{-2}\). These two scales distinguish the localized response at the scalar target's change in slope from the contribution distributed away from that point.

The second-difference bound includes the Fourier endpoints explicitly. Together with the spatial bound, it controls the entire periodic twice antiderivative, including the tails measured by distance from the periodic origin. The normalized packet used below combines this localization with a calibration that leaves sufficient mass and first moment for the common-atom completion.

\subsection{Weighted separation and statistical consequences}

The packet localization lemma supplies low-frequency cancellation and spatial concentration, and the weighted-separation theorem converts the packet into common-mean moment priors and an approximation lower bound. \Cref{sec:appendix-lower-bound} uses those priors in the statistical comparison. The detailed approximation, sparse, dense, transfer, and concluding arguments are divided explicitly in \cref{sec:deferred-proofs}.

The next result collects the approximation comparison and the completed-prior construction. It also records their statistical consequences alongside the attaining upper bound. Within the scalar parameter range, weighted approximation error and common-mean moment-prior separation have the same order, uniformly over the overlap parameter.

\begin{theoremv}{thm:weighted-separation-and-frontier}[Weighted separation and minimax frontier]
There exist universal real constants \(0<c<C<\infty\) and an integer \(A_0\geq1\) such that the following approximation and separation guarantees hold for every choice of parameters satisfying:
\begin{itemize}
\item \textbf{(Degree.)} \(K\) is an integer with \(K\geq2\).
\item \textbf{(Support.)} \(M\) is real and \(2\leq M\leq K^2\).
\item \textbf{(Overlap.)} \(0\leq\epsilon\leq1/2\).
\end{itemize}
For the weighted error and constrained-prior separation in
\cref{obj:def:weighted-error,obj:def:constrained-prior-separation}, with scalar functional
\[
\phi_\epsilon(z)=
\frac{((1-\epsilon)+\epsilon z)\max\{z-1,0\}}{1+z},
\]
one has
\[
c\frac{\sqrt M}{K}
\leq \mathcal E_{K,M,\epsilon}^{w}
\leq \Delta_{K,M,\epsilon}
\leq 4\mathcal E_{K,M,\epsilon}^{w}
\leq C\frac{\sqrt M}{K}.
\]
Moreover, let \(\sigma_+\) and \(\sigma_-\) be the positive and negative Jordan parts of the normalized Jackson packet with calibration \(A_0\), pushed to \([0,M]\). These are finite measures with integrable first moments. Their positive mass and common completion location satisfy
\[
t=\sigma_+(\mathbb R)\leq\frac12,
\qquad
u=\int z\,d\sigma_+(z),
\qquad
z_0=\frac{1-u}{1-t},
\qquad
\frac12\leq z_0\leq2.
\]
There exists \((\nu_0,\nu_1)\in\mathfrak N_{K,M}\), the constrained prior class of \cref{obj:def:constrained-prior-separation}, such that the common-atom completion of \cref{obj:def:weighted-handle} satisfies
\[
\mathscr H^w(\sigma_+,\sigma_-)
=(\nu_0,\nu_1)
=\bigl(\sigma_-+(1-t)\delta_{z_0},
       \sigma_++(1-t)\delta_{z_0}\bigr),
\]
and
\[
c\frac{\sqrt M}{K}
\leq
\left|
\int\phi_\epsilon(z)\,d\nu_1(z)
-
\int\phi_\epsilon(z)\,d\nu_0(z)
\right|.
\]

There also exist universal real constants
\(0<c_\star\leq C_\star<\infty\), \(H_0>0\), and \(\kappa>0\), and a universal integer \(D_0\geq2\), such that the armwise estimator
\(\widehat V_{n,d,\epsilon}^{\mathrm{AF}}\) with these calibrations satisfies the following risk guarantees for every choice of parameters satisfying:
\begin{itemize}
\item \textbf{(Sample size.)} \(n\) is an integer with \(n\geq1\).
\item \textbf{(Alphabet.)} \(d\) is an integer with \(d\geq2\).
\item \textbf{(Overlap.)} \(0<\epsilon\leq1/2\).
\end{itemize}
The observed risks obey
\[
c_\star\min\left\{1,\frac{d}{n\epsilon\log(ed)}\right\}
\leq \mathfrak R_{n,d,\epsilon}
\leq
\sup_{\mathbb P\in\mathcal D_{d,\epsilon}^{\mathrm{obs}}}
E_{\mathbb P}\!\left[
\left\{\widehat V_{n,d,\epsilon}^{\mathrm{AF}}
-\Psi(\mathbb P)\right\}^2
\right]
\leq
C_\star\min\left\{1,\frac{d}{n\epsilon\log(ed)}\right\}.
\]
The causal minimax risk obeys
\[
c_\star\min\left\{1,\frac{d}{n\epsilon\log(ed)}\right\}
\leq
\inf_{\widehat V}
\sup_{P\in\mathcal P_{d,\epsilon}}
E_{\operatorname{Obs}(P)}\!\left[
\left\{\widehat V-V^\star(P)\right\}^2
\right]
\leq
C_\star\min\left\{1,\frac{d}{n\epsilon\log(ed)}\right\},
\]
where the infimum ranges over measurable real-valued estimators based on \(n\) observed units.

Finally, consider any sequences satisfying:
\begin{itemize}
\item \textbf{(Sample sizes.)} The integer sample sizes \(n_k\) tend to infinity.
\item \textbf{(Alphabets.)} Each \(d_k\) is an integer with \(d_k\geq2\).
\item \textbf{(Overlap.)} \(0<\epsilon_k\leq1/2\) for every \(k\).
\end{itemize}
There exists a sequence of measurable real-valued estimators
\(\widehat V_k\), each based on \(n_k\) observed units with alphabet size \(d_k\), such that
\[
\sup_{\mathbb P\in\mathcal D_{d_k,\epsilon_k}^{\mathrm{obs}}}
E_{\mathbb P}\!\left[
\left\{\widehat V_k-\Psi(\mathbb P)\right\}^2
\right]\longrightarrow0
\]
if and only if
\[
\frac{d_k}{n_k\epsilon_k\log(ed_k)}\longrightarrow0.
\]
\end{theoremv}

The approximation comparison expresses the same difficulty in two forms. Polynomial approximation measures how accurately the scalar target can be represented through low-order moments; constrained separation measures how far target values can differ when those moments agree. The packet construction realizes a separation of order \(\sqrt M/K\) using probability priors with common mean one. Adding the shared completion atom preserves the target contrast, while the bounds on its location place the completed priors within the prescribed support.

The factor four in the upper comparison reflects the normalization in \cref{obj:def:constrained-prior-separation}: each prior assigns integral two to the approximation weight \(1+z\). The remaining inequalities in \cref{obj:thm:weighted-separation-and-frontier} establish both the reverse comparison and the uniform approximation scale. In particular, the explicit completed pair achieves the stated lower order under the probability and mean constraints.

The statistical clauses connect this scalar construction to the main-body results. The constrained priors supply the separation used by the all-estimator lower bound in \cref{obj:thm:all-estimator-lower}, while \cref{obj:thm:observable-upper} supplies the attaining risk bound for the estimator of \cref{obj:def:armwise-estimator}. Their combination gives the matched observed and causal rates stated in \cref{obj:thm:matched-frontier}, with the causal interpretation supported by \cref{obj:prop:identification}. Along the admissible sequences, the same two-sided risk comparison yields the necessary and sufficient condition for uniform consistency in \cref{obj:thm:consistency-threshold}.


\section{Statistical lower bound and verification scope}\label{sec:appendix-lower-bound}

This appendix specifies the experiments used to pass from the constrained priors to \cref{obj:thm:all-estimator-lower}. The detailed sparse and dense comparisons, fixed-sample transfer, and final consequences appear in \cref{sec:deferred-proofs}.

The constrained priors of \cref{obj:def:constrained-prior-separation} connect the scalar separation in \cref{obj:thm:weighted-separation-and-frontier} to a comparison between observed-data experiments. The statistical construction has three components: normalization into admissible observed laws, comparison of their full armwise likelihoods, and transfer from an auxiliary count experiment to the fixed-sample risk in \cref{obj:def:minimax-risk}. We first specify the family on which these components operate.

\begin{definitionv}{def:sparse-family}[Sparse observed law]
The sparse observed law \(\mathbb P_{\mathbf z}^{\mathrm{sp}}\) is defined for a covariate alphabet size \(d\in\mathbb N\), an overlap parameter \(\epsilon\in\mathbb R\), an intensity support endpoint \(M\in\mathbb R\), and a vector of cell intensities \(\mathbf z\), with
\begin{itemize}
\item \textbf{(Alphabet.)} \(d>0\).
\item \textbf{(Overlap.)} \(0<\epsilon\le 1/2\).
\item \textbf{(Intensities.)} \(\mathbf z\in[0,M]^d\).
\end{itemize}
For each covariate cell \(x\in[d]\), define the unnormalized atom masses by
\[
\widetilde q_{11,x}=\frac{\epsilon(1-\epsilon)z_x}{d},
\qquad
\widetilde q_{10,x}=\frac{\epsilon(1-\epsilon)}{d},
\qquad
\widetilde q_{00,x}=\widetilde q_{01,x}
=\frac{(1-\epsilon)^2+\epsilon^2z_x}{2d}.
\]
Their total mass is
\[
S(\mathbf z)
=\sum_{x\in[d]}\sum_{a=0}^1\sum_{y=0}^1\widetilde q_{ay,x}.
\]
The normalized probability law of the categorical covariate \(X\), binary treatment \(A\), and binary observed outcome \(Y\) is
\[
\mathbb P_{\mathbf z}^{\mathrm{sp}}(X=x,A=a,Y=y)
=\frac{\widetilde q_{ay,x}}{S(\mathbf z)},
\qquad x\in[d],\quad a,y\in\{0,1\}.
\]
\end{definitionv}

The normalizing mass \leanref{sym:S(\mathbf z)}{\ensuremath{S(\mathbf z)}} in \cref{obj:def:sparse-family} equals
\[
S(\mathbf z)=1-\epsilon+\frac{\epsilon}{d}\sum_{x\in[d]}z_x.
\]
Consequently, the common-mean-one restriction on the constrained priors centers this mass at one when intensities are drawn independently across cells. The deterministic bound \(S(\mathbf z)\ge 1-\epsilon\ge 1/2\) also holds throughout the parameter range used below. These two properties give normalization a distinct role: it produces a probability law for every intensity vector while preserving a controlled connection between the scalar contrast and the observed value.

For the auxiliary experiment, fix \leanref{sym:B}{\ensuremath{B}}, the Poisson inflation parameter, with \(B>2\). Conditional on an intensity vector, observe independent counts in every cell and arm-outcome category, with
\[
N_{ay,x}\sim\operatorname{Poisson}
\bigl(Bn\widetilde q_{ay,x}(z_x)\bigr).
\]
The total count, denoted \(N_{\mathrm{tot}}=\sum_{x,a,y}N_{ay,x}\), has mean \(BnS(\mathbf z)\). Conditional on that total, the atom counts have the multinomial law associated with \cref{obj:def:sparse-family}. Thus the auxiliary experiment retains both treatment arms and all outcome categories. Its one-cell comparison uses the reference intensity \leanref{sym:z_\star}{\ensuremath{z_\star}}, equal to \(M/2\), and the likelihood ratio \leanref{sym:L_z}{\ensuremath{L_z}}, defined explicitly in the following statement.

\begin{lemmav}{lem:sparse-likelihood}[Sparse law and likelihood identities]
Let the parameters satisfy:
\begin{itemize}
\item \textbf{(Sample and alphabet sizes.)} \(n,d\in\mathbb N\), \(n\ge1\), and \(d\ge2\), with the zero-based cell alphabet \([d]\) of \cref{obj:synth_1}.
\item \textbf{(Overlap.)} \(0<\epsilon\le1/2\).
\item \textbf{(Intensity support.)} \(M\ge2\) and \(\mathbf z\in[0,M]^d\).
\item \textbf{(Poisson inflation.)} \(B>2\).
\end{itemize}
The normalized sparse law \(\mathbb P_{\mathbf z}^{\mathrm{sp}}\) of \cref{obj:def:sparse-family} belongs to \(\mathcal D_{d,\epsilon}^{\mathrm{obs}}\) of \cref{obj:def:observed-class}. For every cell \(x\in[d]\), its cell probability, treatment propensity, and conditional outcome means satisfy
\[
p_x=\frac{1-\epsilon+\epsilon z_x}{dS(\mathbf z)},
\qquad
\pi_x=\frac{\epsilon(1-\epsilon)(1+z_x)}
{1-\epsilon+\epsilon z_x},
\qquad
\mu_{0x}=\frac12,
\qquad
\mu_{1x}=\frac{z_x}{1+z_x}.
\]
Define the scalar functional
\[
\phi_\epsilon(z)
=\frac{(1-\epsilon+\epsilon z)\max\{z-1,0\}}{1+z}.
\]
The observed value of \cref{obj:def:observed-value} is
\[
\Psi(\mathbb P_{\mathbf z}^{\mathrm{sp}})
=\frac12+\frac{1}{2dS(\mathbf z)}
\sum_{x\in[d]}\phi_\epsilon(z_x).
\]

For the one-cell Poisson experiment, write
\(\widetilde q_{ay,x}(z)\) for the unnormalized atom mass in
\cref{obj:def:sparse-family} with cell intensity \(z\), and set the reference intensity \(z_\star=M/2\). Under intensity \(z\), the four counts \(N_{ay,x}\), \(a,y\in\{0,1\}\), are independent Poisson variables with means \(Bn\widetilde q_{ay,x}(z)\). Define their likelihood ratio relative to intensity \(z_\star\) by
\[
L_z(N)
=
\prod_{a,y\in\{0,1\}}
\exp\!\left\{
Bn\bigl(\widetilde q_{ay,x}(z_\star)
-\widetilde q_{ay,x}(z)\bigr)
\right\}
\left(
\frac{\widetilde q_{ay,x}(z)}
{\widetilde q_{ay,x}(z_\star)}
\right)^{N_{ay,x}}.
\]
Writing \(E_{z_\star}\) for expectation under these independent Poisson counts at reference intensity \(z_\star\), for every \(z,w\in[0,M]\),
\[
E_{z_\star}[L_zL_w]
=\exp\!\left\{\alpha(z-z_\star)(w-z_\star)\right\},
\]
where the Gram coefficient is
\[
\alpha
=\frac{Bn}{d}
\left\{
\frac{\epsilon(1-\epsilon)}{z_\star}
+\frac{\epsilon^4}{(1-\epsilon)^2+\epsilon^2z_\star}
\right\}
\le \frac{2Bn\epsilon}{dM}.
\]
\end{lemmav}

\begin{proof}[Proof of \cref{obj:lem:sparse-likelihood}]
\begin{enumerate}
\item By \cref{obj:def:sparse-family}, the normalized atom probabilities are
\[
q_{ay,x}=\frac{\widetilde q_{ay,x}}{S(\mathbf z)}.
\]
Summing the four unnormalized masses in each cell gives
\[
\sum_{a,y\in\{0,1\}}\widetilde q_{ay,x}
=\frac{\epsilon(1-\epsilon)(1+z_x)+(1-\epsilon)^2+\epsilon^2z_x}{d}
=\frac{1-\epsilon+\epsilon z_x}{d}.
\]
Consequently,
\[
S(\mathbf z)
=1-\epsilon+\frac{\epsilon}{d}\sum_{x\in[d]}z_x
\ge 1-\epsilon\ge\frac12>0.
\]
Since \(d>0\) and \(z_x\ge0\), we also have
\[
1-\epsilon+\epsilon z_x>0.
\]
Summing the normalized atom probabilities therefore yields
\[
p_x=\frac{1-\epsilon+\epsilon z_x}{dS(\mathbf z)},
\qquad
\sum_{y\in\{0,1\}}q_{1y,x}
=\frac{\epsilon(1-\epsilon)(1+z_x)}{dS(\mathbf z)}.
\]
Taking their ratio, as in \cref{obj:def:observed-class}, gives
\[
\pi_x
=\frac{\epsilon(1-\epsilon)(1+z_x)}
{1-\epsilon+\epsilon z_x}.
\]
The overlap inequalities follow by clearing the positive denominator:
\[
\pi_x-\epsilon
=\frac{\epsilon(1-2\epsilon)z_x}{1-\epsilon+\epsilon z_x}\ge0,
\qquad
1-\epsilon-\pi_x
=\frac{(1-\epsilon)(1-2\epsilon)}
{1-\epsilon+\epsilon z_x}\ge0.
\]
Here \(1-2\epsilon\ge0\), \(z_x\ge0\), and \(1-\epsilon>0\).
Thus the sparse law satisfies \cref{obj:ass:shrinking-overlap} and belongs to the observed class of \cref{obj:def:observed-class}.
% lean: CausalSmith.Stat.OptvalueVanishingoverlapRate.sparse_likelihood

\item The preceding calculations establish the asserted cell probability and propensity. For the conditional outcome means, the arm masses and outcome-one atom probabilities are
\[
\begin{aligned}
\sum_{y\in\{0,1\}}q_{0y,x}
&=\frac{(1-\epsilon)^2+\epsilon^2z_x}{dS(\mathbf z)},
&
q_{01,x}
&=\frac{(1-\epsilon)^2+\epsilon^2z_x}{2dS(\mathbf z)},\\
\sum_{y\in\{0,1\}}q_{1y,x}
&=\frac{\epsilon(1-\epsilon)(1+z_x)}{dS(\mathbf z)},
&
q_{11,x}
&=\frac{\epsilon(1-\epsilon)z_x}{dS(\mathbf z)}.
\end{aligned}
\]
Both arm masses are strictly positive, because
\[
(1-\epsilon)^2+\epsilon^2z_x>0,
\qquad
\epsilon(1-\epsilon)(1+z_x)>0.
\]
Dividing the outcome-one probabilities by their respective arm masses, using \cref{obj:def:observed-class}, gives
\[
\mu_{0x}=\frac12,
\qquad
\mu_{1x}=\frac{z_x}{1+z_x}.
\]
These divisions remain valid at \(z_x=0\).
% lean: CausalSmith.Stat.OptvalueVanishingoverlapRate.sparse_likelihood

\item To evaluate the observed value, split each cell according to whether \(z_x\le1\) or \(z_x\ge1\). Since \(1+z_x>0\),
\[
z_x\le1\ \Longrightarrow\ \frac{z_x}{1+z_x}\le\frac12,
\qquad
z_x\ge1\ \Longrightarrow\ \frac{z_x}{1+z_x}\ge\frac12.
\]
In the first case \(\max\{z_x-1,0\}=0\), and in the second it equals \(z_x-1\). Hence, in both cases,
\[
\max\left\{\frac12,\frac{z_x}{1+z_x}\right\}
=\frac12+\frac{\max\{z_x-1,0\}}{2(1+z_x)}.
\]
Using the definition of \(\phi_\epsilon\) in the statement, the cell contribution becomes
\[
p_x\max\{\mu_{0x},\mu_{1x}\}
=\frac{p_x}{2}
+\frac{\phi_\epsilon(z_x)}{2dS(\mathbf z)}.
\]
Moreover, the normalizer identity from the first step gives
\[
\sum_{x\in[d]}p_x
=\frac{\sum_{x\in[d]}(1-\epsilon+\epsilon z_x)}
{dS(\mathbf z)}
=\frac{dS(\mathbf z)}{dS(\mathbf z)}
=1.
\]
Summing the cell contributions according to \cref{obj:def:observed-value} therefore yields
\[
\Psi(\mathbb P_{\mathbf z}^{\mathrm{sp}})
=\frac12+\frac{1}{2dS(\mathbf z)}
\sum_{x\in[d]}\phi_\epsilon(z_x).
\]
% lean: CausalSmith.Stat.OptvalueVanishingoverlapRate.sparse_likelihood

\item Fix a cell \(x\). For \(z\in[0,M]\) and \(a,y\in\{0,1\}\), define the Poisson means and their reference values by
\[
\lambda_{ay,x}(z):=Bn\widetilde q_{ay,x}(z),
\qquad
\lambda_{\star,ay,x}:=\lambda_{ay,x}(z_\star).
\]
Explicitly,
\[
\begin{aligned}
\lambda_{11,x}(z)&=\frac{Bn}{d}\epsilon(1-\epsilon)z,
&
\lambda_{10,x}(z)&=\frac{Bn}{d}\epsilon(1-\epsilon),\\
\lambda_{00,x}(z)&=\lambda_{01,x}(z)
=\frac{Bn}{2d}\bigl((1-\epsilon)^2+\epsilon^2z\bigr).
\end{aligned}
\]
These means are nonnegative throughout \([0,M]\). Since \(B>2\), \(n\ge1\), \(d\ge2\), and \(z_\star=M/2\ge1\), all four reference means are strictly positive. Powers at zero are interpreted with \(0^0=1\).

For each coordinate and every \(\tau\in\mathbb R\), the Poisson exponential series gives
\[
\begin{aligned}
E_{z_\star}\!\left[\tau^{N_{ay,x}}\right]
&=\exp(-\lambda_{\star,ay,x})
\sum_{\rho=0}^{\infty}
\frac{(\lambda_{\star,ay,x}\tau)^\rho}{\rho!}\\
&=\exp\!\left\{\lambda_{\star,ay,x}(\tau-1)\right\}.
\end{aligned}
\]
Here \(\rho\) ranges over the nonnegative integers; the series converges absolutely, as is seen by replacing \(\tau\) with \(|\tau|\). Combining the two likelihood factors for this coordinate and applying the identity with
\(\tau=\lambda_{ay,x}(z)\lambda_{ay,x}(w)/\lambda_{\star,ay,x}^2\) gives
\[
\begin{aligned}
&E_{z_\star}\!\left[
\left\{\exp(\lambda_{\star,ay,x}-\lambda_{ay,x}(z))
\left(\frac{\lambda_{ay,x}(z)}{\lambda_{\star,ay,x}}\right)^{N_{ay,x}}\right\}
\left\{\exp(\lambda_{\star,ay,x}-\lambda_{ay,x}(w))
\left(\frac{\lambda_{ay,x}(w)}{\lambda_{\star,ay,x}}\right)^{N_{ay,x}}\right\}
\right]\\
&\quad=
\exp\!\left\{
2\lambda_{\star,ay,x}-\lambda_{ay,x}(z)-\lambda_{ay,x}(w)
+\lambda_{\star,ay,x}
\left(\frac{\lambda_{ay,x}(z)\lambda_{ay,x}(w)}
{\lambda_{\star,ay,x}^2}-1\right)
\right\}\\
&\quad=
\exp\!\left\{
\frac{(\lambda_{ay,x}(z)-\lambda_{\star,ay,x})
(\lambda_{ay,x}(w)-\lambda_{\star,ay,x})}
{\lambda_{\star,ay,x}}
\right\}.
\end{aligned}
\]
% lean: CausalSmith.Stat.OptvalueVanishingoverlapRate.poisson_likelihood_gram_scalar

Independence of the four reference counts now implies
\[
E_{z_\star}[L_zL_w]
=\exp\!\left\{
\sum_{a,y\in\{0,1\}}
\frac{(\lambda_{ay,x}(z)-\lambda_{\star,ay,x})
(\lambda_{ay,x}(w)-\lambda_{\star,ay,x})}
{\lambda_{\star,ay,x}}
\right\}.
\]
The treated-success contribution is
\[
\frac{(\lambda_{11,x}(z)-\lambda_{\star,11,x})
(\lambda_{11,x}(w)-\lambda_{\star,11,x})}
{\lambda_{\star,11,x}}
=
\frac{Bn}{d}\frac{\epsilon(1-\epsilon)}{z_\star}
(z-z_\star)(w-z_\star).
\]
The treated-failure contribution equals zero because its mean is constant. The two control contributions sum to
\[
\sum_{y\in\{0,1\}}
\frac{(\lambda_{0y,x}(z)-\lambda_{\star,0y,x})
(\lambda_{0y,x}(w)-\lambda_{\star,0y,x})}
{\lambda_{\star,0y,x}}
=
\frac{Bn}{d}
\frac{\epsilon^4}{(1-\epsilon)^2+\epsilon^2z_\star}
(z-z_\star)(w-z_\star).
\]
Thus their total is \(\alpha(z-z_\star)(w-z_\star)\), with \(\alpha\) exactly as in the statement, proving the likelihood identity.
% lean: CausalSmith.Stat.OptvalueVanishingoverlapRate.sparse_likelihood

\item Finally, \(M>0\) and
\[
(1-\epsilon)^2+\epsilon^2M/2>0.
\]
Clearing these positive denominators, the required inner bound follows from
\[
\begin{aligned}
&2\epsilon^2\bigl((1-\epsilon)^2+\epsilon^2M/2\bigr)
-M\epsilon^4\\
&\qquad=2\epsilon^2(1-\epsilon)^2\ge0.
\end{aligned}
\]
Indeed, this is precisely the numerator obtained by subtracting the left-hand side from the right-hand side of
\[
\frac{\epsilon(1-\epsilon)}{M/2}
+\frac{\epsilon^4}{(1-\epsilon)^2+\epsilon^2M/2}
\le\frac{2\epsilon}{M}
\]
and multiplying by \(M\bigl((1-\epsilon)^2+\epsilon^2M/2\bigr)\).
Multiplication by the nonnegative factor \(Bn/d\), together with \(z_\star=M/2\), gives
\[
\alpha
\le\frac{Bn}{d}\frac{2\epsilon}{M}
=\frac{2Bn\epsilon}{dM}.
\]
% lean: CausalSmith.Stat.OptvalueVanishingoverlapRate.sparseGramCoefficient_bound
\end{enumerate}
\end{proof}

The Gram coefficient \leanref{sym:\alpha}{\ensuremath{\alpha}} in \cref{obj:lem:sparse-likelihood} records the information carried by the complete one-cell experiment. Its first term comes from treated successes, and its second term includes the intensity dependence of both control counts. The bound therefore supports a comparison using the full observed likelihood. The same statement places every normalized alternative within the public overlap class, including intensities at the support endpoints.

The constrained-prior construction in \cref{obj:def:weighted-handle,obj:thm:weighted-separation-and-frontier} supplies probability priors with common mean one and matching moments through the prescribed degree. Drawing cell intensities independently under each prior gives two product mixtures of the auxiliary experiment. Their statistical role is to combine a scalar target contrast with likelihood agreement governed by \cref{obj:lem:sparse-likelihood}. Because the observed value divides the scalar sum by the random normalizing mass, target separation is assessed for that normalized value. The common first moment and the bounded support are the conditions controlling this additional source of variation.

The many-cell comparison uses the constrained priors over their admissible support range. In the dense regime, where rare-arm observations per cell support local comparisons, the boundary-propensity construction places treatment probabilities at the overlap floor and outcome means near the point at which the preferred arm changes. This supplies the corresponding local comparison for the maximum of the two conditional means. For bounded alphabets, a rare-arm two-point comparison supplies the risk scale governed by the total rare-arm information \(n\epsilon\); its bounded-separation branch also covers small \(n\epsilon\). These comparisons serve complementary parameter ranges in the universal bound below.

The fixed-sample transfer concerns bounded squared loss. An estimator may be projected onto \([0,1]\), since the observed value lies in that interval. In the auxiliary experiment, randomly ordering the recorded atoms gives independent observations from the normalized sparse law conditional on the total count. Applying the fixed-sample estimator to the first \(n\) records on \(N_{\mathrm{tot}}\ge n\), and assigning a value in \([0,1]\) otherwise, charges at most one unit of squared loss on the insufficient-count event. Its probability is uniformly bounded by
\[
\Pr\!\left\{\operatorname{Poisson}(Bn/2)<n\right\},
\]
using the deterministic normalization bound. The inflation parameter controls this explicit transfer cost. Accordingly, the auxiliary comparison yields a statement about exactly \(n\) observations and the full estimator class in \cref{obj:def:minimax-risk}.

The resulting lower bound combines the normalized many-cell comparison with the dense and two-point regimes.

\begin{theoremv}{thm:all-estimator-lower}[Universal minimax risk lower bound]
There exists a universal constant \(c_\star>0\) such that, for every sample size \(n\in\mathbb N\), alphabet size \(d\in\mathbb N\), and overlap parameter \(\epsilon\in\mathbb R\) satisfying
\begin{itemize}
\item \textbf{(Sample size.)} \(n\ge1\).
\item \textbf{(Alphabet size.)} \(d\ge2\).
\item \textbf{(Overlap.)} \(0<\epsilon\le1/2\).
\end{itemize}
the minimax squared-error risk \(\mathfrak R_{n,d,\epsilon}\) satisfies
\[
\mathfrak R_{n,d,\epsilon}
\ge c_\star\min\left\{1,\frac{d}{n\epsilon L_d}\right\},
\qquad L_d=\log(ed).
\]
\end{theoremv}

The lower bound expresses the joint effect of dividing observations among covariate cells and observing the informative treatment arm with probability near the overlap floor. Matching low-order moments makes the scalar contrast difficult to distinguish across many cells, while the full likelihood identity accounts for the information available in both arms. The truncation at one describes the bounded-risk scale when the combined information is small. Because the risk in \cref{obj:def:minimax-risk} ranges over all measurable estimators, the comparison applies throughout that estimator class.

Combining \cref{obj:thm:all-estimator-lower,obj:thm:observable-upper} gives the two-sided observed-risk comparison in \cref{obj:thm:matched-frontier}. The identification and completion statement in \cref{obj:prop:identification} supplies its causal interpretation under the stated causal assumptions. Along the admissible sequences of \cref{obj:thm:consistency-threshold}, the two-sided comparison makes vanishing risk equivalent to
\[
\frac{d_n}{n\epsilon_n\log(ed_n)}\longrightarrow0.
\]
Here \(d_n\) is the sequence of alphabet sizes and \(\epsilon_n\) is the sequence of public overlap floors. The attaining estimator supplies sufficiency, and the universal lower bound supplies necessity.

\paragraph{Verification scope.}
The formal theorem statements and their supplied proofs are machine-checked in Lean 4. The sampling, overlap, and causal assumptions specify the models to which those statements apply. The theorem-local verification disclosures separately identify any cited dependency inputs and their trust boundaries.


\section{Proofs of the main results}\label{sec:deferred-proofs}

\begin{proof}[Proof of \cref{obj:prop:identification}]
\begin{enumerate}
\item \emph{Completion.}
Fix \(\mathbb P\in\mathcal D_{d,\epsilon}^{\mathrm{obs}}\), and use its cell quantities from \cref{obj:def:observed-class}. Define the Bernoulli probability weights by
\[
\operatorname{Bern}(\vartheta;y)
=
\begin{cases}
1-\vartheta,&y=0,\\
\vartheta,&y=1,
\end{cases}
\qquad
\vartheta\in[0,1],\quad y\in\{0,1\}.
\]
For \(x\in[d]\) and \(a,y,y_0,y_1\in\{0,1\}\), define the full-data atom probabilities by
\[
\begin{aligned}
&P(X=x,A=a,Y=y,Y(0)=y_0,Y(1)=y_1)\\
&\qquad=
\begin{cases}
q_{ay,x}\operatorname{Bern}(\mu_{1-a,x};y_{1-a}),&y=y_a,\\
0,&y\ne y_a,
\end{cases}
\end{aligned}
\]
where
\[
y_a=
\begin{cases}
y_0,&a=0,\\
y_1,&a=1,
\end{cases}
\qquad
y_{1-a}=
\begin{cases}
y_1,&a=0,\\
y_0,&a=1.
\end{cases}
\]
The bounds \(q_{ay,x}\ge0\) and \(\mu_{ax}\in[0,1]\) from \cref{obj:def:observed-class} make these masses nonnegative. Summing first over the potential-outcome coordinates gives
\[
\begin{aligned}
&\sum_{x\in[d]}\sum_{a,y,y_0,y_1\in\{0,1\}}
P(X=x,A=a,Y=y,Y(0)=y_0,Y(1)=y_1)\\
&\quad=
\sum_{x\in[d]}\sum_{a,y\in\{0,1\}}
q_{ay,x}
\bigl\{
\operatorname{Bern}(\mu_{1-a,x};0)
+\operatorname{Bern}(\mu_{1-a,x};1)
\bigr\}\\
&\quad=\sum_{x\in[d]}\sum_{a,y\in\{0,1\}}q_{ay,x}=1.
\end{aligned}
\]
Thus the construction defines a probability law. Its atoms with \(y\ne y_a\) have zero mass, so
\[
P\{Y=Y(A)\}=1,
\]
which is \cref{obj:ass:consistency}.
% lean: CausalSmith.Stat.OptvalueVanishingoverlapRate.completionFullMass_nonneg_sum
% lean: CausalSmith.Stat.OptvalueVanishingoverlapRate.bernoulliCompletion_consistency

Summing over the observed outcome, and then using the observed-atom factorization supplied by \cref{obj:def:observed-class}, yields
\[
\begin{aligned}
&P(X=x,A=a,Y(0)=y_0,Y(1)=y_1)\\
&\quad=q_{a\,y_a,x}\operatorname{Bern}(\mu_{1-a,x};y_{1-a})\\
&\quad=p_x\operatorname{Bern}(\pi_x;a)
\operatorname{Bern}(\mu_{0x};y_0)
\operatorname{Bern}(\mu_{1x};y_1).
\end{aligned}
\]
Consequently, summing these factors over the remaining binary coordinates gives, for \(a,a',y\in\{0,1\}\),
\[
\begin{aligned}
P(X=x)&=p_x,\\
P(X=x,A=a')&=p_x\operatorname{Bern}(\pi_x;a'),\\
P(X=x,Y(a)=y)&=p_x\operatorname{Bern}(\mu_{ax};y),\\
P(X=x,A=a',Y(a)=y)
&=p_x\operatorname{Bern}(\pi_x;a')
\operatorname{Bern}(\mu_{ax};y).
\end{aligned}
\]
Multiplication therefore gives
\[
\begin{aligned}
&P(X=x,A=a',Y(a)=y)\,P(X=x)\\
&\qquad=P(X=x,A=a')\,P(X=x,Y(a)=y).
\end{aligned}
\]
On positive-mass cells, division by \(p_x^2\) proves \cref{obj:ass:conditional-exchangeability}; on null cells both products are zero.
% lean: CausalSmith.Stat.OptvalueVanishingoverlapRate.bernoulliCompletion_poAtom_factorization
% lean: CausalSmith.Stat.OptvalueVanishingoverlapRate.bernoulliCompletion_exchangeability

For each observed atom, the construction also gives
\[
\begin{aligned}
P(X=x,A=a,Y=y)
&=q_{ay,x}
\bigl\{
\operatorname{Bern}(\mu_{1-a,x};0)
+\operatorname{Bern}(\mu_{1-a,x};1)
\bigr\}\\
&=q_{ay,x}.
\end{aligned}
\]
Hence
\[
\operatorname{Obs}(P)=\mathbb P.
\]
The observed marginal inherits the overlap restriction of \(\mathbb P\). Together with consistency and conditional exchangeability, this puts \(P\) in the class of \cref{obj:def:causal-class}, proving the completion assertion.
% lean: CausalSmith.Stat.OptvalueVanishingoverlapRate.bernoulliCompletion_observedMarginal

\item \emph{Value identification.}
Fix \(P\in\mathcal P_{d,\epsilon}\), and use \(q_{ay,x},p_x,\pi_x,\mu_{ax}\) for the quantities of its observed marginal defined in \cref{obj:def:observed-class}. Marginalization gives
\[
P(X=x)
=\sum_{a,y\in\{0,1\}}
\operatorname{Obs}(P)(X=x,A=a,Y=y)
=p_x.
\]
We compare the two value sums cell by cell. If \(p_x=0\), their cell contributions are both zero:
\[
p_x\max_{a\in\{0,1\}}\mathbb E_P[Y(a)\mid X=x]
=
p_x\max_{a\in\{0,1\}}\mu_{ax}
=0.
\]
Here the conditional potential-outcome means use the stated zero convention.
% lean: CausalSmith.Stat.OptvalueVanishingoverlapRate.observedMarginal_cellMass

Suppose now that \(p_x>0\). The lower overlap bound in \cref{obj:ass:shrinking-overlap} gives
\[
q_{10,x}+q_{11,x}=p_x\pi_x\ge\epsilon p_x>0.
\]
The decomposition of the cell mass and the upper overlap bound give
\[
\begin{aligned}
p_x&=(q_{00,x}+q_{01,x})+(q_{10,x}+q_{11,x}),\\
q_{10,x}+q_{11,x}&\le(1-\epsilon)p_x,
\end{aligned}
\]
and hence
\[
q_{00,x}+q_{01,x}
=p_x-(q_{10,x}+q_{11,x})
\ge\epsilon p_x>0.
\]
Thus both arm masses are positive.

For either \(a\in\{0,1\}\), \cref{obj:ass:conditional-exchangeability} supplies
\[
P(X=x,A=a,Y(a)=1)\,p_x
=
P(X=x,A=a)\,P(X=x,Y(a)=1).
\]
Consistency, as specified in \cref{obj:ass:consistency}, identifies the selected atoms:
\[
P(X=x,A=a,Y(a)=y)=q_{ay,x},
\qquad y\in\{0,1\}.
\]
Summing this identity over \(y\), and substituting its outcome-one instance into the exchangeability identity, gives
\[
P(X=x,A=a)=q_{a0,x}+q_{a1,x},
\qquad
q_{a1,x}p_x
=(q_{a0,x}+q_{a1,x})P(X=x,Y(a)=1).
\]
Division by the positive cell and arm masses now yields
\[
\mathbb E_P[Y(a)\mid X=x]
=
\frac{P(X=x,Y(a)=1)}{p_x}
=
\frac{q_{a1,x}}{q_{a0,x}+q_{a1,x}}
=\mu_{ax}.
\]
Combining the positive- and zero-mass cases and summing over \(x\) proves
\[
\begin{aligned}
V^\star(P)
&=\sum_{x\in[d]}P(X=x)
\max_{a\in\{0,1\}}\mathbb E_P[Y(a)\mid X=x]\\
&=\sum_{x\in[d]}p_x\max_{a\in\{0,1\}}\mu_{ax}
=\Psi\{\operatorname{Obs}(P)\},
\end{aligned}
\]
where the last equality is the definition in \cref{obj:def:observed-value}.
% lean: CausalSmith.Stat.OptvalueVanishingoverlapRate.poRegression_eq_outcomeMean_of_pos
% lean: CausalSmith.Stat.OptvalueVanishingoverlapRate.identification

\item \emph{The armwise extension on observed cells.}
Fix \(\mathbb P\in\mathcal D_{d,\epsilon}^{\mathrm{obs}}\) and \(x\in[d]\). The cell vector and the quantities in \cref{obj:def:armwise-extension} satisfy
\[
s(q_x)=p_x,\qquad
s_a(q_x)=q_{a0,x}+q_{a1,x},\qquad
(q_x)_{a1}=q_{a1,x},
\qquad a\in\{0,1\}.
\]
Every coordinate is nonnegative, so \(p_x\ge0\). If \(p_x=0\), then
\[
0\le q_{ay,x}\le p_x=0
\qquad(a,y\in\{0,1\}),
\]
and therefore \(q_x=0\). The zero branch of \cref{obj:def:armwise-extension} gives
\[
F_\epsilon(q_x)=0
=p_x\max_{a\in\{0,1\}}\mu_{ax}.
\]

Suppose instead that \(p_x>0\). The lower overlap bound gives
\[
\epsilon p_x\le q_{10,x}+q_{11,x}=s_1(q_x).
\]
For the other arm, the upper overlap bound and the cell-mass decomposition give
\[
s_1(q_x)\le(1-\epsilon)p_x,
\qquad
s_0(q_x)=p_x-s_1(q_x)\ge\epsilon p_x.
\]
Thus, for each \(a\in\{0,1\}\),
\[
s_a(q_x)\ge\epsilon p_x>0,
\qquad
D_a(q_x)=\max\{s_a(q_x),\epsilon p_x\}=s_a(q_x).
\]
Each armwise term consequently reduces to
\[
\frac{s(q_x)(q_x)_{a1}}{D_a(q_x)}
=
\frac{p_xq_{a1,x}}{q_{a0,x}+q_{a1,x}}
=
p_x\mu_{ax}.
\]
Since \(s(q_x)=p_x>0\), the vector \(q_x\) is nonzero. Using the nonzero branch of \cref{obj:def:armwise-extension} and factoring the nonnegative \(p_x\) out of the maximum gives
\[
F_\epsilon(q_x)
=
\max_{a\in\{0,1\}}p_x\mu_{ax}
=
p_x\max_{a\in\{0,1\}}\mu_{ax}.
\]
This proves the displayed armwise identity, including its null-cell case.
% lean: CausalSmith.Stat.OptvalueVanishingoverlapRate.armwiseExtension_cellVector
\end{enumerate}
\end{proof}

\begin{proof}[Proof of \cref{obj:thm:matched-frontier}]
\begin{enumerate}
\item Choose the universal constants \(c_\star,C_\star,H_0,\kappa,D_0\) supplied by the risk guarantees in \cref{obj:thm:weighted-separation-and-frontier}. They satisfy
\[
0<c_\star\le C_\star,\qquad H_0>0,\qquad \kappa>0,\qquad D_0\ge2.
\]
Fix integers \(n\ge1\), \(d\ge2\), and an overlap parameter \(0<\epsilon\le1/2\). Define
\[
L_d=\log(ed),
\qquad
\rho_{n,d,\epsilon}
=\min\left\{1,\frac{d}{n\epsilon L_d}\right\}.
\]
For the estimator of \cref{obj:def:armwise-estimator} with the chosen calibrations, \cref{obj:thm:weighted-separation-and-frontier} gives
\[
c_\star\rho_{n,d,\epsilon}
\le \mathfrak R_{n,d,\epsilon}
\le
\sup_{\mathbb P\in\mathcal D_{d,\epsilon}^{\mathrm{obs}}}
E_{\mathbb P}\!\left[
\bigl(\widehat V_{n,d,\epsilon}^{\mathrm{AF}}-\Psi(\mathbb P)\bigr)^2
\right]
\le C_\star\rho_{n,d,\epsilon},
\]
and
\[
c_\star\rho_{n,d,\epsilon}
\le
\inf_{\widehat V}
\sup_{P\in\mathcal P_{d,\epsilon}}
E_{\operatorname{Obs}(P)}\!\left[
\bigl(\widehat V-V^\star(P)\bigr)^2
\right]
\le C_\star\rho_{n,d,\epsilon}.
\]
The infimum is over measurable real-valued estimators based on the observed sample. Thus the observed chain and the two causal minimax bounds are already established with the same universal constants.
% lean: CausalSmith.Stat.OptvalueVanishingoverlapRate.weighted_separation_and_frontier

\item We next establish equality of the causal and observed worst-case risks for each estimator. Fix any measurable function
\[
\widehat V:
\bigl([d]\times\{0,1\}\times\{0,1\}\bigr)^n
\longrightarrow\mathbb R.
\]
By \cref{obj:prop:identification}, every \(P\in\mathcal P_{d,\epsilon}\) satisfies
\[
V^\star(P)=\Psi\{\operatorname{Obs}(P)\}.
\]
Consequently,
\[
E_{\operatorname{Obs}(P)}\!\left[
\bigl(\widehat V-V^\star(P)\bigr)^2
\right]
=
E_{\operatorname{Obs}(P)}\!\left[
\bigl(\widehat V-\Psi\{\operatorname{Obs}(P)\}\bigr)^2
\right].
\]
Here both expectations use the independent sample of size \(n\) from the observed marginal, as in \cref{obj:ass:iid-sampling}.
% lean: CausalSmith.Stat.OptvalueVanishingoverlapRate.causalRisk_eq_observedRisk

Every causal law has an observed marginal in \(\mathcal D_{d,\epsilon}^{\mathrm{obs}}\) by \cref{obj:def:causal-class}. Conversely, \cref{obj:prop:identification} supplies a causal completion for every law in that observed class. The displayed risk identity therefore gives both inclusions in the equality of attained risk sets:
\[
\begin{aligned}
&\left\{
E_{\operatorname{Obs}(P)}\!\left[
\bigl(\widehat V-V^\star(P)\bigr)^2
\right]:
P\in\mathcal P_{d,\epsilon}
\right\}
\\
&\qquad=
\left\{
E_{\mathbb P}\!\left[
\bigl(\widehat V-\Psi(\mathbb P)\bigr)^2
\right]:
\mathbb P\in\mathcal D_{d,\epsilon}^{\mathrm{obs}}
\right\}.
\end{aligned}
\]
% lean: CausalSmith.Stat.OptvalueVanishingoverlapRate.causalRisk_range_eq_observedRisk_range
Taking suprema yields
\[
\sup_{P\in\mathcal P_{d,\epsilon}}
E_{\operatorname{Obs}(P)}\!\left[
\bigl(\widehat V-V^\star(P)\bigr)^2
\right]
=
\sup_{\mathbb P\in\mathcal D_{d,\epsilon}^{\mathrm{obs}}}
E_{\mathbb P}\!\left[
\bigl(\widehat V-\Psi(\mathbb P)\bigr)^2
\right].
\]
% lean: CausalSmith.Stat.OptvalueVanishingoverlapRate.causalWorstRisk_eq_observedWorstRisk

\item Taking the infimum of this identity over the common estimator class gives
\[
\inf_{\widehat V}
\sup_{P\in\mathcal P_{d,\epsilon}}
E_{\operatorname{Obs}(P)}\!\left[
\bigl(\widehat V-V^\star(P)\bigr)^2
\right]
=\mathfrak R_{n,d,\epsilon}.
\]
% lean: CausalSmith.Stat.OptvalueVanishingoverlapRate.causalMinimaxRisk_eq_minimaxRisk
Applying the same identity to the prescribed armwise estimator gives
\[
\begin{aligned}
&\sup_{P\in\mathcal P_{d,\epsilon}}
E_{\operatorname{Obs}(P)}\!\left[
\bigl(\widehat V_{n,d,\epsilon}^{\mathrm{AF}}-V^\star(P)\bigr)^2
\right]
\\
&\qquad=
\sup_{\mathbb P\in\mathcal D_{d,\epsilon}^{\mathrm{obs}}}
E_{\mathbb P}\!\left[
\bigl(\widehat V_{n,d,\epsilon}^{\mathrm{AF}}-\Psi(\mathbb P)\bigr)^2
\right].
\end{aligned}
\]
% lean: CausalSmith.Stat.OptvalueVanishingoverlapRate.causalArmwiseWorstRisk_eq_armwiseWorstRisk
Substituting these two equalities into the middle observed inequality from the first step, and then substituting the second equality into its upper inequality, yields
\[
\begin{aligned}
\inf_{\widehat V}
\sup_{P\in\mathcal P_{d,\epsilon}}
E_{\operatorname{Obs}(P)}\!\left[
\bigl(\widehat V-V^\star(P)\bigr)^2
\right]
&\le
\sup_{P\in\mathcal P_{d,\epsilon}}
E_{\operatorname{Obs}(P)}\!\left[
\bigl(\widehat V_{n,d,\epsilon}^{\mathrm{AF}}-V^\star(P)\bigr)^2
\right]
\\
&\le C_\star\rho_{n,d,\epsilon}.
\end{aligned}
\]
Together with the causal lower bound from the first step, these are the asserted causal inequalities. Since \(n,d,\epsilon\) were arbitrary in the stated range, all bounds hold with the chosen universal constants.
% lean: CausalSmith.Stat.OptvalueVanishingoverlapRate.matched_frontier
\end{enumerate}
\end{proof}

\begin{proof}[Proof of \cref{obj:thm:consistency-threshold}]
\begin{enumerate}
\item Fix the universal constants and estimator supplied by
\cref{obj:thm:matched-frontier}. For integers \(n\ge1\), \(d\ge2\), and
\(0<\epsilon\le1/2\), write
\[
L_d:=\log(ed),
\qquad
\rho_{n,d,\epsilon}
:=\min\left\{1,\frac{d}{n\epsilon L_d}\right\}.
\]
The supplied constants satisfy \(0<c_\star\le C_\star\), and the observed-risk bounds are
\[
c_\star\rho_{n,d,\epsilon}
\le \mathfrak R_{n,d,\epsilon}
\le
\sup_{\mathbb P\in\mathcal D_{d,\epsilon}^{\mathrm{obs}}}
E_{\mathbb P}\!\left[
\bigl(\widehat V_{n,d,\epsilon}^{\mathrm{AF}}-\Psi(\mathbb P)\bigr)^2
\right]
\le C_\star\rho_{n,d,\epsilon}.
\]
% lean: CausalSmith.Stat.OptvalueVanishingoverlapRate.matched_frontier

\item Apply the final equivalence in
\cref{obj:thm:weighted-separation-and-frontier} with the sample size equal to the sequence index. These sample sizes tend to infinity, and the stipulated alphabet and overlap conditions satisfy its remaining hypotheses. It gives precisely
\[
\begin{gathered}
\exists\,(\widehat V_n)_{n\in\mathbb N}
\text{ of measurable real-valued estimators based on \(n\) observations}:\\
\sup_{\mathbb P\in\mathcal D_{d_n,\epsilon_n}^{\mathrm{obs}}}
E_{\mathbb P}\!\left[
\bigl(\widehat V_n-\Psi(\mathbb P)\bigr)^2
\right]\longrightarrow0
\end{gathered}
\quad\Longleftrightarrow\quad
\frac{d_n}{n\epsilon_n L_{d_n}}\longrightarrow0.
\]
Here the ratio is evaluated for \(n\ge1\); the limit is unchanged by the initial index. At \(n=0\), the estimator may be assigned the constant value
\(\widehat V_0=0\). Since \(L_{d_n}=\log(ed_n)\), this proves the consistency assertion.
% lean: CausalSmith.Stat.OptvalueVanishingoverlapRate.weighted_separation_and_frontier

\item Fix an integer \(d\ge2\). The logarithmic bounds
\[
0\le\log d\le d-1
\]
follow respectively from \(d\ge1\) and the standard inequality
\(\log d\le d-1\) for \(d>0\). Consequently,
\[
L_d=1+\log d,
\qquad
0<L_d\le d,
\qquad
\lambda_d:=\frac{d}{L_d}\ge1.
\]
Choose
\[
c_d:=c_\star,
\qquad
C_d:=C_\star\frac{d}{L_d}.
\]
Because \(C_\star\ge c_\star>0\) and \(\lambda_d\ge1\),
\[
0<c_d=c_\star\le C_\star
\le C_\star\lambda_d=C_d.
\]

We will use the elementary comparison
\[
\min\{1,\xi_{\mathrm{cmp}}\}
\le \min\{1,\lambda_{\mathrm{cmp}}\xi_{\mathrm{cmp}}\}
\le \lambda_{\mathrm{cmp}}\min\{1,\xi_{\mathrm{cmp}}\},
\qquad
\lambda_{\mathrm{cmp}}\in[1,\infty),\quad
\xi_{\mathrm{cmp}}\in[0,\infty).
\]
For the first inequality,
\(\xi_{\mathrm{cmp}}\le\lambda_{\mathrm{cmp}}\xi_{\mathrm{cmp}}\), so monotonicity of the minimum gives
\[
\min\{1,\xi_{\mathrm{cmp}}\}
\le\min\{1,\lambda_{\mathrm{cmp}}\xi_{\mathrm{cmp}}\}.
\]
For the second inequality, if \(\xi_{\mathrm{cmp}}\le1\), then
\[
\min\{1,\lambda_{\mathrm{cmp}}\xi_{\mathrm{cmp}}\}
\le\lambda_{\mathrm{cmp}}\xi_{\mathrm{cmp}}
=\lambda_{\mathrm{cmp}}\min\{1,\xi_{\mathrm{cmp}}\}.
\]
If \(\xi_{\mathrm{cmp}}\ge1\), then
\[
\min\{1,\lambda_{\mathrm{cmp}}\xi_{\mathrm{cmp}}\}
\le1\le\lambda_{\mathrm{cmp}}
=\lambda_{\mathrm{cmp}}\min\{1,\xi_{\mathrm{cmp}}\}.
\]
These two cases include the common boundary \(\xi_{\mathrm{cmp}}=1\).

Now let \(n\ge1\) and \(0<\epsilon\le1/2\), and define
\[
\xi_{n,\epsilon}:=\frac{1}{n\epsilon}\ge0.
\]
The denominator is positive, and direct rearrangement gives
\[
\rho_{n,d,\epsilon}
=\min\{1,\lambda_d\xi_{n,\epsilon}\}.
\]
Applying the comparison with
\(\lambda_{\mathrm{cmp}}=\lambda_d\) and
\(\xi_{\mathrm{cmp}}=\xi_{n,\epsilon}\), the frontier bounds yield
\[
c_d\min\{1,\xi_{n,\epsilon}\}
\le c_\star\min\{1,\lambda_d\xi_{n,\epsilon}\}
\le\mathfrak R_{n,d,\epsilon}
\]
and
\[
\begin{aligned}
\mathfrak R_{n,d,\epsilon}
&\le
\sup_{\mathbb P\in\mathcal D_{d,\epsilon}^{\mathrm{obs}}}
E_{\mathbb P}\!\left[
\bigl(\widehat V_{n,d,\epsilon}^{\mathrm{AF}}-\Psi(\mathbb P)\bigr)^2
\right]\\
&\le C_\star\min\{1,\lambda_d\xi_{n,\epsilon}\}\\
&\le C_\star\lambda_d\min\{1,\xi_{n,\epsilon}\}
=C_d\min\left\{1,\frac{1}{n\epsilon}\right\}.
\end{aligned}
\]
Thus \(c=c_d\) and \(C=C_d\) give the fixed-alphabet bounds throughout the required range.
% lean: CausalSmith.Stat.OptvalueVanishingoverlapRate.logAlphabet_bounds
% lean: CausalSmith.Stat.OptvalueVanishingoverlapRate.min_one_mul_bounds
% lean: CausalSmith.Stat.OptvalueVanishingoverlapRate.consistency_threshold

\item Fix \(0<\epsilon\le1/2\), and choose
\[
c_\epsilon:=c_\star,
\qquad
C_\epsilon:=\frac{C_\star}{\epsilon}.
\]
Since \(\epsilon>0\) and \(\epsilon\le1/2\le1\),
\[
\lambda_\epsilon:=\frac1\epsilon\ge1.
\]
It follows that
\[
0<c_\epsilon=c_\star\le C_\star
\le C_\star\lambda_\epsilon=C_\epsilon.
\]
For integers \(n\ge1\) and \(d\ge2\), the logarithmic bounds established above give \(L_d>0\). Define
\[
\xi_{n,d}:=\frac{d}{nL_d}\ge0.
\]
Then
\[
\rho_{n,d,\epsilon}
=\min\{1,\lambda_\epsilon\xi_{n,d}\}.
\]
Applying the same elementary comparison with
\(\lambda_{\mathrm{cmp}}=\lambda_\epsilon\) and
\(\xi_{\mathrm{cmp}}=\xi_{n,d}\), we obtain
\[
c_\epsilon\min\{1,\xi_{n,d}\}
\le c_\star\min\{1,\lambda_\epsilon\xi_{n,d}\}
\le\mathfrak R_{n,d,\epsilon}
\]
and
\[
\begin{aligned}
\mathfrak R_{n,d,\epsilon}
&\le
\sup_{\mathbb P\in\mathcal D_{d,\epsilon}^{\mathrm{obs}}}
E_{\mathbb P}\!\left[
\bigl(\widehat V_{n,d,\epsilon}^{\mathrm{AF}}-\Psi(\mathbb P)\bigr)^2
\right]\\
&\le C_\star\min\{1,\lambda_\epsilon\xi_{n,d}\}\\
&\le C_\star\lambda_\epsilon\min\{1,\xi_{n,d}\}
=C_\epsilon\min\left\{1,\frac{d}{n\log(ed)}\right\}.
\end{aligned}
\]
Taking \(c=c_\epsilon\) and \(C=C_\epsilon\) proves the fixed-overlap bounds.
% lean: CausalSmith.Stat.OptvalueVanishingoverlapRate.consistency_threshold
\end{enumerate}
\end{proof}

\begin{proof}[Proof of \cref{obj:thm:observable-upper}]
\textit{1. Pilot calibration and auxiliary experiment.}
Choose
\[
H_0=1025.
\]
Thus \(H_0>0\), \(H_0\ge1\), and \(H_0\ge1024\), the universal self-normalization multiplier.

We first establish the auxiliary risk bound for integers \(n,d\ge1\), an overlap parameter \(0<\epsilon\le1/2\), and a law \(\mathbb P\in\mathcal D_{d,\epsilon}^{\mathrm{obs}}\). Write
\[
L=L_d=\log(ed),\qquad m=\frac n8,\qquad
\delta_{\mathrm p}=\frac Lm,\qquad
q_x=(q_{ay,x})_{a,y\in\{0,1\}}.
\]
The atom probabilities are those of \cref{obj:def:observed-class}; in particular, \(m>0\) and \(L=1+\log d\ge1\).

Consider the auxiliary experiment that draws \(\mathsf M\sim\operatorname{Pois}(n/4)\) and, conditional on \(\mathsf M\), draws \(\mathsf M\) independent observations with law \(\mathbb P\), together with independent fair marks. Its law on the disjoint union of all finite marked samples is
\[
\mathcal Q_{\mathbb P}
=
\sum_{k=0}^{\infty}
e^{-n/4}\frac{(n/4)^k}{k!}
\left(\mathbb P\otimes\operatorname{Bernoulli}(1/2)\right)^{\otimes k}.
\]
Write \(\mathbb E_{\mathrm{aux}}\) for expectation under this law. Construct the counts and cell statistics of \cref{obj:def:armwise-estimator} for every finite marked sample, with the auxiliary count unrestricted. For arrays
\(\boldsymbol t^{\mathrm p},\boldsymbol t^{\mathrm e}\in[0,1]^{\mathcal J_d}\), indexed by the observed atom set defined in \cref{obj:synth_1}, the count generating function is
\[
\mathbb E_{\mathrm{aux}}
\prod_{x,a,y}
(t^{\mathrm p}_{ay,x})^{N'_{ay,x}}
(t^{\mathrm e}_{ay,x})^{N_{ay,x}}
=
\exp\left\{
m\sum_{x,a,y}q_{ay,x}(t^{\mathrm p}_{ay,x}-1)
+
m\sum_{x,a,y}q_{ay,x}(t^{\mathrm e}_{ay,x}-1)
\right\}.
\]
Indeed, conditional on the auxiliary count, each marked observation contributes one category factor; summing over the Poisson count gives the displayed exponential. Consequently,
\[
N'_{ay,x},N_{ay,x}\sim\operatorname{Pois}(mq_{ay,x}),
\]
and all these counts are mutually independent. In particular, the complete cell statistics \(T_x\) are independent across distinct cells.
% lean: CausalSmith.Stat.OptvalueVanishingoverlapRate.markedCount_poisson_hasLaw
% lean: CausalSmith.Stat.OptvalueVanishingoverlapRate.markedPoissonCellValue_pairwise_indepFun

\textit{2. Polynomial oscillation and coefficient control.}
Use the centers, radii, rectangles, polynomial, and clipping scales of \cref{obj:def:armwise-estimator}. Every coordinate radius is positive, since
\[
h_{ay,x}=H_0\left(\sqrt{c_{ay,x}\delta_{\mathrm p}}+\delta_{\mathrm p}\right)>0,
\qquad
\frac{h_{ay,x}}2\le r_{ay,x}\le h_{ay,x}.
\]
Moreover,
\[
c_{ay,x}\le b_{ay,x}\le c_{ay,x}+h_{ay,x}.
\]
Thus \(b_x\ne0\). Applying \cref{obj:lem:armwise-modulus} with comparison vector \(b_x\), for every \(v\in Q_x\), gives
\[
|F_\epsilon(v)-F_\epsilon(b_x)|
\le
2\sum_{a=0}^1
\left(1+\frac{s(b_x)}{D_a(b_x)}\right)R_{a,x}
=S_x^{\mathrm{aw}}.
\]
The Jackson convolution is positive and normalized, so the same bound holds for its polynomial:
\[
\sup_{v\in Q_x}
|P_{x,K_n}(v)-F_\epsilon(b_x)|
\le S_x^{\mathrm{aw}}.
\]

For completeness, define the Chebyshev polynomials by
\[
\mathsf{Ch}_0(\xi_\circ)=1,\qquad
\mathsf{Ch}_1(\xi_\circ)=\xi_\circ,\qquad
\mathsf{Ch}_{j+1}(\xi_\circ)
=2\xi_\circ\mathsf{Ch}_j(\xi_\circ)-\mathsf{Ch}_{j-1}(\xi_\circ),
\quad j\ge1.
\]
Here \(\xi_\circ\) is a real scalar. The recurrence implies
\[
\sum_{\ell=0}^{j}
\left|[\xi_\circ^\ell]\mathsf{Ch}_j(\xi_\circ)\right|
\le(1+\sqrt2)^j,
\]
by induction and the identity
\((1+\sqrt2)^2=2(1+\sqrt2)+1\).

With \(D=2(K_n-1)\), expand the normalized polynomial as
\[
P_{x,K_n}(b_x+r_x\odot\xi)-F_\epsilon(b_x)
=
\sum_{\boldsymbol\nu\in\{0,\ldots,D\}^4}
\gamma_{x,\boldsymbol\nu}
\prod_{a,y}\mathsf{Ch}_{\nu_{ay}}(\xi_{ay}),
\qquad \xi\in[-1,1]^4.
\]
Cosine orthogonality expresses each coefficient as an integral of the bounded polynomial against cosine factors, with normalization factor at most \(16\). Hence
\[
|\gamma_{x,\boldsymbol\nu}|\le16S_x^{\mathrm{aw}}.
\]
Converting the tensor Chebyshev expansion to the monomial expansion defining \(\beta_{x,\boldsymbol\alpha}\) yields
\[
\sum_{\boldsymbol\alpha}|\beta_{x,\boldsymbol\alpha}|
\le
16(D+1)^4(1+\sqrt2)^{4D}S_x^{\mathrm{aw}}.
\]
The numerical factor satisfies
\[
\left\{16(D+1)^4(1+\sqrt2)^{4D}\right\}^2
\le e^{256+24D}\le e^{304K_n}.
\]
Indeed, \(D+1\le e^D\), \(1+\sqrt2\le e^2\), \(256\le e^{256}\), and \(D\le2K_n\), \(K_n\ge1\).
% lean: CausalSmith.Stat.OptvalueVanishingoverlapRate.jacksonCoeff_sum_abs_le
% lean: CausalSmith.Stat.OptvalueVanishingoverlapRate.jacksonCoeff_factor_square_le_exp

\textit{3. Factorial second moments and universal degree tuning.}
Define the successful-localization event by
\[
\mathcal G_x=\{q_x\in Q_x\}.
\]
On this event,
\[
|q_{ay,x}-b_{ay,x}|\le r_{ay,x},
\qquad
\frac{q_{ay,x}}{m r_{ay,x}^2}\le\frac8L.
\]
To verify the second inequality, the coordinate formulas give
\[
c_{ay,x}\delta_{\mathrm p}\le h_{ay,x}^2,\qquad
\delta_{\mathrm p}\le h_{ay,x},\qquad
r_{ay,x}\ge\frac{h_{ay,x}}2,
\]
because \(H_0\ge1\). Since \(q_{ay,x}\le c_{ay,x}+h_{ay,x}\),
\[
q_{ay,x}\delta_{\mathrm p}
\le(c_{ay,x}+h_{ay,x})\delta_{\mathrm p}
\le2h_{ay,x}^2
\le8r_{ay,x}^2,
\]
which proves the claim.

For a Poisson count \(\mathsf N_\circ\) with mean \(m q_\circ\), where \(q_\circ\ge0\), the factorial identities are
\[
\mathbb E(\mathsf N_\circ)_j=(m q_\circ)^j,
\]
\[
(\mathsf N_\circ)_j(\mathsf N_\circ)_k
=
\sum_{\ell=0}^{\min(j,k)}
\binom j\ell\binom k\ell\ell!\,
(\mathsf N_\circ)_{j+k-\ell},
\qquad j,k\in\mathbb N.
\]
The first follows by shifting the Poisson series. The second counts the overlap of two ordered selections of distinct indices. Substituting the first identity into the centered factorial statistic of \cref{obj:def:armwise-estimator}, for every real \(\zeta\), and applying the binomial theorem gives
\[
\mathbb E U_j(\mathsf N_\circ;\zeta)
=
\sum_{j_{\mathrm{raw}}=0}^{j}
\binom j{j_{\mathrm{raw}}}(-\zeta)^{j-j_{\mathrm{raw}}}
q_\circ^{j_{\mathrm{raw}}}
=(q_\circ-\zeta)^j.
\]
For the mixed second moment, substitution of the raw factorial product identity gives the finite triple sum
\[
\begin{aligned}
\mathbb E[U_j(\mathsf N_\circ;\zeta)U_k(\mathsf N_\circ;\zeta)]
={}&
\sum_{j_{\mathrm{raw}}=0}^{j}
\sum_{k_{\mathrm{raw}}=0}^{k}
\binom j{j_{\mathrm{raw}}}\binom k{k_{\mathrm{raw}}}
(-\zeta)^{j-j_{\mathrm{raw}}+k-k_{\mathrm{raw}}}\\
&\quad\times
\sum_{\ell=0}^{\min(j_{\mathrm{raw}},k_{\mathrm{raw}})}
\binom{j_{\mathrm{raw}}}\ell
\binom{k_{\mathrm{raw}}}\ell\ell!
\left(\frac{q_\circ}{m}\right)^\ell
q_\circ^{j_{\mathrm{raw}}+k_{\mathrm{raw}}-2\ell}.
\end{aligned}
\]
Fixing the overlap order and regrouping the finite sums factors this expression as
\[
\begin{aligned}
\mathbb E[U_j(\mathsf N_\circ;\zeta)U_k(\mathsf N_\circ;\zeta)]
={}&\sum_{\ell=0}^{\min(j,k)}
\ell!\left(\frac{q_\circ}{m}\right)^\ell\\
&\times\left\{
\sum_{j_{\mathrm{raw}}=\ell}^{j}
\binom j{j_{\mathrm{raw}}}\binom{j_{\mathrm{raw}}}\ell
(-\zeta)^{j-j_{\mathrm{raw}}}
q_\circ^{j_{\mathrm{raw}}-\ell}
\right\}\\
&\times\left\{
\sum_{k_{\mathrm{raw}}=\ell}^{k}
\binom k{k_{\mathrm{raw}}}\binom{k_{\mathrm{raw}}}\ell
(-\zeta)^{k-k_{\mathrm{raw}}}
q_\circ^{k_{\mathrm{raw}}-\ell}
\right\}.
\end{aligned}
\]
Terms with a raw order smaller than \(\ell\) vanish because the corresponding binomial coefficient is zero. For the remaining terms, the shifted binomial identity is
\[
\begin{aligned}
\sum_{j_{\mathrm{raw}}=\ell}^{j}
\binom j{j_{\mathrm{raw}}}\binom{j_{\mathrm{raw}}}\ell
(-\zeta)^{j-j_{\mathrm{raw}}}
q_\circ^{j_{\mathrm{raw}}-\ell}
&=
\binom j\ell
\sum_{i_{\mathrm{shift}}=0}^{j-\ell}
\binom{j-\ell}{i_{\mathrm{shift}}}
(-\zeta)^{j-\ell-i_{\mathrm{shift}}}
q_\circ^{i_{\mathrm{shift}}}\\
&=\binom j\ell(q_\circ-\zeta)^{j-\ell},
\qquad 0\le\ell\le j.
\end{aligned}
\]
The first equality substitutes \(j_{\mathrm{raw}}=\ell+i_{\mathrm{shift}}\) and uses
\[
\binom j{\ell+i_{\mathrm{shift}}}
\binom{\ell+i_{\mathrm{shift}}}\ell
=
\binom j\ell\binom{j-\ell}{i_{\mathrm{shift}}},
\qquad
 i_{\mathrm{shift}}\in\{0,\ldots,j-\ell\};
\]
the second is the binomial theorem. Applying this identity to both factors yields
\[
\mathbb E[U_j(\mathsf N_\circ;\zeta)U_k(\mathsf N_\circ;\zeta)]
=
\sum_{\ell=0}^{\min(j,k)}
\binom j\ell\binom k\ell\ell!
\left(\frac{q_\circ}{m}\right)^\ell
(q_\circ-\zeta)^{j+k-2\ell}.
\]
Taking \(k=j\) therefore gives
\[
\mathbb E U_j(\mathsf N_\circ;\zeta)^2
=
\sum_{\ell=0}^{j}
\binom j\ell^2\ell!
\left(\frac{q_\circ}{m}\right)^\ell
(q_\circ-\zeta)^{2j-2\ell}.
\]
Thus, conditional on a pilot in \(\mathcal G_x\), for \(0\le j\le D\),
\[
\mathbb E_{\mathrm{aux}}
\left[
\left.
\frac{U_j(N_{ay,x};b_{ay,x})^2}{r_{ay,x}^{2j}}
\right|N'
\right]
\le
\sum_{\ell=0}^{j}\frac{j^{2\ell}}{\ell!}
\left(\frac8L\right)^\ell
\le e^{8j^2/L}.
\]
Here \(N'\) denotes the complete pilot-count array. Evaluation-coordinate independence therefore gives
\[
\mathbb E_{\mathrm{aux}}[Z_x\mid N']
=P_{x,K_n}(q_x),
\]
\[
\mathbb E_{\mathrm{aux}}
[(Z_x-F_\epsilon(b_x))^2\mid N']
\le
\left(\sum_{\boldsymbol\alpha}|\beta_{x,\boldsymbol\alpha}|\right)^2
e^{32D^2/L}
\quad\text{on }\mathcal G_x.
\]
For the last inequality, each normalized factorial product has second moment at most \(e^{32D^2/L}\); Cauchy--Schwarz bounds each cross term in the polynomial square by this same envelope.

Choose the numerical constants
\[
\kappa=\frac1{17952},
\qquad
C_{\mathrm{fac}}=\exp\left(1632+\frac1{16}\right).
\]
This choice satisfies
\[
304\kappa+256\kappa^2\le\frac1{16}.
\]
The degree definition gives, for every \(d\ge1\),
\[
\frac{\kappa L}{2}\le K_n\le2+\kappa L.
\]
For the lower bound, split according as \(\kappa L\le4\) or \(\kappa L>4\), and use respectively \(K_n\ge2\) or
\(\lfloor\kappa L\rfloor>\kappa L-1\). Also,
\[
K_n^2\le8+2\kappa^2L^2.
\]
Combining these inequalities with the coefficient estimate gives
\[
\begin{aligned}
\left(\sum_{\boldsymbol\alpha}|\beta_{x,\boldsymbol\alpha}|\right)^2
e^{32D^2/L}
&\le
(S_x^{\mathrm{aw}})^2
e^{304K_n+128K_n^2/L}\\
&\le
(S_x^{\mathrm{aw}})^2e^{1632+L/16}\\
&=
C_{\mathrm{fac}}d^{1/16}(S_x^{\mathrm{aw}})^2.
\end{aligned}
\]
Thus the degree tuning is uniform over the entire auxiliary alphabet range.
% lean: CausalSmith.Stat.OptvalueVanishingoverlapRate.pilotFactorialCellValue_square_envelope
% lean: CausalSmith.Stat.OptvalueVanishingoverlapRate.pilotJacksonCoeff_square_tuning_exists

\textit{4. Clipping on successful pilots.}
On \(\mathcal G_x\), the polynomial oscillation bound places \(P_{x,K_n}(q_x)\) in the clipping interval, since \(d^{1/4}\ge1\). Projection onto that interval consequently gives
\[
\mathbb E_{\mathrm{aux}}
[(T_x-P_{x,K_n}(q_x))^2\mid N']
\le
C_{\mathrm{fac}}d^{1/16}(S_x^{\mathrm{aw}})^2.
\]
Indeed, projection contracts distance from \(P_{x,K_n}(q_x)\), and
\[
\mathbb E_{\mathrm{aux}}
[(Z_x-P_{x,K_n}(q_x))^2\mid N']
\le
\mathbb E_{\mathrm{aux}}
[(Z_x-F_\epsilon(b_x))^2\mid N'].
\]
For the clipping bias, the scalar projection inequality
\[
|\Pi_{[-\mathfrak u_{\mathrm{clip}},\mathfrak u_{\mathrm{clip}}]}(\mathfrak v_{\mathrm{clip}})-\mathfrak v_{\mathrm{clip}}|
\le\frac{\mathfrak v_{\mathrm{clip}}^2}{\mathfrak u_{\mathrm{clip}}},
\qquad \mathfrak u_{\mathrm{clip}}>0,\quad
\mathfrak v_{\mathrm{clip}}\in\mathbb R,
\]
follows by considering \(|\mathfrak v_{\mathrm{clip}}|\le\mathfrak u_{\mathrm{clip}}\) and \(|\mathfrak v_{\mathrm{clip}}|>\mathfrak u_{\mathrm{clip}}\). Applying it with clipping radius \(d^{1/4}S_x^{\mathrm{aw}}>0\), and using factorial unbiasedness, yields
\[
\left|
\mathbb E_{\mathrm{aux}}[T_x\mid N']-P_{x,K_n}(q_x)
\right|
\le C_{\mathrm{fac}}d^{-3/16}S_x^{\mathrm{aw}}
\quad\text{on }\mathcal G_x.
\]
% lean: CausalSmith.Stat.OptvalueVanishingoverlapRate.pilotClippedCellValue_tuning_exists

\textit{5. Approximation error in occupied cells.}
For \(p_x>0\), define
\[
\rho_{1x}=\pi_x,\qquad \rho_{0x}=1-\pi_x.
\]
By \cref{obj:def:observed-class,obj:ass:shrinking-overlap},
\[
s(q_x)=p_x,\qquad
s_a(q_x)=p_x\rho_{ax}>0,\qquad
D_a(q_x)=s_a(q_x),\qquad
\rho_{ax}\ge\epsilon.
\]
It follows from \cref{obj:def:armwise-extension,obj:def:observed-value} that
\[
F_\epsilon(q_x)=p_x\max_a\mu_{ax},
\qquad
\Psi(\mathbb P)=\sum_{x\in[d]}F_\epsilon(q_x).
\]
The same cell identity holds when \(p_x=0\), because then \(q_x=0\).

We record the kernel moments used in the approximation. Define, for \(t\in[-\pi,\pi]\), the continuously extended sine ratio and its fourth-power mass by
\[
\mathcal S_{K_n}(t)=
\begin{cases}
\displaystyle\frac{\sin(K_nt/2)}{\sin(t/2)},&t\ne0,\\[4pt]
K_n,&t=0,
\end{cases}
\qquad
\mathcal A_{K_n}=\int_{-\pi}^{\pi}\mathcal S_{K_n}(t)^4\,dt.
\]
The finite cosine expansion is
\[
\mathcal S_{K_n}(t)^2
=K_n+2\sum_{j=1}^{K_n-1}(K_n-j)\cos(jt).
\]
Indeed, multiplying the right side by \(2(1-\cos t)\) and using the cosine product identity telescopes to \(2(1-\cos(K_nt))\); division by \(4\sin^2(t/2)\) proves the identity for \(t\ne0\), and continuity supplies its value at zero. Cosine orthogonality and the finite square sum give
\[
\sum_{j=1}^{K_n-1}\{2(K_n-j)\}^2
=4\sum_{j=1}^{K_n-1}j^2
=\frac23(K_n-1)K_n(2K_n-1),
\]
where the square-sum formula follows by summing the successive differences of cubes. Consequently,
\[
\begin{aligned}
\mathcal A_{K_n}
&=2\pi K_n^2+
\pi\sum_{j=1}^{K_n-1}\{2(K_n-j)\}^2\\
&=\frac{2\pi}{3}K_n(2K_n^2+1)
\ge\frac{4\pi}{3}K_n^3,
\qquad
J_{K_n}(t)=\frac{\mathcal S_{K_n}(t)^4}{\mathcal A_{K_n}}.
\end{aligned}
\]
The same cosine expansion, integrated without squaring, yields
\[
\int_{-\pi}^{\pi}\mathcal S_{K_n}(t)^2\,dt=2\pi K_n.
\]
For \(|t|\le\pi\), the sine lower bound gives
\[
|t|\le\pi|\sin(t/2)|,
\qquad
t^2\le\pi^2\sin^2(t/2).
\]
For \(t\ne0\), it follows that
\[
\begin{aligned}
t^2\mathcal S_{K_n}(t)^4
&\le\pi^2\sin^2(K_nt/2)\mathcal S_{K_n}(t)^2\\
&\le\pi^2\mathcal S_{K_n}(t)^2.
\end{aligned}
\]
At \(t=0\), this inequality follows directly from its zero left side. Integrating and normalizing therefore gives
\[
\int_{-\pi}^{\pi}t^2\mathcal S_{K_n}(t)^4\,dt
\le2\pi^3K_n,
\qquad
\int_{-\pi}^{\pi}t^2J_{K_n}(t)\,dt
\le\frac{24}{K_n^2}\le\frac{64}{K_n^2},
\]
using the mass lower bound and \(\pi^2\le16\).
Finally, the nonnegative square \((K_n|t|-2)^2\) implies
\[
|t|\le\frac{K_n}{4}t^2+\frac1{K_n}.
\]
Since the kernel is nonnegative and integrates to one, we obtain
\[
\begin{aligned}
\int_{-\pi}^{\pi}|t|J_{K_n}(t)\,dt
&\le\frac{K_n}{4}
\int_{-\pi}^{\pi}t^2J_{K_n}(t)\,dt+\frac1{K_n}\\
&\le\frac7{K_n}\le\frac{32}{K_n}.
\end{aligned}
\]
% lean: Causalean.Mathlib.Analysis.JacksonApproximation.jrawMass_eq
% lean: Causalean.Mathlib.Analysis.JacksonApproximation.jackson_second_moment
% lean: Causalean.Mathlib.Analysis.JacksonApproximation.jackson_first_moment
For \(q_{ay,x}=b_{ay,x}+r_{ay,x}\cos\theta_{ay,x}\), define
\[
\theta_{ay,x}
=
\arccos\left(\frac{q_{ay,x}-b_{ay,x}}{r_{ay,x}}\right)
\quad\text{on }\mathcal G_x.
\]
Taylor's formula gives
\[
\begin{aligned}
r_{ay,x}
|\cos(\theta_{ay,x}+t)-\cos\theta_{ay,x}|
\le{}&
\sqrt{(q_{ay,x}-\ell_{ay,x})(u_{ay,x}-q_{ay,x})}\,|t|\\
&+\frac{r_{ay,x}}2t^2.
\end{aligned}
\]
Applying \cref{obj:lem:armwise-modulus} at \(q_x\) inside the tensor convolution therefore gives
\[
\begin{aligned}
|P_{x,K_n}(q_x)-F_\epsilon(q_x)|
\le64\sum_{a=0}^1
\left(1+\frac{p_x}{s_a(q_x)}\right)
\left\{
\frac{\sum_y\sqrt{(q_{ay,x}-\ell_{ay,x})(u_{ay,x}-q_{ay,x})}}{K_n}
+\frac{R_{a,x}}{K_n^2}
\right\}.
\end{aligned}
\]

The localized coordinate formulas imply
\[
h_{ay,x}
\le H_0(H_0+2)
\left(\sqrt{q_{ay,x}\delta_{\mathrm p}}+\delta_{\mathrm p}\right).
\]
To see this, the lower-endpoint inequality implies
\[
c_{ay,x}-H_0\left(\sqrt{c_{ay,x}\delta_{\mathrm p}}+\delta_{\mathrm p}\right)
\le q_{ay,x}.
\]
Solving the resulting quadratic inequality in \(\sqrt{c_{ay,x}}\) gives
\[
\sqrt{c_{ay,x}\delta_{\mathrm p}}
\le
\sqrt{q_{ay,x}\delta_{\mathrm p}}+(H_0+1)\delta_{\mathrm p},
\]
which yields the stated half-width bound.

There is also the boundary-sensitive estimate
\[
\sqrt{(q_{ay,x}-\ell_{ay,x})(u_{ay,x}-q_{ay,x})}
\le2H_0(H_0+2)\sqrt{q_{ay,x}\delta_{\mathrm p}}.
\]
For \(q_{ay,x}\le\delta_{\mathrm p}\), bound the left distance by \(q_{ay,x}\), the right distance by \(2h_{ay,x}\), and the half width by
\(2H_0(H_0+2)\delta_{\mathrm p}\). For \(q_{ay,x}>\delta_{\mathrm p}\), the product of the two distances is at most \(h_{ay,x}^2\), and the half-width bound applies. Summing coordinate radii also gives
\[
R_{a,x}
\le2H_0(H_0+2)
\left(\sqrt{s_a(q_x)\delta_{\mathrm p}}+\delta_{\mathrm p}\right).
\]

For \(d\ge2\), the mass bound in \cref{obj:lem:armwise-modulus} supplies
\[
\sum_{a=0}^1\frac{p_x}{s_a(q_x)}
\sum_{y=0}^1\sqrt{q_{ay,x}}
\le2\sqrt2\sqrt{\frac{p_x}{\epsilon}}.
\]
For \(d=1\), the unique cell has \(p_x=1\), and the same bound follows directly from the two arm masses. For each arm, Cauchy--Schwarz and \(s_a(q_x)\ge\epsilon p_x>0\) give
\[
\sum_{y=0}^1\sqrt{q_{ay,x}}\le\sqrt{2s_a(q_x)},
\qquad
\frac{p_x}{s_a(q_x)}\sum_{y=0}^1\sqrt{q_{ay,x}}
\le\frac{\sqrt2\,p_x}{\sqrt{s_a(q_x)}}
\le\sqrt{\frac{2p_x}{\epsilon}}.
\]
Adding the two arm inequalities proves the displayed mass bound also in this case.
% lean: CausalSmith.Stat.OptvalueVanishingoverlapRate.observed_armwise_sqrt_bound
Together with \(1+p_x/s_a(q_x)\le2p_x/s_a(q_x)\), this bound now yields
\[
\begin{aligned}
|P_{x,K_n}(q_x)-F_\epsilon(q_x)|
\le512H_0(H_0+2)
\left\{
\frac{\sqrt2}{K_n}\sqrt{\frac{p_xL}{m\epsilon}}
+
\frac1{K_n^2}
\left(\sqrt{\frac{p_xL}{m\epsilon}}+\frac L{m\epsilon}\right)
\right\}.
\end{aligned}
\]
Define
\[
\eta_x=
\sqrt{\frac{p_x}{m\epsilon L}}+\frac1{m\epsilon L},
\]
\[
C_{\mathrm{occ}}
=
512H_0(H_0+2)\sqrt2
\left(\frac1{\kappa/2}+\frac1{(\kappa/2)^2}\right).
\]
Using \(K_n\ge(\kappa/2)L\) and \(L\ge1\), we conclude that
\[
|P_{x,K_n}(q_x)-F_\epsilon(q_x)|
\le C_{\mathrm{occ}}\eta_x
\quad\text{on }\mathcal G_x,\quad p_x>0.
\]
Combining this estimate with the clipping bounds and integrating over the pilot gives
\[
\mathbb E_{\mathrm{aux}}
[(T_x-F_\epsilon(q_x))^2\mathbf1_{\mathcal G_x}]
\le
2C_{\mathrm{fac}}d^{1/16}
\mathbb E_{\mathrm{aux}}(S_x^{\mathrm{aw}})^2
+2C_{\mathrm{occ}}^2\eta_x^2,
\]
\[
\left|
\mathbb E_{\mathrm{aux}}
[(T_x-F_\epsilon(q_x))\mathbf1_{\mathcal G_x}]
\right|
\le
C_{\mathrm{occ}}\eta_x
+
C_{\mathrm{fac}}d^{-3/16}
\mathbb E_{\mathrm{aux}}S_x^{\mathrm{aw}}.
\]
The squared bound uses
\((v+w)^2\le2v^2+2w^2\); the mean bound uses the triangle inequality.
% lean: CausalSmith.Stat.OptvalueVanishingoverlapRate.observed_pilotJackson_mean_bias_normalized_le
% lean: CausalSmith.Stat.OptvalueVanishingoverlapRate.markedPoissonCellValue_occupied_good_moments_tuning_exists

\textit{6. Null cells and uniform successful-pilot bounds.}
If \(p_x=0\), every pilot and evaluation count in that cell is zero almost surely. Hence
\[
q_x=0,\qquad
b_{ay,x}=r_{ay,x}=\frac{H_0L}{2m},\qquad
\ell_{ay,x}=0,
\qquad
\mathbb P_{\mathrm{aux}}(\mathcal G_x)=1.
\]
At the lower corner, the factorial statistic equals the polynomial value:
\[
Z_x=P_{x,K_n}(0).
\]
This follows either by substituting zero counts in the factorial expansion or by its expectation identity, since the evaluation counts are constant. The polynomial oscillation bound places this value inside the clipping interval, so
\[
T_x=P_{x,K_n}(0).
\]
From \cref{obj:def:armwise-extension},
\[
0\le F_\epsilon(v)\le s(v),\qquad v\in\mathbb R_+^4.
\]
At the null corner, a displaced coordinate in the convolution is
\(r_{ay,x}(1-\cos t)\le r_{ay,x}t^2/2\). The second kernel moment thus gives
\[
|T_x-F_\epsilon(0)|
\le\frac{32}{K_n^2}\sum_{a,y}r_{ay,x}
=\frac{64H_0L}{mK_n^2}.
\]
Define
\[
C_{\mathrm{null}}=\frac{64H_0}{(\kappa/2)^2},
\qquad
C_{\mathrm{app}}=C_{\mathrm{occ}}+C_{\mathrm{null}}.
\]
Since \(K_n\ge(\kappa/2)L\) and \(\epsilon\le1\), the null-cell error is at most \(C_{\mathrm{null}}\eta_x\).

Therefore, for every cell, occupied or null, the successful-pilot estimates take the common form
\[
\mathbb E_{\mathrm{aux}}
[(T_x-F_\epsilon(q_x))^2\mathbf1_{\mathcal G_x}]
\le
C_{\mathrm v}d^{1/16}
\mathbb E_{\mathrm{aux}}(S_x^{\mathrm{aw}})^2
+C_{\mathrm e}\eta_x^2,
\]
\[
\left|
\mathbb E_{\mathrm{aux}}
[(T_x-F_\epsilon(q_x))\mathbf1_{\mathcal G_x}]
\right|
\le
C_{\mathrm b}
\left(\eta_x+d^{-3/16}\mathbb E_{\mathrm{aux}}S_x^{\mathrm{aw}}\right),
\]
where the universal constants are
\[
C_{\mathrm v}=2C_{\mathrm{fac}},\qquad
C_{\mathrm e}=2C_{\mathrm{app}}^2,\qquad
C_{\mathrm b}=C_{\mathrm{app}}+C_{\mathrm{fac}}.
\]
These enlargements use the nonnegativity of the clipping scales and of \(\eta_x\).
% lean: CausalSmith.Stat.OptvalueVanishingoverlapRate.markedPoissonCellValue_null_good_moments_le
% lean: CausalSmith.Stat.OptvalueVanishingoverlapRate.markedPoissonStatistic_active_rate_tuning_exists

\textit{7. Unconditional clipping-scale moments.}
For each pilot, the coordinate inequalities \(r_{ay,x}\le h_{ay,x}\) and \(c_{ay,x}\le b_{ay,x}\) imply
\[
R_{a,x}
\le2H_0\left(\sqrt{s_a(b_x)\delta_{\mathrm p}}+\delta_{\mathrm p}\right).
\]
Using \(D_a(b_x)\ge s_a(b_x)\) and
\(D_a(b_x)\ge\epsilon s(b_x)\), separately according as
\(s_a(b_x)\ge\epsilon s(b_x)\) or
\(s_a(b_x)<\epsilon s(b_x)\), gives
\[
\frac{s(b_x)}{D_a(b_x)}R_{a,x}
\le
2H_0\left(
\sqrt{\frac{s(b_x)\delta_{\mathrm p}}{\epsilon}}
+\frac{\delta_{\mathrm p}}{\epsilon}
\right).
\]
The unweighted radius satisfies this same bound, since
\(s_a(b_x)\le s(b_x)\) and \(\epsilon\le1\). Consequently,
\[
S_x^{\mathrm{aw}}
\le16H_0\left(
\sqrt{\frac{s(b_x)\delta_{\mathrm p}}{\epsilon}}
+\frac{\delta_{\mathrm p}}{\epsilon}
\right),
\]
\[
(S_x^{\mathrm{aw}})^2
\le512H_0^2\left(
\frac{s(b_x)\delta_{\mathrm p}}{\epsilon}
+\frac{\delta_{\mathrm p}^2}{\epsilon^2}
\right).
\]

The midpoint formulas also give
\[
b_{ay,x}\le(1+H_0)c_{ay,x}+2H_0\delta_{\mathrm p},
\qquad
s(b_x)\le(1+H_0)\sum_{a,y}c_{ay,x}+8H_0\delta_{\mathrm p}.
\]
For example, use \(b_{ay,x}\le c_{ay,x}+h_{ay,x}\) and
\(\sqrt{c_{ay,x}\delta_{\mathrm p}}\le c_{ay,x}+\delta_{\mathrm p}\).
Because \(\mathbb E_{\mathrm{aux}}c_{ay,x}=q_{ay,x}\), integration yields
\[
\mathbb E_{\mathrm{aux}}(S_x^{\mathrm{aw}})^2
\le512H_0^2
\left\{
(1+H_0)\frac{p_x\delta_{\mathrm p}}{\epsilon}
+8H_0\frac{\delta_{\mathrm p}^2}{\epsilon}
+\frac{\delta_{\mathrm p}^2}{\epsilon^2}
\right\}.
\]
Define, for every cell,
\[
\mathfrak a_x
=
\frac{p_xL}{m\epsilon}
+\frac{L^2}{m^2\epsilon^2},
\qquad
C_{\mathrm s}=512H_0^2(1+9H_0).
\]
Since \(\epsilon\le1\),
\[
\mathbb E_{\mathrm{aux}}(S_x^{\mathrm{aw}})^2
\le C_{\mathrm s}\mathfrak a_x.
\]
% lean: CausalSmith.Stat.OptvalueVanishingoverlapRate.markedPoissonClippingScale_integral_sq_le

\textit{8. Failed-pilot moments.}
Define the nonnegative coordinate scores by
\[
e_{ay,x}^{\mathrm p}
=
|c_{ay,x}-q_{ay,x}|+h_{ay,x},
\qquad
e_x^{\mathrm p}=\sum_{a,y}e_{ay,x}^{\mathrm p}.
\]
We need both a probability bound and moments retaining the cell mass:
\[
\mathbb P_{\mathrm{aux}}(\mathcal G_x^c)\le8e^{-40L},
\]
\[
\mathbb E_{\mathrm{aux}}(e_{ay,x}^{\mathrm p})^4
\le
\left(\frac{H_0}{1024}\right)^4
10^{20}
\left(\sqrt{q_{ay,x}\delta_{\mathrm p}}+\delta_{\mathrm p}\right)^4,
\]
\[
\mathbb E_{\mathrm{aux}}
[e_x^{\mathrm p}\mathbf1_{\mathcal G_x^c}]
\le
\frac{H_0}{1024}\,10^{160}e^{-20L}
\left(\sqrt{p_x\delta_{\mathrm p}}+\delta_{\mathrm p}\right).
\]

Here are the underlying calculations. For any
\(\lambda_\circ\ge0\) and
\(\mathsf N_\circ\sim\operatorname{Pois}(\lambda_\circ)\), define
\[
\mathfrak s_\circ=\sqrt{\lambda_\circ L}+L,\qquad
\mathfrak e_\circ
=
|\mathsf N_\circ-\lambda_\circ|
+1024\left(\sqrt{\mathsf N_\circ L}+L\right),
\qquad
\vartheta_\circ=\frac1{2\mathfrak s_\circ}.
\]
The centered Poisson generating function is
\[
\mathbb E e^{t(\mathsf N_\circ-\lambda_\circ)}
=
\exp\{\lambda_\circ(e^t-1-t)\}.
\]
If \(\lambda_\circ=0\), then \(\mathsf N_\circ=0\) almost surely, and both tail events below have probability zero. For \(\lambda_\circ>0\), define the deviation threshold and exponential tilt at an arbitrary nonnegative level by
\[
z_{\mathrm{tail}}\in[0,\infty),\qquad
\mathfrak x_{\mathrm{up}}
=\sqrt{2\lambda_\circ z_{\mathrm{tail}}}+2z_{\mathrm{tail}},
\qquad
\vartheta_{\mathrm{up}}
=\log\left(1+\frac{\mathfrak x_{\mathrm{up}}}{\lambda_\circ}\right)\ge0.
\]
Exponential Markov inequality and the generating function give
\[
\begin{aligned}
\Pr\{\mathsf N_\circ-\lambda_\circ>\mathfrak x_{\mathrm{up}}\}
&\le
\exp\left\{
-\vartheta_{\mathrm{up}}\mathfrak x_{\mathrm{up}}
+\lambda_\circ(e^{\vartheta_{\mathrm{up}}}-1-\vartheta_{\mathrm{up}})
\right\}\\
&=
\exp\left\{
-\bigl[(\lambda_\circ+\mathfrak x_{\mathrm{up}})
\vartheta_{\mathrm{up}}-\mathfrak x_{\mathrm{up}}\bigr]
\right\}.
\end{aligned}
\]
The logarithmic lower bound follows from
\[
\frac{d}{d\xi_{\mathrm{log}}}
\left\{\log(1+\xi_{\mathrm{log}})
-\frac{2\xi_{\mathrm{log}}}{\xi_{\mathrm{log}}+2}\right\}
=
\frac{\xi_{\mathrm{log}}^2}
{(1+\xi_{\mathrm{log}})(\xi_{\mathrm{log}}+2)^2}
\ge0,
\qquad \xi_{\mathrm{log}}\in[0,\infty),
\]
since the difference is zero at \(\xi_{\mathrm{log}}=0\). Applying it at \(\mathfrak x_{\mathrm{up}}/\lambda_\circ\) gives
\[
\vartheta_{\mathrm{up}}
\ge\frac{2\mathfrak x_{\mathrm{up}}}{2\lambda_\circ+\mathfrak x_{\mathrm{up}}},
\qquad
(\lambda_\circ+\mathfrak x_{\mathrm{up}})\vartheta_{\mathrm{up}}
-\mathfrak x_{\mathrm{up}}
\ge
\frac{\mathfrak x_{\mathrm{up}}^2}
{2\lambda_\circ+\mathfrak x_{\mathrm{up}}}
\ge z_{\mathrm{tail}}.
\]
The last inequality follows by expanding the threshold:
\[
\mathfrak x_{\mathrm{up}}^2
-z_{\mathrm{tail}}(2\lambda_\circ+\mathfrak x_{\mathrm{up}})
=
3z_{\mathrm{tail}}\sqrt{2\lambda_\circ z_{\mathrm{tail}}}
+2z_{\mathrm{tail}}^2\ge0.
\]
Thus the upper-tail probability is at most \(e^{-z_{\mathrm{tail}}}\). Taking \(z_{\mathrm{tail}}=40L\), and applying the lower-tail exponential Markov inequality with
\(e^{-t}-1+t\le t^2/2\), \(t\ge0\), gives
\[
\Pr\!\left\{
\mathsf N_\circ-\lambda_\circ>
\sqrt{80\lambda_\circ L}+80L
\right\}\le e^{-40L},
\qquad
\Pr\!\left\{
\lambda_\circ-\mathsf N_\circ>
\sqrt{80\lambda_\circ L}
\right\}\le e^{-40L}.
\]
The event
\[
\left\{
|\mathsf N_\circ-\lambda_\circ|
>\frac{1024}{4}\left(\sqrt{\mathsf N_\circ L}+L\right)
\right\}
\]
is contained in the union of these two tail events: on their complement, substitute
\(\lambda_\circ\le\mathsf N_\circ+|\mathsf N_\circ-\lambda_\circ|\) and square the nonnegative terms. Since \(H_0\ge1024\), failure of a pilot coordinate implies the corresponding event above. A union bound over the four coordinates proves the probability estimate.

For the moments, square-root subadditivity and
\(\sqrt{|\mathsf N_\circ-\lambda_\circ|L}
\le\{|\mathsf N_\circ-\lambda_\circ|+L\}/2\) imply
\[
\mathfrak e_\circ
\le1536\left(\mathfrak s_\circ+
|\mathsf N_\circ-\lambda_\circ|\right).
\]
Also \(\mathfrak s_\circ\ge1\),
\(\lambda_\circ\le\mathfrak s_\circ^2\), and
\(\lambda_\circ\vartheta_\circ^2\le1/4\). Using
\(|e^t-1-t|\le t^2\) for \(|t|\le1\) in the generating function therefore gives
\[
\mathbb E\exp\left(\frac{\mathfrak e_\circ}{4096\mathfrak s_\circ}\right)
\le
e^{1/2}\left(
\mathbb E e^{\vartheta_\circ(\mathsf N_\circ-\lambda_\circ)}
+
\mathbb E e^{-\vartheta_\circ(\mathsf N_\circ-\lambda_\circ)}
\right)
\le2e^{3/4}\le8.
\]
Thus, for \(j=2,4\), the inequality \(v^j\le j!e^v\), \(v\ge0\), yields
\[
\mathbb E\mathfrak e_\circ^j
\le8j!\,4096^j\mathfrak s_\circ^j
\le10^{4j+4}\mathfrak s_\circ^j.
\]
Substituting \(\lambda_\circ=mq_{ay,x}\) gives the displayed fourth-moment bound, because
\[
e_{ay,x}^{\mathrm p}
\le\frac{H_0}{1024m}\mathfrak e_\circ.
\]
We next localize the scalar moments before summing over failing coordinates. For the scalar Poisson count above, define
\[
\Delta_\circ=|\mathsf N_\circ-\lambda_\circ|,
\qquad
\mathcal B_\circ=
\left\{
\Delta_\circ>256\left(\sqrt{\mathsf N_\circ L}+L\right)
\right\}.
\]
Applying the exponential Markov calculations at level \(80L\) gives
\[
\Pr\!\left\{
\mathsf N_\circ-\lambda_\circ>
\sqrt{160\lambda_\circ L}+160L
\right\}\le e^{-80L},
\qquad
\Pr\!\left\{
\lambda_\circ-\mathsf N_\circ>
\sqrt{160\lambda_\circ L}
\right\}\le e^{-80L}.
\]
To compare \(\mathcal B_\circ\) with these events, first suppose \(\mathsf N_\circ\ge\lambda_\circ\). On \(\mathcal B_\circ\),
\[
\sqrt{160\lambda_\circ L}+160L
\le16\sqrt{\mathsf N_\circ L}+160L
\le256\left(\sqrt{\mathsf N_\circ L}+L\right)
<\Delta_\circ.
\]
If \(\mathsf N_\circ<\lambda_\circ\), then \(\lambda_\circ=\mathsf N_\circ+\Delta_\circ\). The defining inequality of \(\mathcal B_\circ\), and monotonicity of the quadratic above that threshold, give
\[
\begin{aligned}
\Delta_\circ^2-160L\Delta_\circ-160\mathsf N_\circ L
>{}&65376\mathsf N_\circ L
+90112L\sqrt{\mathsf N_\circ L}+24576L^2>0.
\end{aligned}
\]
Thus \(\Delta_\circ>\sqrt{160\lambda_\circ L}\). The two cases prove
\[
\Pr(\mathcal B_\circ)\le2e^{-80L}.
\]
The exponential score bound above also yields, for every integer \(1\le j_{\mathrm{mom}}\le8\),
\[
\mathbb E\mathfrak e_\circ^{j_{\mathrm{mom}}}
\le8j_{\mathrm{mom}}!\,4096^{j_{\mathrm{mom}}}
\mathfrak s_\circ^{j_{\mathrm{mom}}}
\le10^{4j_{\mathrm{mom}}+4}\mathfrak s_\circ^{j_{\mathrm{mom}}}.
\]
For \(j_{\mathrm{mom}}\in\{1,2,4\}\), define the positive balancing coefficient by
\[
\upsilon_{\circ,j_{\mathrm{mom}}}
=\frac{e^{-40L}}{\mathfrak s_\circ^{j_{\mathrm{mom}}}}.
\]
The nonnegative square
\((\upsilon_{\circ,j_{\mathrm{mom}}}\mathfrak e_\circ^{j_{\mathrm{mom}}}-1)^2\), on \(\mathcal B_\circ\), proves the pointwise inequality
\[
\mathfrak e_\circ^{j_{\mathrm{mom}}}\mathbf1_{\mathcal B_\circ}
\le\frac12\left(
\upsilon_{\circ,j_{\mathrm{mom}}}\mathfrak e_\circ^{2j_{\mathrm{mom}}}
+\upsilon_{\circ,j_{\mathrm{mom}}}^{-1}\mathbf1_{\mathcal B_\circ}
\right).
\]
Outside \(\mathcal B_\circ\), its left side is zero and its right side is nonnegative. Integration therefore gives the localized scalar bound
\[
\begin{aligned}
\mathbb E\!\left[
\mathfrak e_\circ^{j_{\mathrm{mom}}}\mathbf1_{\mathcal B_\circ}
\right]
&\le\frac{10^{8j_{\mathrm{mom}}+4}+2}{2}
 e^{-40L}\mathfrak s_\circ^{j_{\mathrm{mom}}}\\
&\le10^{8j_{\mathrm{mom}}+8}
 e^{-40L}\mathfrak s_\circ^{j_{\mathrm{mom}}}.
\end{aligned}
\]
% lean: Causalean.Stat.Concentration.PoissonSelfNormalized.poisson_weighted_bad_moment

For each \((a,y)\in\{0,1\}^2\) and each cell \(x\), define the raw score, raw local scale, and raw failure event by
\[
\mathfrak e_{ay,x}
=|N'_{ay,x}-mq_{ay,x}|
+1024\left(\sqrt{N'_{ay,x}L}+L\right),
\qquad
\mathfrak s_{ay,x}=\sqrt{mq_{ay,x}L}+L,
\]
\[
\mathcal B_{ay,x}
=\left\{
|N'_{ay,x}-mq_{ay,x}|>
256\left(\sqrt{N'_{ay,x}L}+L\right)
\right\},
\qquad
\mathcal B_x^{\cup}=\bigcup_{a,y}\mathcal B_{ay,x}.
\]
The interval-failure criterion and \(H_0\ge1024\) imply
\[
\mathcal G_x^c\subseteq\mathcal B_x^{\cup},
\qquad
e_{ay,x}^{\mathrm p}
\le\frac{H_0}{1024m}\mathfrak e_{ay,x}.
\]
For \(j_{\mathrm{mom}}\in\{1,2,4\}\), consider the contribution of one coordinate's score on another coordinate's failure event. If \((a,y)=(a',y')\), the localized scalar bound gives
\[
\mathbb E_{\mathrm{aux}}\!\left[
\mathfrak e_{ay,x}^{j_{\mathrm{mom}}}\mathbf1_{\mathcal B_{ay,x}}
\right]
\le10^{8j_{\mathrm{mom}}+8}e^{-40L}
\mathfrak s_{ay,x}^{j_{\mathrm{mom}}}.
\]
If \((a,y)\ne(a',y')\), independence of the pilot coordinates factors the expectation:
\[
\begin{aligned}
\mathbb E_{\mathrm{aux}}\!\left[
\mathfrak e_{ay,x}^{j_{\mathrm{mom}}}\mathbf1_{\mathcal B_{a'y',x}}
\right]
&=\mathbb E_{\mathrm{aux}}\mathfrak e_{ay,x}^{j_{\mathrm{mom}}}
\Pr(\mathcal B_{a'y',x})\\
&\le2\cdot10^{4j_{\mathrm{mom}}+4}e^{-40L}
\mathfrak s_{ay,x}^{j_{\mathrm{mom}}},
\end{aligned}
\]
where the scalar probability bound was enlarged from \(2e^{-80L}\) to \(2e^{-40L}\). Hence both cases satisfy
\[
\mathbb E_{\mathrm{aux}}\!\left[
\mathfrak e_{ay,x}^{j_{\mathrm{mom}}}\mathbf1_{\mathcal B_{a'y',x}}
\right]
\le
\left(10^{8j_{\mathrm{mom}}+8}+2\cdot10^{4j_{\mathrm{mom}}+4}\right)
e^{-40L}\mathfrak s_{ay,x}^{j_{\mathrm{mom}}}.
\]
The union indicator and the finite-sum power inequality give the pointwise bound
\[
\left(\sum_{a,y}\mathfrak e_{ay,x}\right)^{j_{\mathrm{mom}}}
\mathbf1_{\mathcal B_x^{\cup}}
\le4^{j_{\mathrm{mom}}-1}
\sum_{a',y'}\sum_{a,y}
\mathfrak e_{ay,x}^{j_{\mathrm{mom}}}
\mathbf1_{\mathcal B_{a'y',x}}.
\]
To sum the local scales, Cauchy--Schwarz and \(\sum_{a,y}q_{ay,x}=p_x\) give
\[
\sum_{a,y}\mathfrak s_{ay,x}
\le2\sqrt{mp_xL}+4L
\le4\left(\sqrt{mp_xL}+L\right),
\qquad
\sum_{a,y}\mathfrak s_{ay,x}^{j_{\mathrm{mom}}}
\le4\left(\sum_{a,y}\mathfrak s_{ay,x}\right)^{j_{\mathrm{mom}}}.
\]
Integrating the pointwise bound and inserting the two coordinate cases therefore yields
\[
\begin{aligned}
\mathbb E_{\mathrm{aux}}\!\left[
\left(\sum_{a,y}\mathfrak e_{ay,x}\right)^{j_{\mathrm{mom}}}
\mathbf1_{\mathcal B_x^{\cup}}
\right]
\le{}&4^{2j_{\mathrm{mom}}+1}
\left(10^{8j_{\mathrm{mom}}+8}+2\cdot10^{4j_{\mathrm{mom}}+4}\right)
e^{-40L}\left(\sqrt{mp_xL}+L\right)^{j_{\mathrm{mom}}}\\
\le{}&10^{80(j_{\mathrm{mom}}+1)}e^{-20L}
\left(\sqrt{mp_xL}+L\right)^{j_{\mathrm{mom}}}.
\end{aligned}
\]
The last enlargement uses, for the three moment orders,
\[
4^{2j_{\mathrm{mom}}+1}
\left(10^{8j_{\mathrm{mom}}+8}+2\cdot10^{4j_{\mathrm{mom}}+4}\right)
\le10^{80(j_{\mathrm{mom}}+1)},
\qquad e^{-40L}\le e^{-20L}.
\]
Finally, the raw-score comparison and
\[
\frac{\sqrt{mp_xL}+L}{m}
=\sqrt{p_x\delta_{\mathrm p}}+\delta_{\mathrm p}
\]
prove the normalized aggregate bound
\[
\mathbb E_{\mathrm{aux}}\!\left[
(e_x^{\mathrm p})^{j_{\mathrm{mom}}}\mathbf1_{\mathcal G_x^c}
\right]
\le
\left(\frac{H_0}{1024}\right)^{j_{\mathrm{mom}}}
10^{80(j_{\mathrm{mom}}+1)}e^{-20L}
\left(\sqrt{p_x\delta_{\mathrm p}}+\delta_{\mathrm p}\right)^{j_{\mathrm{mom}}}.
\]
Taking \(j_{\mathrm{mom}}=1\) gives the stated aggregate first-moment estimate with numerical constant \(10^{160}\).
% lean: CausalSmith.Stat.OptvalueVanishingoverlapRate.pilotRectangle_bad_probability_le
% lean: CausalSmith.Stat.OptvalueVanishingoverlapRate.pilotRectangle_bad_aggregate_moment_le

The midpoint error satisfies
\[
|b_{ay,x}-q_{ay,x}|\le e_{ay,x}^{\mathrm p},
\qquad
s(b_x)\le p_x+e_x^{\mathrm p}.
\]
Using the clipping-scale square bound from the preceding step, we obtain
\[
\begin{aligned}
\mathbb E_{\mathrm{aux}}
[(S_x^{\mathrm{aw}})^2\mathbf1_{\mathcal G_x^c}]
\le512H_0^2\bigg\{
8e^{-40L}\mathfrak a_x
+\frac{\delta_{\mathrm p}}{\epsilon}
\frac{H_0}{1024}10^{160}e^{-20L}
\left(\sqrt{p_x\delta_{\mathrm p}}+\delta_{\mathrm p}\right)
\bigg\}.
\end{aligned}
\]
Since
\[
\frac{\delta_{\mathrm p}}{\epsilon}
\left(\sqrt{p_x\delta_{\mathrm p}}+\delta_{\mathrm p}\right)
\le2\mathfrak a_x,
\]
where \(\sqrt{p_x\delta_{\mathrm p}}\le p_x+\delta_{\mathrm p}\) and \(\epsilon\le1\), this becomes
\[
\mathbb E_{\mathrm{aux}}
[(S_x^{\mathrm{aw}})^2\mathbf1_{\mathcal G_x^c}]
\le
C_{\mathrm{fail,s}}e^{-20L}\mathfrak a_x,
\]
with
\[
C_{\mathrm{fail,s}}
=
512H_0^2\left(8+2\frac{H_0}{1024}10^{160}\right).
\]

For an occupied cell, \cref{obj:lem:armwise-modulus} gives
\[
|F_\epsilon(b_x)-F_\epsilon(q_x)|
\le
4\sum_{a,y}\frac{p_x}{s_a(q_x)}e_{ay,x}^{\mathrm p}.
\]
The fourth-moment and probability estimates imply
\[
\begin{aligned}
\mathbb E_{\mathrm{aux}}
[(e_{ay,x}^{\mathrm p})^2\mathbf1_{\mathcal G_x^c}]
&\le
\left\{\mathbb E_{\mathrm{aux}}(e_{ay,x}^{\mathrm p})^4\right\}^{1/2}
\Pr(\mathcal G_x^c)^{1/2}\\
&\le
\sqrt{8\cdot10^{20}}
\left(\frac{H_0}{1024}\right)^2
e^{-20L}
\left(\sqrt{q_{ay,x}\delta_{\mathrm p}}+\delta_{\mathrm p}\right)^2.
\end{aligned}
\]
For each arm, the overlap floor gives
\[
\left(\frac{p_x}{s_a(q_x)}\right)^2
\sum_y
\left(\sqrt{q_{ay,x}\delta_{\mathrm p}}+\delta_{\mathrm p}\right)^2
\le8\mathfrak a_x.
\]
Indeed, the left side is at most
\[
2\frac{p_x^2\delta_{\mathrm p}}{s_a(q_x)}
+
4\frac{p_x^2\delta_{\mathrm p}^2}{s_a(q_x)^2}
\le
2\frac{p_x\delta_{\mathrm p}}{\epsilon}
+
4\frac{\delta_{\mathrm p}^2}{\epsilon^2}.
\]
Squaring the four-coordinate sum consequently yields
\[
\mathbb E_{\mathrm{aux}}
[(F_\epsilon(b_x)-F_\epsilon(q_x))^2\mathbf1_{\mathcal G_x^c}]
\le C_{\mathrm{fail,f}}e^{-20L}\mathfrak a_x,
\]
where
\[
C_{\mathrm{fail,f}}
=
1024\sqrt{8\cdot10^{20}}
\left(\frac{H_0}{1024}\right)^2.
\]

Clipping supplies the pathwise bound
\[
|T_x-F_\epsilon(q_x)|
\le
d^{1/4}S_x^{\mathrm{aw}}
+
|F_\epsilon(b_x)-F_\epsilon(q_x)|.
\]
Hence, with
\[
C_{\mathrm d}=2(C_{\mathrm{fail,s}}+C_{\mathrm{fail,f}}),
\]
we have
\[
\mathbb E_{\mathrm{aux}}
[(T_x-F_\epsilon(q_x))^2\mathbf1_{\mathcal G_x^c}]
\le C_{\mathrm d}d^{-4}\mathfrak a_x.
\]
Here
\[
d^{1/2}e^{-20L}\le d^{-4},
\qquad e^{-20L}\le d^{-4}
\]
for every \(d\ge1\). For null cells, \(\mathcal G_x^c\) has probability zero, so the same bound follows directly.
% lean: CausalSmith.Stat.OptvalueVanishingoverlapRate.markedPoissonCellValue_bad_sq_le

\textit{9. Independent-cell risk assembly.}
Define, on the whole auxiliary sample space,
\[
E_x^{\mathrm{good}}
=
(T_x-F_\epsilon(q_x))\mathbf1_{\mathcal G_x},
\qquad
E_x^{\mathrm{bad}}
=
(T_x-F_\epsilon(q_x))\mathbf1_{\mathcal G_x^c}.
\]
Then
\[
T_x-F_\epsilon(q_x)=E_x^{\mathrm{good}}+E_x^{\mathrm{bad}},
\qquad
E_x^{\mathrm{good}}E_x^{\mathrm{bad}}=0.
\]
All these variables have finite second moments by the preceding estimates. Independence of the complete cell statistics gives
\[
\mathbb E_{\mathrm{aux}}
\left(\sum_xT_x-\Psi(\mathbb P)\right)^2
=
\sum_x\operatorname{Var}_{\mathrm{aux}}(T_x)
+
\left\{\sum_x
\mathbb E_{\mathrm{aux}}(T_x-F_\epsilon(q_x))\right\}^2.
\]
The disjoint error decomposition implies
\[
\operatorname{Var}_{\mathrm{aux}}(T_x)
\le
\mathbb E_{\mathrm{aux}}(E_x^{\mathrm{good}})^2
+
\mathbb E_{\mathrm{aux}}(E_x^{\mathrm{bad}})^2.
\]
Applying \((v+w)^2\le2v^2+2w^2\) to the total bias, and Cauchy--Schwarz to its failed-pilot part, therefore gives
\[
\begin{aligned}
\mathbb E_{\mathrm{aux}}
\left(\sum_xT_x-\Psi(\mathbb P)\right)^2
\le{}&
\sum_x\mathbb E_{\mathrm{aux}}(E_x^{\mathrm{good}})^2
+
2\left\{\sum_x\mathbb E_{\mathrm{aux}}E_x^{\mathrm{good}}\right\}^2\\
&+
(1+2d)\sum_x\mathbb E_{\mathrm{aux}}(E_x^{\mathrm{bad}})^2.
\end{aligned}
\]
% lean: CausalSmith.Stat.OptvalueVanishingoverlapRate.independentCell_good_bad_sqRisk_le

\textit{10. Absorption of the pilot scales.}
Assume the active-branch condition \(d/L\le n\epsilon\), and define
\[
\mathfrak r_{\mathrm{aux}}=\frac{d}{m\epsilon L},
\qquad
\mathfrak m_{\mathrm{sum}}
=\sum_x\mathfrak a_x
=\frac L{m\epsilon}+\frac{dL^2}{m^2\epsilon^2}.
\]
Then
\[
0<\mathfrak r_{\mathrm{aux}}\le8,
\qquad
\mathfrak m_{\mathrm{sum}}
\le\frac{L^2+8L^4}{m\epsilon L}.
\]
The second inequality follows by substituting
\(d\le8m\epsilon L\) into its second summand.

For a positive real \(\delta_{\mathrm{dec}}\), define the numerical envelope
\[
\mathcal M(\delta_{\mathrm{dec}})
=
\left(\frac{e^{\delta_{\mathrm{dec}}/2}}{\delta_{\mathrm{dec}}/2}\right)^2
+
8\left(\frac{e^{\delta_{\mathrm{dec}}/4}}{\delta_{\mathrm{dec}}/4}\right)^4.
\]
The elementary inequality \(\log z\le z^u/u\), \(z\ge1\), \(u>0\), applied with \(z=ed\), gives
\[
(L^2+8L^4)d^{-\delta_{\mathrm{dec}}}
\le\mathcal M(\delta_{\mathrm{dec}}).
\]
Set
\[
C_{\mathrm{var}}=\mathcal M(15/16),\qquad
C_{\mathrm{bias}}=\mathcal M(3/8),\qquad
C_{\mathrm{bad,sum}}=\mathcal M(1).
\]
Thus
\[
d^{1/16}\sum_x\mathbb E_{\mathrm{aux}}(S_x^{\mathrm{aw}})^2
\le C_{\mathrm s}C_{\mathrm{var}}\mathfrak r_{\mathrm{aux}}.
\]

The probability-mass identity and Cauchy--Schwarz give
\[
\sum_x\eta_x\le\sqrt{\mathfrak r_{\mathrm{aux}}}+\mathfrak r_{\mathrm{aux}},
\qquad
\sum_x\eta_x^2\le\left(\sum_x\eta_x\right)^2
\le18\mathfrak r_{\mathrm{aux}}.
\]
Also,
\[
\left(\sum_x\mathbb E_{\mathrm{aux}}S_x^{\mathrm{aw}}\right)^2
\le d\sum_x\mathbb E_{\mathrm{aux}}(S_x^{\mathrm{aw}})^2,
\]
and hence
\[
d^{-3/8}
\left(\sum_x\mathbb E_{\mathrm{aux}}S_x^{\mathrm{aw}}\right)^2
\le C_{\mathrm s}C_{\mathrm{bias}}\mathfrak r_{\mathrm{aux}}.
\]
The successful-pilot estimates now imply
\[
\sum_x\mathbb E_{\mathrm{aux}}(E_x^{\mathrm{good}})^2
\le
(C_{\mathrm v}C_{\mathrm s}C_{\mathrm{var}}+18C_{\mathrm e})
\mathfrak r_{\mathrm{aux}},
\]
\[
\left(\sum_x\mathbb E_{\mathrm{aux}}E_x^{\mathrm{good}}\right)^2
\le
2C_{\mathrm b}^2(18+C_{\mathrm s}C_{\mathrm{bias}})
\mathfrak r_{\mathrm{aux}}.
\]
For the failed-pilot term, \(d\,d^{-4}\le1\) gives
\[
d\sum_x\mathbb E_{\mathrm{aux}}(E_x^{\mathrm{bad}})^2
\le C_{\mathrm d}\mathfrak m_{\mathrm{sum}}
\le C_{\mathrm d}C_{\mathrm{bad,sum}}\mathfrak r_{\mathrm{aux}},
\]
so, since \(d\ge1\),
\[
(1+2d)\sum_x\mathbb E_{\mathrm{aux}}(E_x^{\mathrm{bad}})^2
\le3C_{\mathrm d}C_{\mathrm{bad,sum}}\mathfrak r_{\mathrm{aux}}.
\]

Define
\[
\begin{aligned}
C_{\mathrm P}={}&
C_{\mathrm v}C_{\mathrm s}C_{\mathrm{var}}
+4C_{\mathrm b}^2(18+C_{\mathrm s}C_{\mathrm{bias}})
+3C_{\mathrm d}C_{\mathrm{bad,sum}}
+18C_{\mathrm e}+1.
\end{aligned}
\]
This is a positive finite universal constant. The independent-cell assembly yields
\[
\mathbb E_{\mathrm{aux}}
\left(\sum_xT_x-\Psi(\mathbb P)\right)^2
\le C_{\mathrm P}\frac{d}{m\epsilon L}.
\]
Finally, define the projected auxiliary statistic by
\[
V_{\mathrm{aux}}=\Pi_{[0,1]}\left(\sum_xT_x\right).
\]
Since \(\Psi(\mathbb P)\in[0,1]\), projection contracts its squared error, giving
\[
\mathbb E_{\mathrm{aux}}
(V_{\mathrm{aux}}-\Psi(\mathbb P))^2
\le C_{\mathrm P}\frac{d}{m\epsilon L}.
\]
% lean: CausalSmith.Stat.OptvalueVanishingoverlapRate.activeBranch_good_bad_approximation_sqRisk_rate_le
% lean: CausalSmith.Stat.OptvalueVanishingoverlapRate.markedPoissonStatistic_goodPilotEstimates_rate_le
% lean: CausalSmith.Stat.OptvalueVanishingoverlapRate.markedPoissonStatistic_active_rate_tuning_exists

\textit{11. Fixed-sample transfer.}
Set
\[
D_0=2,\qquad
C_\star=3+8\log(2e)+8C_{\mathrm P}.
\]
Then \(D_0\ge2\) and \(0<C_\star<\infty\). For the admissible theorem parameters \(n\ge1,d\ge2\), the finite sample alphabet makes the estimator in \cref{obj:def:armwise-estimator} measurable. Its auxiliary variables are integrated by the specified finite sum, so it is a deterministic function of the observations.

Fix an admissible observed law and its sampling law from \cref{obj:ass:iid-sampling}. First consider \(d<D_0\). With the chosen cutoff this branch has an empty parameter range, because \(d\ge2=D_0\).

Next suppose \(d\ge D_0\) and \(n\epsilon<d/L\). The estimator equals \(1/2\), and
\[
\mathbb E_{\mathbb P}
\left(\widehat V_{n,d,\epsilon}^{\mathrm{AF}}-\Psi(\mathbb P)\right)^2
\le\frac14,
\qquad
\min\left\{1,\frac{d}{n\epsilon L}\right\}=1.
\]
This proves the desired bound on that branch.

On the active branch, put the auxiliary count and marks alongside the fixed observed sample. Conditional averaging of squared error gives
\[
\begin{aligned}
\mathbb E_{\mathbb P}
\left(\widehat V_{n,d,\epsilon}^{\mathrm{AF}}-\Psi(\mathbb P)\right)^2
\le{}&
\mathbb E_{\mathrm{aux}}
\left[
(V_{\mathrm{aux}}-\Psi(\mathbb P))^2
\mathbf1_{\{\mathsf M\le n\}}
\right]\\
&+\Pr(\mathsf M>n)\Psi(\mathbb P)^2.
\end{aligned}
\]
Indeed, the auxiliary mixture has value zero on overflow. For each \(k\le n\), the first \(k\) observations have law \(\mathbb P^{\otimes k}\); consequently the integrated nonoverflow loss is exactly the corresponding part of the unrestricted finite-Poisson experiment.

The quarter-mean Poisson cap satisfies
\[
\Pr(\mathsf M>n)\le e^{-n/2}.
\]
For example, exponential Markov inequality at \(t=\log4\) gives
\[
\Pr(\mathsf M>n)
\le e^{-n(\log4-3/4)}
\le e^{-n/2}.
\]
Also, \(e^{n/2}\ge n/2\), \(L\le d\), and \(\epsilon\le1/2\) imply
\[
e^{-n/2}\le\frac2n
\le2\frac{d}{n\epsilon L}.
\]
Since \(\Psi(\mathbb P)^2\le1\), the auxiliary risk bound and \(m=n/8\) therefore give
\[
\mathbb E_{\mathbb P}
\left(\widehat V_{n,d,\epsilon}^{\mathrm{AF}}-\Psi(\mathbb P)\right)^2
\le
(8C_{\mathrm P}+2)\frac{d}{n\epsilon L}.
\]
The estimator and target both lie in \([0,1]\), so the risk is also at most one. Splitting according as \(d/(n\epsilon L)\ge1\) or \(d/(n\epsilon L)\le1\), we obtain
\[
\mathbb E_{\mathbb P}
\left(\widehat V_{n,d,\epsilon}^{\mathrm{AF}}-\Psi(\mathbb P)\right)^2
\le
(8C_{\mathrm P}+2)
\min\left\{1,\frac{d}{n\epsilon L}\right\}
\le
C_\star\min\left\{1,\frac{d}{n\epsilon L}\right\}.
\]
% lean: CausalSmith.Stat.OptvalueVanishingoverlapRate.armwiseEstimator_active_le_markedPoisson_sqRisk
% lean: CausalSmith.Stat.OptvalueVanishingoverlapRate.armwiseEstimator_allBranches_rate_le

\textit{12. Uniformity over sampling laws.}
The preceding estimates hold for every admissible observed law and every family of sampling laws satisfying \cref{obj:ass:iid-sampling}. If the admissible class is empty, its worst-risk value, with the empty-class convention zero, satisfies the bound because the right side is nonnegative. Otherwise, taking the supremum of the pointwise bound gives
\[
\sup_{\mathbb P\in\mathcal D_{d,\epsilon}^{\mathrm{obs}}}
\mathbb E_{\mathbb P}
\left[
\left\{
\widehat V_{n,d,\epsilon}^{\mathrm{AF}}-\Psi(\mathbb P)
\right\}^{2}
\right]
\le
C_\star\min\left\{1,\frac{d}{n\epsilon L_d}\right\}.
\]
Together with measurability and determinism, this proves the claim.
% lean: CausalSmith.Stat.OptvalueVanishingoverlapRate.armwiseEstimator_allBranches_worstRisk_le
% lean: CausalSmith.Stat.OptvalueVanishingoverlapRate.observable_upper
\end{proof}

\begin{proof}[Proof of \cref{obj:thm:weighted-separation-and-frontier}]
\begin{enumerate}
\item
\textbf{Approximation and constrained priors.}\quad
We first establish the packet lower bound. Throughout the approximation argument, assume
\[
K\ge2,\qquad 2\le M\le K^2,\qquad 0\le\epsilon\le\tfrac12,
\]
and write
\[
\mathcal E=\mathcal E_{K,M,\epsilon}^{w},
\qquad
\Delta=\Delta_{K,M,\epsilon},
\qquad
z_M(\omega)=\frac M2(1+\cos\omega),
\qquad
\vartheta_M=\arccos(2/M-1).
\]
Thus \(z_M(\vartheta_M)=1\). For integers \(N\ge2\), use
\(\mathcal J_N,k_N,G_N\) from
\cref{obj:lem:jackson-packet-localization}, and define the Lebesgue-normalized kernel
\[
J_N^{\mathrm J}(\omega)=\frac{\mathcal J_N(\omega)}{2\pi},
\qquad \omega\in\mathbb R.
\]
We will use the following two kernel moments:
\[
\int_{-\pi}^{\pi}\omega^2J_N^{\mathrm J}(\omega)\,d\omega
\le\frac{64}{N^2},
\qquad
\int_{-\pi}^{\pi}|\omega|J_N^{\mathrm J}(\omega)\,d\omega
\le\frac{32}{N}.
\]
Here is a direct derivation. Set, with removable values filled continuously,
\[
D_N(\omega)=\frac{\sin(N\omega/2)}{\sin(\omega/2)},
\qquad
\mathcal Z_N^{\mathrm J}
=\int_{-\pi}^{\pi}D_N(\omega)^4\,d\omega.
\]
The finite geometric-sum identity and orthogonality give
\[
D_N(\omega)^2
=N+2\sum_{\ell=1}^{N-1}(N-\ell)\cos(\ell\omega),
\]
\[
\int_{-\pi}^{\pi}D_N(\omega)^2\,d\omega=2\pi N,
\qquad
\mathcal Z_N^{\mathrm J}
=\frac{2\pi}{3}N(2N^2+1)
\ge\frac{4\pi}{3}N^3,
\qquad
J_N^{\mathrm J}=\frac{D_N^4}{\mathcal Z_N^{\mathrm J}}.
\]
For \(|\omega|\le\pi\), the inequality
\(|\sin(\omega/2)|\ge|\omega|/\pi\) implies
\[
\omega^2D_N(\omega)^4\le\pi^2D_N(\omega)^2.
\]
Consequently,
\[
\int_{-\pi}^{\pi}\omega^2J_N^{\mathrm J}(\omega)\,d\omega
\le
\frac{2\pi^3N}{(4\pi/3)N^3}
\le\frac{24}{N^2}
\le\frac{64}{N^2}.
\]
Finally, integrating
\[
|\omega|\le\frac N4\omega^2+\frac1N
\]
against \(J_N^{\mathrm J}\), whose integral is one, proves
\[
\int_{-\pi}^{\pi}|\omega|J_N^{\mathrm J}(\omega)\,d\omega
\le\frac N4\frac{24}{N^2}+\frac1N
\le\frac{32}{N}.
\]
% lean: Causalean.Mathlib.Analysis.JacksonApproximation.jackson_first_moment

\item
Define the nonlinear remainder and its angular version by
\[
g_\epsilon(z)
=\epsilon(1-z)_+
 +(1-2\epsilon)\frac{(z-1)_+}{1+z},
\qquad z\ge0,
\qquad
f_{M,\epsilon}(\omega)=g_\epsilon\{z_M(\omega)\},
\qquad \omega\in\mathbb R.
\]
Splitting at \(z=1\) gives
\[
\phi_\epsilon(z)=\epsilon(z-1)+g_\epsilon(z).
\]
For later integration, define the smooth curvature on the entire circle, including the kink locations, by
\[
a_{M,\epsilon}(\omega)=
\begin{cases}
\displaystyle \epsilon\frac M2\cos\omega,
&z_M(\omega)<1,\\[4pt]
\displaystyle
-\frac{4(1-2\epsilon)\{z_M'(\omega)\}^2}
       {\{1+z_M(\omega)\}^3}
+\frac{2(1-2\epsilon)z_M''(\omega)}
       {\{1+z_M(\omega)\}^2},
&z_M(\omega)>1,\\[6pt]
0,&z_M(\omega)=1.
\end{cases}
\]
On each open smooth arc this equals \(f_{M,\epsilon}''\). Since
\[
z_M'(\omega)=-\frac M2\sin\omega,
\qquad
z_M''(\omega)=-\frac M2\cos\omega,
\qquad
\{z_M'(\omega)\}^2=z_M(\omega)\{M-z_M(\omega)\},
\]
we have
\[
|a_{M,\epsilon}(\omega)|\le5M.
\]
Indeed, the lower branch is bounded by \(M/4\), while on the upper branch
\[
|a_{M,\epsilon}(\omega)|
\le
\frac{4Mz_M(\omega)}{\{1+z_M(\omega)\}^3}
+\frac{M}{\{1+z_M(\omega)\}^2}
\le5M.
\]
The scalar derivative jump is
\[
g_\epsilon'(1+)-g_\epsilon'(1-)
=\frac{1-2\epsilon}{2}+\epsilon=\frac12.
\]
Thus the angular derivative has a positive jump of size
\[
J_M^{\mathrm{jump}}=\frac{\sqrt{M-1}}2
\]
at each of \(-\vartheta_M\) and \(\vartheta_M\).

Twice integrating by parts on the three arcs
\([-\pi,-\vartheta_M]\), \([-\vartheta_M,\vartheta_M]\), and
\([\vartheta_M,\pi]\) therefore yields
\[
\begin{aligned}
\int_{-\pi}^{\pi}
f_{M,\epsilon}(\omega)k_N(\omega-\vartheta_M)\,\frac{d\omega}{2\pi}
={}&
\int_{-\pi}^{\pi}
a_{M,\epsilon}(\omega)G_N(\omega-\vartheta_M)\,
\frac{d\omega}{2\pi}\\
&+\frac{\sqrt{M-1}}{4\pi}
 \{G_N(0)+G_N(2\vartheta_M)\}.
\end{aligned}
\]
The endpoint terms at \(-\pi\) and \(\pi\) cancel by periodicity; the two remaining boundary terms are the displayed derivative jumps. The bounds in
\cref{obj:lem:jackson-packet-localization} supply universal constants
\(c_{\mathrm{loc}},C_{\mathrm{loc}}>0\) such that
\[
\frac{c_{\mathrm{loc}}}{N}\le|G_N(0)|,
\qquad
\int_{-\pi}^{\pi}|G_N(\omega)|\,\frac{d\omega}{2\pi}
\le\frac{C_{\mathrm{loc}}}{N^2}.
\]
In particular, with
\[
C_{\mathrm{sm}}=5C_{\mathrm{loc}},
\]
the smooth contribution satisfies
\[
\left|
\int_{-\pi}^{\pi}
a_{M,\epsilon}(\omega)G_N(\omega-\vartheta_M)\,
\frac{d\omega}{2\pi}
\right|
\le\frac{C_{\mathrm{sm}}M}{N^2}.
\]

To control the other kink, put
\[
\beta_M=\pi-\vartheta_M.
\]
Since \(M\ge2\), one has \(\vartheta_M\in[\pi/2,\pi)\), and hence
\[
\|2\vartheta_M\|_{\mathbb T}=2\beta_M,
\qquad
\cos\beta_M=1-\frac2M.
\]
The inequality \(1-\cos\beta_M\le\beta_M^2/2\) gives
\[
\beta_M\ge\frac2{\sqrt M},
\qquad
\|2\vartheta_M\|_{\mathbb T}\ge\frac4{\sqrt M}.
\]
Applying the localized bound for \(G_N\) gives, with
\[
C_{\mathrm{opp}}=\frac{C_{\mathrm{loc}}}{16},
\]
\[
|G_N(2\vartheta_M)|
\le
\frac{C_{\mathrm{loc}}}
     {N\{1+N\|2\vartheta_M\|_{\mathbb T}\}^2}
\le\frac{C_{\mathrm{opp}}M}{N^3}.
\]

Choose a universal integer \(A_0\) satisfying
\[
A_0>
\max\left\{
1,\frac{2C_{\mathrm{opp}}}{c_{\mathrm{loc}}},
\frac{32\pi C_{\mathrm{sm}}}{c_{\mathrm{loc}}}
\right\},
\qquad
N=A_0K.
\]
Then
\[
2C_{\mathrm{opp}}\le c_{\mathrm{loc}}A_0^2,
\qquad
32\pi C_{\mathrm{sm}}\le c_{\mathrm{loc}}A_0,
\]
so \(M\le K^2\) implies
\[
|G_N(2\vartheta_M)|\le\frac{c_{\mathrm{loc}}}{2N},
\qquad
|G_N(0)+G_N(2\vartheta_M)|
\ge\frac{c_{\mathrm{loc}}}{2N}.
\]
Also,
\[
\sqrt{M-1}\ge\frac{\sqrt M}{2},
\qquad
\frac{C_{\mathrm{sm}}M}{N^2}
\le\frac{c_{\mathrm{loc}}\sqrt M}{32\pi N}.
\]
The integration-by-parts identity consequently proves
\[
\left|
\int_{-\pi}^{\pi}
f_{M,\epsilon}(\omega)k_N(\omega-\vartheta_M)\,
\frac{d\omega}{2\pi}
\right|
\ge
\frac{c_{\mathrm{loc}}\sqrt M}{32\pi N}.
\]
% lean: CausalSmith.Stat.OptvalueVanishingoverlapRate.packetRemainder_response_lower

\item
Define the two-center packet and its normalization by
\[
H_{K,M}(\omega)
=k_N(\omega-\vartheta_M)+k_N(\omega+\vartheta_M),
\qquad
W_{K,M}
=\int_{-\pi}^{\pi}
\{1+z_M(\omega)\}|H_{K,M}(\omega)|\,\frac{d\omega}{2\pi}.
\]
The frequency exclusions in
\cref{obj:lem:jackson-packet-localization} imply
\[
\int_{-\pi}^{\pi}
P_{\mathrm{test}}\{z_M(\omega)\}H_{K,M}(\omega)\,
\frac{d\omega}{2\pi}=0,
\qquad
\deg P_{\mathrm{test}}\le K.
\]
Indeed, composing such a polynomial with \(z_M\) produces a trigonometric polynomial with frequencies of absolute value at most \(K\), whereas the translated packets have frequencies between \(2N+2\) and \(6N-2\). Translation preserves these exclusions.

Both \(f_{M,\epsilon}\) and \(k_N\) are even, so their pairings with the two translates agree. The affine term in
\(\phi_\epsilon=\epsilon(z-1)+g_\epsilon\) is annihilated. Thus
\[
\begin{aligned}
\left|
\int_{-\pi}^{\pi}
\phi_\epsilon\{z_M(\omega)\}H_{K,M}(\omega)\,
\frac{d\omega}{2\pi}
\right|
&=
2\left|
\int_{-\pi}^{\pi}
f_{M,\epsilon}(\omega)k_N(\omega-\vartheta_M)\,
\frac{d\omega}{2\pi}
\right|\\
&\ge
\frac{c_{\mathrm{loc}}}{16\pi A_0}\frac{\sqrt M}{K}.
\end{aligned}
\]

We next verify the normalization over the whole parameter range. The nonnegative Jackson kernel has a strict maximum at the periodic origin: the geometric-sum representation of \(D_N\) gives
\[
|D_N(\omega)|<N
\quad\text{when }\cos\omega<1,
\qquad
|D_N(0)|=N.
\]
Since \(\cos(2\vartheta_M)<1\),
\[
H_{K,M}(\vartheta_M)
=\mathcal J_N(0)
 +\mathcal J_N(2\vartheta_M)\cos(8N\vartheta_M)
\ge\mathcal J_N(0)-\mathcal J_N(2\vartheta_M)>0.
\]
Continuity and nonnegativity of the integrand defining \(W_{K,M}\) therefore give \(W_{K,M}>0\).

For the upper bound, use
\[
|H_{K,M}(\omega)|
\le
\mathcal J_N(\omega-\vartheta_M)
+\mathcal J_N(\omega+\vartheta_M)
\]
and the identity
\[
z_M(\omega+\vartheta_M)+z_M(\omega-\vartheta_M)
=2+(M-2)(1-\cos\omega).
\]
After translating the two periodic integrals,
\[
\begin{aligned}
W_{K,M}
&\le
\int_{-\pi}^{\pi}
\{4+(M-2)(1-\cos\omega)\}J_N^{\mathrm J}(\omega)\,d\omega\\
&\le
4+\frac{M-2}{2}
 \int_{-\pi}^{\pi}\omega^2J_N^{\mathrm J}(\omega)\,d\omega\\
&\le4+\frac{32(M-2)}{N^2}
\le36.
\end{aligned}
\]
Define
\[
c_{\mathrm p}=\frac{c_{\mathrm{loc}}}{576\pi A_0}>0.
\]
Dividing the packet response by \(W_{K,M}\le36\) proves
\[
\left|
\int_{-\pi}^{\pi}
\phi_\epsilon\{z_M(\omega)\}
\frac{H_{K,M}(\omega)}{W_{K,M}}\,
\frac{d\omega}{2\pi}
\right|
\ge c_{\mathrm p}\frac{\sqrt M}{K}.
\]
% lean: CausalSmith.Stat.OptvalueVanishingoverlapRate.packetNormalized_response_lower

\item
The packet parts in the statement are explicitly
\[
\sigma_+
=(z_M)_\#
\left[
\left(\frac{H_{K,M}}{W_{K,M}}\right)_+
\frac{d\omega}{2\pi}
\right],
\qquad
\sigma_-
=(z_M)_\#
\left[
\left(-\frac{H_{K,M}}{W_{K,M}}\right)_+
\frac{d\omega}{2\pi}
\right],
\]
where both measures inside the brackets are restricted to
\([-\pi,\pi]\). They are finite and supported on \([0,M]\), and therefore have integrable first moments. The packet is even, and the fibers of \(z_M\) consist of points with equal packet sign; these are the positive and negative Jordan parts of the pushed signed packet.

The annihilation and normalization established above give
\[
\int z^\ell\,d\sigma_+
=\int z^\ell\,d\sigma_-,
\qquad \ell=0,\ldots,K,
\]
\[
\int(1+z)\,d(\sigma_++\sigma_-)=1.
\]
Writing, as in \cref{obj:def:weighted-handle},
\[
t=\sigma_+(\mathbb R),
\qquad
u=\int z\,d\sigma_+,
\qquad
z_0=\frac{1-u}{1-t},
\]
the degree-zero and degree-one identities imply
\[
\sigma_-(\mathbb R)=t,
\qquad
\int z\,d\sigma_-=u,
\qquad
t+u=\frac12.
\]
Since \(t,u\ge0\),
\[
0\le t\le\frac12,
\qquad
z_0=\frac{1/2+t}{1-t},
\qquad
\frac12\le z_0\le2\le M.
\]
In particular \(t<1\), so the completion map is defined, and its positive-part coefficient equals \(1-t\). Set
\[
\nu_0=\sigma_-+(1-t)\delta_{z_0},
\qquad
\nu_1=\sigma_++(1-t)\delta_{z_0}.
\]
Their masses and first moments satisfy
\[
\nu_0(\mathbb R)=\nu_1(\mathbb R)=t+(1-t)=1,
\qquad
\int z\,d\nu_0=\int z\,d\nu_1=u+(1-t)z_0=1.
\]
The common atom cancels in every moment difference, so
\[
(\nu_0,\nu_1)\in\mathfrak N_{K,M},
\qquad
\mathscr H^w(\sigma_+,\sigma_-)=(\nu_0,\nu_1).
\]
It also cancels in the target difference:
\[
\begin{aligned}
\int\phi_\epsilon\,d\nu_1-\int\phi_\epsilon\,d\nu_0
&=
\int_{-\pi}^{\pi}
\phi_\epsilon\{z_M(\omega)\}
\frac{H_{K,M}(\omega)}{W_{K,M}}\,
\frac{d\omega}{2\pi}.
\end{aligned}
\]
Hence the explicit completion has gap at least
\(c_{\mathrm p}\sqrt M/K\).
% lean: CausalSmith.Stat.OptvalueVanishingoverlapRate.packetCompletion_separation_lower

\item
For any real polynomial \(P_{\mathrm{test}}\) of degree at most \(K\), define
\[
e(P_{\mathrm{test}})
=\sup_{z\in[0,M]}
\frac{|\phi_\epsilon(z)-P_{\mathrm{test}}(z)|}{1+z}.
\]
The supremum is finite by continuity. For any feasible prior pair, polynomial moment matching gives
\[
\int P_{\mathrm{test}}\,d\nu_1
=\int P_{\mathrm{test}}\,d\nu_0,
\]
and probability normalization with mean one gives
\[
\int(1+z)\,d\nu_0=\int(1+z)\,d\nu_1=2.
\]
Therefore
\[
\begin{aligned}
\left|\int\phi_\epsilon\,d\nu_1-\int\phi_\epsilon\,d\nu_0\right|
&=
\left|
\int(\phi_\epsilon-P_{\mathrm{test}})\,d\nu_1
-\int(\phi_\epsilon-P_{\mathrm{test}})\,d\nu_0
\right|\\
&\le
e(P_{\mathrm{test}})
\sum_{\iota=0}^1\int(1+z)\,d\nu_\iota
=4e(P_{\mathrm{test}}).
\end{aligned}
\]
Applying this bound to the packet completion and then taking the polynomial infimum yields, with
\[
c_{\mathrm E}=\frac{c_{\mathrm p}}4,
\]
\[
c_{\mathrm E}\frac{\sqrt M}{K}\le\mathcal E.
\]
In particular \(\mathcal E>0\).
% lean: CausalSmith.Stat.OptvalueVanishingoverlapRate.packetWeightedApproxError_lower

\item
We now prove the reverse comparison \(\mathcal E\le\Delta\) by a finite weighted dual construction. Define
\[
F_{\mathrm w}(z)=\frac{\phi_\epsilon(z)}{1+z},
\qquad
\mathcal V_K^{\mathrm w}
=\left\{
z\mapsto\frac{P_{\mathrm{test}}(z)}{1+z}:
\deg P_{\mathrm{test}}\le K
\right\}
\subset C([0,M],\mathbb R).
\]
This is a finite-dimensional subspace. Its closed ball of radius
\[
R_{\mathrm{best}}=2\|F_{\mathrm w}\|_\infty+1
\]
is compact, so the distance from \(F_{\mathrm w}\) attains a minimum there. The zero function gives distance at most \(\|F_{\mathrm w}\|_\infty\), while a function outside this ball has distance at least
\[
\|v_{\mathrm w}\|_\infty-\|F_{\mathrm w}\|_\infty
>\|F_{\mathrm w}\|_\infty+1.
\]
Thus a polynomial \(P_{\mathrm{best}}\) attains the global infimum:
\[
\deg P_{\mathrm{best}}\le K,
\qquad
e(P_{\mathrm{best}})=\mathcal E.
\]
Define its weighted residual and extremal set by
\[
r_{\mathrm{best}}(z)
=\frac{\phi_\epsilon(z)-P_{\mathrm{best}}(z)}{1+z},
\qquad
\mathcal A_{\mathrm{ext}}
=\{z\in[0,M]:|r_{\mathrm{best}}(z)|=\mathcal E\}.
\]

There exist ordered nodes and an orientation satisfying
\[
0\le z_0^{\mathrm{dual}}<\cdots<z_{K+1}^{\mathrm{dual}}\le M,
\qquad
o_{\mathrm{dual}}\in\{-1,1\},
\]
\[
r_{\mathrm{best}}(z_\ell^{\mathrm{dual}})
=o_{\mathrm{dual}}(-1)^\ell\mathcal E,
\qquad \ell=0,\ldots,K+1.
\]
To see this, suppose such nodes are absent. The positive and negative extremal sets are disjoint compact sets. If one is empty, the constant polynomial with the other sign agrees with the residual sign on all extrema. If both are nonempty, their positive separation permits a finite ordered cover by sufficiently small neighborhoods. Consecutive sign blocks in that cover have at most \(K\) transitions: \(K+1\) transitions would furnish the forbidden \(K+2\) alternating nodes. Placing one root in each transition gap and choosing the overall sign produces a polynomial \(Q_{\mathrm{sign}}\) such that
\[
\deg Q_{\mathrm{sign}}\le K,
\qquad
r_{\mathrm{best}}(z)Q_{\mathrm{sign}}(z)>0
\quad(z\in\mathcal A_{\mathrm{ext}}).
\]
Put
\[
v_{\mathrm{sign}}(z)=\frac{Q_{\mathrm{sign}}(z)}{1+z}.
\]
Define the set where the perturbation has the wrong sign or vanishes by
\[
\mathcal B_{\mathrm{bad}}
=\{z\in[0,M]:r_{\mathrm{best}}(z)v_{\mathrm{sign}}(z)\le0\}.
\]
This set is compact. If it is nonempty, the continuous absolute residual attains its maximum there, and
\[
\max_{z\in\mathcal B_{\mathrm{bad}}}|r_{\mathrm{best}}(z)|<\mathcal E.
\]
Indeed, equality at a maximizing point would make that point an extremum, where the residual and perturbation have strictly positive product. Define the threshold in both cases by
\[
b_{\mathrm{ext}}=
\begin{cases}
\displaystyle\frac{\max_{z\in\mathcal B_{\mathrm{bad}}}|r_{\mathrm{best}}(z)|+\mathcal E}{2},
&\mathcal B_{\mathrm{bad}}\ne\varnothing,\\[6pt]
\mathcal E/2,&\mathcal B_{\mathrm{bad}}=\varnothing.
\end{cases}
\]
Thus \(b_{\mathrm{ext}}<\mathcal E\), and in either case
\[
|r_{\mathrm{best}}(z)|\ge b_{\mathrm{ext}}
\quad\Longrightarrow\quad
r_{\mathrm{best}}(z)v_{\mathrm{sign}}(z)>0.
\]
The high-residual set
\[
\mathcal A_{\mathrm{high}}
=\{z\in[0,M]:|r_{\mathrm{best}}(z)|\ge b_{\mathrm{ext}}\}
\]
is compact and nonempty, since the residual attains its norm \(\mathcal E\). Consequently the following extrema exist:
\[
m_{\mathrm{ext}}
=\min_{z\in\mathcal A_{\mathrm{high}}}
 r_{\mathrm{best}}(z)v_{\mathrm{sign}}(z)>0,
\qquad
B_{\mathrm{sign}}
=\max_{z\in[0,M]}|v_{\mathrm{sign}}(z)|>0.
\]
The second constant is positive because the perturbation is nonzero at every residual extremum. These definitions give
\[
r_{\mathrm{best}}(z)v_{\mathrm{sign}}(z)\ge m_{\mathrm{ext}}
\quad\text{if }|r_{\mathrm{best}}(z)|\ge b_{\mathrm{ext}},
\qquad
|v_{\mathrm{sign}}(z)|\le B_{\mathrm{sign}}
\quad(z\in[0,M]).
\]
Choose
\[
\eta_{\mathrm{pert}}
=\min\left\{
\frac{m_{\mathrm{ext}}}{B_{\mathrm{sign}}^2},
\frac{\mathcal E-b_{\mathrm{ext}}}{2B_{\mathrm{sign}}}
\right\}>0.
\]
On the first region,
\[
\begin{aligned}
\{r_{\mathrm{best}}-\eta_{\mathrm{pert}}v_{\mathrm{sign}}\}^2
&\le
\mathcal E^2-2\eta_{\mathrm{pert}}m_{\mathrm{ext}}
+\eta_{\mathrm{pert}}^2B_{\mathrm{sign}}^2\\
&\le\mathcal E^2-\eta_{\mathrm{pert}}m_{\mathrm{ext}}
<\mathcal E^2.
\end{aligned}
\]
On its complement,
\[
|r_{\mathrm{best}}-\eta_{\mathrm{pert}}v_{\mathrm{sign}}|
<
b_{\mathrm{ext}}+\frac{\mathcal E-b_{\mathrm{ext}}}{2}
<\mathcal E.
\]
Continuity on the compact interval makes the supremum strictly smaller than \(\mathcal E\), contradicting optimality of
\(P_{\mathrm{best}}+\eta_{\mathrm{pert}}Q_{\mathrm{sign}}\). This proves the alternating-node assertion.
% lean: CausalSmith.Stat.OptvalueVanishingoverlapRate.weightedApproxError_alternating_witness

\item
For these nodes, define the Lagrange weights and their two normalizations by
\[
\lambda_\ell^{\mathrm{Lag}}
=\prod_{\substack{0\le r\le K+1\\r\ne\ell}}
(z_\ell^{\mathrm{dual}}-z_r^{\mathrm{dual}})^{-1},
\qquad
\lambda_\ell^{\mathrm{norm}}
=\frac{\lambda_\ell^{\mathrm{Lag}}}
       {\sum_{r=0}^{K+1}|\lambda_r^{\mathrm{Lag}}|},
\]
\[
D_{\mathrm{dual}}
=\sum_{\ell=0}^{K+1}
|\lambda_\ell^{\mathrm{norm}}|(1+z_\ell^{\mathrm{dual}}),
\qquad
\lambda_\ell^{\mathrm{dual}}
=\frac{\lambda_\ell^{\mathrm{norm}}}{D_{\mathrm{dual}}}.
\]
The node differences are nonzero, and the two normalization denominators are positive; moreover, \(D_{\mathrm{dual}}\ge1\).
Comparing the coefficient of degree \(K+1\) in the Lagrange interpolation formula gives
\[
\sum_{\ell=0}^{K+1}
\lambda_\ell^{\mathrm{norm}}P_{\mathrm{test}}(z_\ell^{\mathrm{dual}})
=0
\qquad(\deg P_{\mathrm{test}}\le K).
\]
The ordered nodes also give
\[
\lambda_\ell^{\mathrm{norm}}
=(-1)^{K+1-\ell}|\lambda_\ell^{\mathrm{norm}}|.
\]
Combining this sign with the alternating residual yields
\[
\begin{aligned}
\sum_{\ell=0}^{K+1}
\lambda_\ell^{\mathrm{norm}}\phi_\epsilon(z_\ell^{\mathrm{dual}})
&=
\sum_{\ell=0}^{K+1}
\lambda_\ell^{\mathrm{norm}}
\{\phi_\epsilon(z_\ell^{\mathrm{dual}})
-P_{\mathrm{best}}(z_\ell^{\mathrm{dual}})\}\\
&=o_{\mathrm{dual}}(-1)^{K+1}\mathcal E D_{\mathrm{dual}}.
\end{aligned}
\]
Consequently,
\[
\sum_{\ell=0}^{K+1}
\lambda_\ell^{\mathrm{dual}}(z_\ell^{\mathrm{dual}})^r=0
\quad(r=0,\ldots,K),
\]
\[
\sum_{\ell=0}^{K+1}
|\lambda_\ell^{\mathrm{dual}}|(1+z_\ell^{\mathrm{dual}})=1,
\qquad
\left|
\sum_{\ell=0}^{K+1}
\lambda_\ell^{\mathrm{dual}}\phi_\epsilon(z_\ell^{\mathrm{dual}})
\right|=\mathcal E.
\]
Split the weights into positive and negative parts:
\[
\sigma_+^{\mathrm{dual}}
=\sum_{\ell=0}^{K+1}(\lambda_\ell^{\mathrm{dual}})_+
 \delta_{z_\ell^{\mathrm{dual}}},
\qquad
\sigma_-^{\mathrm{dual}}
=\sum_{\ell=0}^{K+1}(-\lambda_\ell^{\mathrm{dual}})_+
 \delta_{z_\ell^{\mathrm{dual}}}.
\]
Define
\[
t_{\mathrm{dual}}=\sigma_+^{\mathrm{dual}}(\mathbb R),
\qquad
u_{\mathrm{dual}}=\int z\,d\sigma_+^{\mathrm{dual}},
\qquad
z_{\mathrm{dual}}^{\mathrm{comp}}
=\frac{1-u_{\mathrm{dual}}}{1-t_{\mathrm{dual}}}.
\]
The degree-zero and degree-one identities and the weighted normalization give
\[
t_{\mathrm{dual}}+u_{\mathrm{dual}}=\frac12,
\qquad
0\le t_{\mathrm{dual}}\le\frac12,
\qquad
\frac12\le z_{\mathrm{dual}}^{\mathrm{comp}}\le2.
\]
The same mass, mean, and moment calculation used for the packet completion shows that
\[
\nu_0^{\mathrm{dual}}
=\sigma_-^{\mathrm{dual}}
 +(1-t_{\mathrm{dual}})\delta_{z_{\mathrm{dual}}^{\mathrm{comp}}},
\qquad
\nu_1^{\mathrm{dual}}
=\sigma_+^{\mathrm{dual}}
 +(1-t_{\mathrm{dual}})\delta_{z_{\mathrm{dual}}^{\mathrm{comp}}}
\]
form a pair in \(\mathfrak N_{K,M}\), with target gap \(\mathcal E\). Hence
\[
\mathcal E\le\Delta.
\]
The feasible set is nonempty, since \((\delta_1,\delta_1)\) is feasible. Its gaps are bounded by \(4e(0)\), by the polynomial residual bound above. Taking the polynomial infimum in that bound, followed by the feasible-pair supremum, therefore gives
\[
\Delta\le4\mathcal E.
\]
% lean: CausalSmith.Stat.OptvalueVanishingoverlapRate.weightedApproxError_le_constrainedPriorSeparation

\item
For the approximation upper bound, the square-root coordinate gives the required uniform angular regularity. Define, for \(\xi\ge0\),
\[
b_{\mathrm{lo}}(\xi)=(1-\xi^2)_+,
\qquad
b_{\mathrm{hi}}(\xi)=\frac{(\xi^2-1)_+}{1+\xi^2}.
\]
For \(0\le\xi\le\eta\), splitting at \(1\) proves
\[
|b_{\mathrm{lo}}(\xi)-b_{\mathrm{lo}}(\eta)|
\le2(\eta-\xi),
\qquad
|b_{\mathrm{hi}}(\xi)-b_{\mathrm{hi}}(\eta)|
\le2(\eta-\xi).
\]
For the lower branch, when \(\eta\le1\), the difference equals
\((\eta-\xi)(\eta+\xi)\); when \(\xi\le1\le\eta\), it equals
\(1-\xi^2\le2(1-\xi)\); when \(1\le\xi\), it vanishes.
For the upper branch, when \(1\le\xi\),
\[
b_{\mathrm{hi}}(\eta)-b_{\mathrm{hi}}(\xi)
=\frac{2(\eta-\xi)(\eta+\xi)}
       {(1+\xi^2)(1+\eta^2)}
\le2(\eta-\xi).
\]
When \(\xi\le1\le\eta\), the bound follows from
\[
\frac{\eta^2-1}{1+\eta^2}\le2(\eta-1)\le2(\eta-\xi);
\]
when \(\eta\le1\), the difference vanishes. Interchanging \(\xi,\eta\) covers the opposite ordering. It follows that
\[
|g_\epsilon(\xi^2)-g_\epsilon(\eta^2)|
\le
2\{\epsilon+(1-2\epsilon)\}|\xi-\eta|
\le2|\xi-\eta|.
\]
Using the half-angle identity
\[
z_M(\omega)=\{\sqrt M\,|\cos(\omega/2)|\}^2
\]
and the Lipschitz bound for cosine gives
\[
|f_{M,\epsilon}(\omega)-f_{M,\epsilon}(\upsilon)|
\le\sqrt M\,|\omega-\upsilon|,
\qquad \omega,\upsilon\in\mathbb R.
\]

Set
\[
N_{\mathrm{up}}=\lfloor K/2\rfloor+1,
\qquad
T_{\mathrm{up}}(\omega)
=\int_{-\pi}^{\pi}
f_{M,\epsilon}(\omega+\upsilon)
J_{N_{\mathrm{up}}}^{\mathrm J}(\upsilon)\,d\upsilon.
\]
Convolution with the even Jackson kernel makes \(T_{\mathrm{up}}\) an even trigonometric polynomial of degree at most
\(2(N_{\mathrm{up}}-1)\le K\). Positivity, normalization, and the first kernel moment give
\[
|f_{M,\epsilon}(\omega)-T_{\mathrm{up}}(\omega)|
\le
\sqrt M\int_{-\pi}^{\pi}
|\upsilon|J_{N_{\mathrm{up}}}^{\mathrm J}(\upsilon)\,d\upsilon
\le\frac{32\sqrt M}{N_{\mathrm{up}}}.
\]
Writing its cosine modes as Chebyshev polynomials supplies a real polynomial \(Q_{\mathrm{up}}\) such that
\[
T_{\mathrm{up}}(\omega)=Q_{\mathrm{up}}(\cos\omega),
\qquad
\deg Q_{\mathrm{up}}\le2(N_{\mathrm{up}}-1).
\]
Define
\[
P_{\mathrm{up}}(z)
=\epsilon(z-1)+Q_{\mathrm{up}}(2z/M-1).
\]
Then \(\deg P_{\mathrm{up}}\le K\), and every \(z\in[0,M]\) can be written as
\(z_M(\omega)\). Since \(K\le2N_{\mathrm{up}}\),
\[
|\phi_\epsilon(z)-P_{\mathrm{up}}(z)|
\le\frac{32\sqrt M}{N_{\mathrm{up}}}
\le\frac{64\sqrt M}{K}.
\]
Division by \(1+z\ge1\), followed by the polynomial infimum, proves
\[
\mathcal E\le\frac{64\sqrt M}{K}.
\]
% lean: CausalSmith.Stat.OptvalueVanishingoverlapRate.weightedApproxError_upper

\item
Choose the universal constants
\[
c=\min\{\min(c_{\mathrm E},c_{\mathrm p}),1\},
\qquad
C=257.
\]
They satisfy \(0<c<C\), and the preceding bounds yield
\[
c\frac{\sqrt M}{K}
\le\mathcal E
\le\Delta
\le4\mathcal E
\le256\frac{\sqrt M}{K}
\le C\frac{\sqrt M}{K}.
\]
The packet completion retains its gap because \(c\le c_{\mathrm p}\).
All the packet estimates used the closed range
\(0\le\epsilon\le1/2\) and the non-strict support bound \(M\le K^2\), so they include both overlap endpoints and \(M=K^2\).
% lean: CausalSmith.Stat.OptvalueVanishingoverlapRate.weighted_approximation_and_packet

\item
\textbf{Sparse comparison.}\quad
We next establish an observed minimax lower bound for all alphabets. For this part, assume
\[
n\ge1,\qquad d\ge2,\qquad 0<\epsilon\le\tfrac12,
\]
and define
\[
L_d=\log(ed),
\qquad
r_{n,d,\epsilon}=\frac{d}{n\epsilon L_d},
\qquad
\rho_{n,d,\epsilon}=\min\{1,r_{n,d,\epsilon}\}.
\]
In particular \(L_d=1+\log d\ge1\) and \(\rho_{n,d,\epsilon}\ge0\).

Here is the likelihood calibration used in both large-alphabet constructions. For integers \(K_{\mathrm{test}}\ge1\) and real \(\tau\ge0\), define
\[
\mathcal T_{K_{\mathrm{test}}}(\tau)
=\sum_{\ell>K_{\mathrm{test}}}\frac{\tau^\ell}{\ell!},
\qquad
\beta_{\mathrm{tail}}=\frac1{16e}.
\]
When \(\tau\le\beta_{\mathrm{tail}}K_{\mathrm{test}}\), the factorial bound
\(\ell!\ge(\ell/e)^\ell\) gives, for \(\ell>K_{\mathrm{test}}\),
\[
\frac{\tau^\ell}{\ell!}
\le\left(\frac{e\tau}{\ell}\right)^\ell
\le4^{-\ell}.
\]
Summing the geometric tail therefore gives
\[
\mathcal T_{K_{\mathrm{test}}}(\tau)
\le4^{-K_{\mathrm{test}}},
\qquad
\sqrt{\mathcal T_{K_{\mathrm{test}}}(\tau)}
\le2^{-K_{\mathrm{test}}}.
\]
Define
\[
C_{\mathrm{tv}}=\frac{|\log8|+2}{\log2}.
\]
If \(K_{\mathrm{test}}\ge C_{\mathrm{tv}}L_d\), then
\[
K_{\mathrm{test}}\log2
\ge(|\log8|+2)(1+\log d)
\ge\log(8d),
\]
and consequently
\[
d\sqrt{\mathcal T_{K_{\mathrm{test}}}(\tau)}
\le d\,2^{-K_{\mathrm{test}}}\le\frac18.
\]
% lean: CausalSmith.Stat.OptvalueVanishingoverlapRate.sparse_geometric_degree_calibration

\item
Consider first the sparse experiment of
\cref{obj:lem:sparse-likelihood}, with inflation \(B>2\) and a constrained pair of degree \(K_{\mathrm{test}}\). Under prior \(\nu_\iota\), \(\iota\in\{0,1\}\), draw cell intensities independently with that prior, and observe all four independent Poisson counts in each cell.

For one cell, let \(\mathbb Q_\star\) be the reference count law at \(z_\star=M/2\), and define the mixture densities
\[
h_{\mathrm{sp},\iota}
=\int L_z\,d\nu_\iota(z),
\qquad \iota\in\{0,1\}.
\]
The likelihood ratios \(L_z\) and Gram coefficient \(\alpha\) are those in
\cref{obj:lem:sparse-likelihood}; at a zero alternative mean, the count factor uses \(0^0=1\). Its Gram identity, followed by expansion of the exponential, gives
\[
\begin{aligned}
\int(h_{\mathrm{sp},1}-h_{\mathrm{sp},0})^2\,d\mathbb Q_\star
&=
\sum_{\ell>K_{\mathrm{test}}}\frac{\alpha^\ell}{\ell!}
\left\{
\int(z-z_\star)^\ell\,d(\nu_1-\nu_0)(z)
\right\}^2\\
&\le
4\mathcal T_{K_{\mathrm{test}}}(\tau_{\mathrm{sp}}),
\qquad
\tau_{\mathrm{sp}}=\frac{\alpha M^2}{4}.
\end{aligned}
\]
The terms through degree \(K_{\mathrm{test}}\) vanish by moment matching, and the remaining moments have absolute value at most
\(2(M/2)^\ell\). Thus Cauchy--Schwarz and telescoping products imply
\[
\operatorname{TV}
\left(\mathbb Q_{\mathrm{sp},0}^{\otimes d},
      \mathbb Q_{\mathrm{sp},1}^{\otimes d}\right)
\le
d\sqrt{\mathcal T_{K_{\mathrm{test}}}(\tau_{\mathrm{sp}})},
\]
where
\[
\mathbb Q_{\mathrm{sp},\iota}
=h_{\mathrm{sp},\iota}\,\mathbb Q_\star.
\]
Set
\[
\eta_{\mathrm{sp}}=2\beta_{\mathrm{tail}}.
\]
The bound on \(\alpha\) in
\cref{obj:lem:sparse-likelihood} gives
\[
\tau_{\mathrm{sp}}\le\frac{Bn\epsilon M}{2d}.
\]
Hence
\[
M\le\frac{\eta_{\mathrm{sp}}dK_{\mathrm{test}}}{Bn\epsilon},
\qquad
K_{\mathrm{test}}\ge C_{\mathrm{tv}}L_d
\]
ensure that the full product-mixture total variation is at most \(1/8\).
% lean: CausalSmith.Stat.OptvalueVanishingoverlapRate.sparse_product_mixture_tv_le_one_eighth

\item
For either sparse product prior, write the independent cell intensities as
\[
\mathbf Z^{\mathrm{sp}}
=(Z_x^{\mathrm{sp}})_{x\in[d]}
\sim\nu_\iota^{\otimes d},
\qquad
\mathbf Z^{\mathrm{sp}}\in[0,M]^d,
\]
and define
\[
H_{\mathrm{sp}}
=\frac1d\sum_{x\in[d]}\phi_\epsilon(Z_x^{\mathrm{sp}}),
\qquad
S_{\mathrm{sp}}
=1-\epsilon+\frac{\epsilon}{d}\sum_{x\in[d]}Z_x^{\mathrm{sp}},
\]
\[
\theta_{\mathrm{sp},\iota}
=\frac12+\frac12\int\phi_\epsilon\,d\nu_\iota,
\qquad
\delta_{\mathrm{sp}}
=|\theta_{\mathrm{sp},1}-\theta_{\mathrm{sp},0}|
=\frac12\left|\int\phi_\epsilon\,d(\nu_1-\nu_0)\right|.
\]
The sparse value identity in
\cref{obj:lem:sparse-likelihood} is
\[
\Psi(\mathbb P_{\mathbf Z^{\mathrm{sp}}}^{\mathrm{sp}})
=\frac12+\frac{H_{\mathrm{sp}}}{2S_{\mathrm{sp}}}.
\]
The mean-one constraint and \(0\le Z_x^{\mathrm{sp}}\le M\) give
\[
E (Z_x^{\mathrm{sp}})^2\le M,
\qquad
E S_{\mathrm{sp}}=1,
\qquad
E(S_{\mathrm{sp}}-1)^2\le\frac{\epsilon^2M}{d}.
\]
For \(0\le z\le M\), both \(z-1\) and zero are at most \(1+z\), so
\[
0\le(z-1)_+\le1+z,
\qquad
0\le1-\epsilon+\epsilon z,
\qquad
1+z>0.
\]
Multiplying the first inequality by the nonnegative factor
\((1-\epsilon+\epsilon z)/(1+z)\) gives
\[
0\le\phi_\epsilon(z)\le1-\epsilon+\epsilon z.
\]
Moreover, \(z^2\le Mz\), \((1-\epsilon)^2\le1\), and the nonnegativity of
\((1-\epsilon-\epsilon z)^2\) yields
\[
\begin{aligned}
\phi_\epsilon(z)^2
&\le(1-\epsilon+\epsilon z)^2\\
&\le2\{(1-\epsilon)^2+\epsilon^2z^2\}\\
&\le2(1+\epsilon^2Mz).
\end{aligned}
\]
Independence therefore yields
\[
E\{H_{\mathrm{sp}}-EH_{\mathrm{sp}}\}^2
\le\frac{2(1+\epsilon^2M)}d.
\]
Since \(S_{\mathrm{sp}}\ge1-\epsilon\ge1/2\) and
\(0\le H_{\mathrm{sp}}\le S_{\mathrm{sp}}\),
\[
\left|\frac{H_{\mathrm{sp}}}{S_{\mathrm{sp}}}
-H_{\mathrm{sp}}\right|
\le|1-S_{\mathrm{sp}}|.
\]
Using \((a+b)^2\le2(a^2+b^2)\) now gives
\[
E\left[
\left\{\Psi(\mathbb P_{\mathbf Z^{\mathrm{sp}}}^{\mathrm{sp}})
-\theta_{\mathrm{sp},\iota}\right\}^2
\right]
\le\frac{2+3\epsilon^2M}{2d}.
\]
Consequently, if
\[
64(2+3\epsilon^2M)\le d\delta_{\mathrm{sp}}^2,
\]
Markov's inequality gives
\[
\Pr_\iota\left(
\left|\Psi(\mathbb P_{\mathbf Z^{\mathrm{sp}}}^{\mathrm{sp}})
-\theta_{\mathrm{sp},\iota}\right|
>\frac{\delta_{\mathrm{sp}}}{4}
\right)\le\frac18.
\]

For completeness, define the two integrated squared risks of a count estimator \(\widehat W\) by
\[
\mathcal B_{\mathrm{sp},\iota}(\widehat W)
=
E_{\mathbf Z^{\mathrm{sp}}\sim\nu_\iota^{\otimes d}}
E_{\mathrm{counts}\mid\mathbf Z^{\mathrm{sp}}}
\left[
\{\widehat W-\Psi(\mathbb P_{\mathbf Z^{\mathrm{sp}}}^{\mathrm{sp}})\}^2
\right].
\]
If \(\theta_{\mathrm{sp},0}\le\theta_{\mathrm{sp},1}\), threshold
\(\widehat W\) at
\[
\frac{\theta_{\mathrm{sp},0}+\theta_{\mathrm{sp},1}}2.
\]
Otherwise interchange the prior indices and use the same midpoint. On an erroneous decision with the target within
\(\delta_{\mathrm{sp}}/4\) of its center, estimation error is at least
\(\delta_{\mathrm{sp}}/4\). The sum of the two testing errors is at least one minus their predictive total variation. Subtracting the two target-deviation probabilities thus proves
\[
\begin{aligned}
\max_{\iota=0,1}\mathcal B_{\mathrm{sp},\iota}(\widehat W)
&\ge
\frac12\left(\frac{\delta_{\mathrm{sp}}}{4}\right)^2
\left(1-\frac18-\frac18-\frac18\right)\\
&=\frac{5\delta_{\mathrm{sp}}^2}{256}.
\end{aligned}
\]

Use the completed packet pair, whose gap is at least
\(c_{\mathrm p}\sqrt M/K_{\mathrm{test}}\). Then
\[
\delta_{\mathrm{sp}}^2
\ge\frac{c_{\mathrm p}^2M}{4K_{\mathrm{test}}^2}.
\]
Thus the condition
\[
256(2+3\epsilon^2M)K_{\mathrm{test}}^2
\le c_{\mathrm p}^2Md
\]
supplies the target-deviation bound, and, with
\[
c_{\mathrm{sp}}=\frac{5c_{\mathrm p}^2}{1024},
\]
every measurable count estimator satisfies
\[
\max_{\iota=0,1}\mathcal B_{\mathrm{sp},\iota}(\widehat W)
\ge c_{\mathrm{sp}}\frac{M}{K_{\mathrm{test}}^2}.
\]
% lean: CausalSmith.Stat.OptvalueVanishingoverlapRate.sparse_packet_bayesRisk_lower

\item
We record the logarithmic degree and alphabet calibration explicitly. For real \(C_{\mathrm{deg}}\ge2\) and integer \(d\ge2\), define
\[
\mathcal K_{\mathrm{cal}}(C_{\mathrm{deg}},d)
=2\left\lceil\frac{C_{\mathrm{deg}}L_d}{2}\right\rceil.
\]
The ceiling inequalities and \(L_d\ge1\) give
\[
2\le\mathcal K_{\mathrm{cal}}(C_{\mathrm{deg}},d),
\qquad
C_{\mathrm{deg}}L_d
\le\mathcal K_{\mathrm{cal}}(C_{\mathrm{deg}},d)
\le(C_{\mathrm{deg}}+2)L_d.
\]
Choose
\[
C_{\mathrm{sp}}=\max\{C_{\mathrm{tv}},2\},
\qquad
K_{\mathrm{sp}}=\mathcal K_{\mathrm{cal}}(C_{\mathrm{sp}},d).
\]
Because \(L_d^2/d\to0\), a universal integer \(D_{\mathrm{sp}}\ge2\) can be chosen so that, for \(d\ge D_{\mathrm{sp}}\),
\[
\frac{L_d^2}{d}
<
\frac{c_{\mathrm p}^2}{448(C_{\mathrm{sp}}+2)^2},
\qquad
448K_{\mathrm{sp}}^2\le c_{\mathrm p}^2d.
\]
The asserted logarithmic limit follows, for example, by writing
\(d=e^{L_d-1}\) and using that an exponential dominates a fixed power. For \(M\ge2\) and \(0\le\epsilon\le1/2\),
\[
256(2+3\epsilon^2M)\le448M.
\]
Hence this single cutoff ensures
\[
256(2+3\epsilon^2M)K_{\mathrm{sp}}^2
\le c_{\mathrm p}^2Md
\]
uniformly in the required support and overlap ranges. The previous two steps therefore apply whenever
\[
2\le M\le K_{\mathrm{sp}}^2,
\qquad
M\le\frac{\eta_{\mathrm{sp}}dK_{\mathrm{sp}}}{Bn\epsilon}.
\]
% lean: CausalSmith.Stat.OptvalueVanishingoverlapRate.sparse_packet_bayesRisk_lower_largeAlphabet

\item
\textbf{Dense comparison.}\quad
The dense construction uses symmetric moment-matched priors on \([-1,1]\). We establish the absolute-moment gap used for them.

For \(K_{\mathrm{abs}}\ge1\), define
\[
E_{\mathrm{abs}}(K_{\mathrm{abs}})
=\inf_{\deg P_{\mathrm{abs}}\le K_{\mathrm{abs}}}
\sup_{-1\le z\le1}\bigl||z|-P_{\mathrm{abs}}(z)\bigr|.
\]
On the circle of period \(\pi\), define, for each integer \(N_{\mathrm{Fej}}\ge1\),
\[
F_{\mathrm{Fej},N_{\mathrm{Fej}}}(\omega)
=\frac1{N_{\mathrm{Fej}}}
\left|\sum_{\ell_{\mathrm{Fej}}=0}^{N_{\mathrm{Fej}}-1}
 e^{2i\ell_{\mathrm{Fej}}\omega}\right|^2,
\qquad \omega\in\mathbb R,
\]
\[
V_{\mathrm{VP},K_{\mathrm{abs}}}
=2F_{\mathrm{Fej},2K_{\mathrm{abs}}}
-F_{\mathrm{Fej},K_{\mathrm{abs}}}.
\]
For a continuous complex-valued \(\pi\)-periodic function \(f_{\mathrm{cusp}}\), define
\[
\Lambda_{\mathrm{cusp},K_{\mathrm{abs}}}(f_{\mathrm{cusp}})
=f_{\mathrm{cusp}}(0)-\int_0^\pi
V_{\mathrm{VP},K_{\mathrm{abs}}}(\omega)
 f_{\mathrm{cusp}}(\omega)\,\frac{d\omega}{\pi}.
\]
Expanding the squared finite sum gives
\[
F_{\mathrm{Fej},N_{\mathrm{Fej}}}(\omega)
=\frac1{N_{\mathrm{Fej}}}
\sum_{\ell_{\mathrm{Fej}},r_{\mathrm{Fej}}=0}^{N_{\mathrm{Fej}}-1}
 e^{2i(\ell_{\mathrm{Fej}}-r_{\mathrm{Fej}})\omega}.
\]
For each integer frequency \(j_{\mathrm{Four}}\), orthogonality retains exactly the pairs whose index difference is \(-j_{\mathrm{Four}}\). There are
\(N_{\mathrm{Fej}}-|j_{\mathrm{Four}}|\) such pairs when
\(|j_{\mathrm{Four}}|<N_{\mathrm{Fej}}\), and none otherwise. Thus the triangular coefficients are
\[
\begin{aligned}
\gamma_{\mathrm{Fej}}(N_{\mathrm{Fej}},j_{\mathrm{Four}})
&:=\int_0^\pi F_{\mathrm{Fej},N_{\mathrm{Fej}}}(\omega)
 e^{2ij_{\mathrm{Four}}\omega}\,\frac{d\omega}{\pi}\\
&=\begin{cases}
1-|j_{\mathrm{Four}}|/N_{\mathrm{Fej}},
&|j_{\mathrm{Four}}|<N_{\mathrm{Fej}},\\
0,&|j_{\mathrm{Four}}|\ge N_{\mathrm{Fej}}.
\end{cases}
\end{aligned}
\]
In particular, both Fej\'er kernels are nonnegative with integral one, so
\[
|\Lambda_{\mathrm{cusp},K_{\mathrm{abs}}}(f_{\mathrm{cusp}})|
\le4\|f_{\mathrm{cusp}}\|_\infty.
\]
Substituting the triangular coefficients into the cusp functional gives
\[
\begin{aligned}
\mu_{\mathrm{cusp}}(j_{\mathrm{Four}})
&:=\Lambda_{\mathrm{cusp},K_{\mathrm{abs}}}
 (\omega\mapsto e^{2ij_{\mathrm{Four}}\omega})\\
&=1-2\gamma_{\mathrm{Fej}}(2K_{\mathrm{abs}},j_{\mathrm{Four}})
 +\gamma_{\mathrm{Fej}}(K_{\mathrm{abs}},j_{\mathrm{Four}})\\
&=\begin{cases}
0,&|j_{\mathrm{Four}}|\le K_{\mathrm{abs}},\\
|j_{\mathrm{Four}}|/K_{\mathrm{abs}}-1,
&K_{\mathrm{abs}}<|j_{\mathrm{Four}}|<2K_{\mathrm{abs}},\\
1,&|j_{\mathrm{Four}}|\ge2K_{\mathrm{abs}}.
\end{cases}
\end{aligned}
\]
These multipliers are nonnegative. The identity
\[
\sin2\omega=\frac{e^{2i\omega}-e^{-2i\omega}}{2i}
\]
and the binomial expansion show that a polynomial of degree at most
\(K_{\mathrm{abs}}\) in \(\sin2\omega\) has frequencies
\(|j_{\mathrm{Four}}|\le K_{\mathrm{abs}}\). Hence
\[
\Lambda_{\mathrm{cusp},K_{\mathrm{abs}}}
 (\omega\mapsto P_{\mathrm{abs}}(\sin2\omega))=0
\qquad(\deg P_{\mathrm{abs}}\le K_{\mathrm{abs}}).
\]

To compute the absolute-sine coefficients, define
\[
S_{\mathrm{abs}}(\omega)=|\sin\omega|,
\qquad
c_{\mathrm{Four},\mathrm{abs}}(j_{\mathrm{Four}})
=\int_0^\pi S_{\mathrm{abs}}(\omega)e^{-2ij_{\mathrm{Four}}\omega}
 \,\frac{d\omega}{\pi},
\qquad j_{\mathrm{Four}}\in\mathbb Z.
\]
On \([0,\pi]\), \(S_{\mathrm{abs}}(\omega)=\sin\omega\). For every integer
\(j_{\mathrm{Four}}\), the denominator \(1-4j_{\mathrm{Four}}^2\) is nonzero, and the function
\[
\mathcal A_{\mathrm{Four},j_{\mathrm{Four}}}(\omega)
=e^{-2ij_{\mathrm{Four}}\omega}
\frac{-2ij_{\mathrm{Four}}\sin\omega-\cos\omega}
 {1-4j_{\mathrm{Four}}^2},
\qquad \omega\in[0,\pi],
\]
satisfies
\[
\mathcal A_{\mathrm{Four},j_{\mathrm{Four}}}'(\omega)
=e^{-2ij_{\mathrm{Four}}\omega}\sin\omega.
\]
Evaluating at the endpoints, using
\(e^{-2ij_{\mathrm{Four}}\pi}=1\), gives
\[
c_{\mathrm{Four},\mathrm{abs}}(j_{\mathrm{Four}})
=\frac{\mathcal A_{\mathrm{Four},j_{\mathrm{Four}}}(\pi)
 -\mathcal A_{\mathrm{Four},j_{\mathrm{Four}}}(0)}{\pi}
=-\frac2{\pi(4j_{\mathrm{Four}}^2-1)}.
\]
The zero coefficient is \(2/\pi\). For \(|j_{\mathrm{Four}}|\ge1\),
\[
|c_{\mathrm{Four},\mathrm{abs}}(j_{\mathrm{Four}})|
=\frac2{\pi(4j_{\mathrm{Four}}^2-1)}
\le\frac1{j_{\mathrm{Four}}^2},
\qquad
\sum_{j_{\mathrm{Four}}\in\mathbb Z}
 |c_{\mathrm{Four},\mathrm{abs}}(j_{\mathrm{Four}})|<\infty.
\]
Thus the Fourier series converges absolutely and uniformly; its coefficients identify its sum with the continuous function \(S_{\mathrm{abs}}\). Pairing opposite frequencies and replacing \(\omega\) by \(2\omega\) yields
\[
|\sin2\omega|
=\frac2\pi-\frac4\pi
 \sum_{\ell=1}^\infty\frac{\cos(4\ell\omega)}{4\ell^2-1}.
\]
Continuity of the cusp functional permits termwise application, giving
\[
-\Lambda_{\mathrm{cusp},K_{\mathrm{abs}}}(|\sin2\cdot|)
=\sum_{\ell=1}^\infty
 \frac{4\mu_{\mathrm{cusp}}(2\ell)}{\pi(4\ell^2-1)}.
\]
Every summand is nonnegative. For
\(K_{\mathrm{abs}}\le\ell<2K_{\mathrm{abs}}\), the multiplier equals one, and \(\pi\le4\) gives
\[
\frac4{\pi(4\ell^2-1)}\ge\frac1{16K_{\mathrm{abs}}^2}.
\] Keeping the
\(K_{\mathrm{abs}}\) terms with
\(K_{\mathrm{abs}}\le\ell<2K_{\mathrm{abs}}\) gives
\[
-\Lambda_{\mathrm{cusp},K_{\mathrm{abs}}}(|\sin2\cdot|)
\ge
\sum_{\ell=K_{\mathrm{abs}}}^{2K_{\mathrm{abs}}-1}
\frac4{\pi(4\ell^2-1)}
\ge\frac1{16K_{\mathrm{abs}}}.
\]
It follows that
\[
\frac1{16K_{\mathrm{abs}}}
\le
4\sup_{-1\le z\le1}\bigl||z|-P_{\mathrm{abs}}(z)\bigr|,
\]
and therefore
\[
E_{\mathrm{abs}}(K_{\mathrm{abs}})
\ge\frac1{100K_{\mathrm{abs}}}.
\]

To construct the priors, let
\[
\mathcal V_{\mathrm{abs}}
=\{P_{\mathrm{abs}}|_{[-1,1]}:\deg P_{\mathrm{abs}}\le K_{\mathrm{abs}}\}.
\]
The norm of the class of \(|z|\) in
\(C([-1,1],\mathbb R)/\mathcal V_{\mathrm{abs}}\) is
\(E_{\mathrm{abs}}(K_{\mathrm{abs}})\). A norming functional on that quotient, pulled back to \(C([-1,1],\mathbb R)\), supplies a functional
\(\Lambda_{\mathrm{abs}}\) satisfying
\[
\|\Lambda_{\mathrm{abs}}\|\le1,
\qquad
\Lambda_{\mathrm{abs}}(z^\ell)=0
\quad(0\le\ell\le K_{\mathrm{abs}}),
\qquad
\Lambda_{\mathrm{abs}}(|z|)=E_{\mathrm{abs}}(K_{\mathrm{abs}}).
\]
Its positive half-split can be obtained by positive extension. In detail, let \(1\) denote the constant function with value one and consider the cone
\[
\mathcal C_{\mathrm{split}}
=\{(f_++f_-,-\Lambda_{\mathrm{abs}}(f_-)):
 f_+,f_-\in C([-1,1],\mathbb R),\ f_+,f_-\ge0\}.
\]
On the subspace of pairs \((a_{\mathrm{cone}}\,1,b_{\mathrm{cone}})\), prescribe the functional
\[
(a_{\mathrm{cone}}\,1,b_{\mathrm{cone}})
\longmapsto a_{\mathrm{cone}}/2+b_{\mathrm{cone}},
\qquad a_{\mathrm{cone}},b_{\mathrm{cone}}\in\mathbb R.
\]
If such a pair belongs to the cone, then
\[
f_++f_-=a_{\mathrm{cone}}\,1,
\qquad
b_{\mathrm{cone}}=-\Lambda_{\mathrm{abs}}(f_-),
\qquad a_{\mathrm{cone}}\ge0.
\]
Consequently,
\[
\|2f_--a_{\mathrm{cone}}\,1\|_\infty\le a_{\mathrm{cone}},
\qquad
\Lambda_{\mathrm{abs}}(f_-)
=\tfrac12\Lambda_{\mathrm{abs}}(2f_--a_{\mathrm{cone}}\,1)
\le a_{\mathrm{cone}}/2,
\]
where the last bound uses \(\|\Lambda_{\mathrm{abs}}\|\le1\) and
\(\Lambda_{\mathrm{abs}}(1)=0\). Thus the prescribed functional is nonnegative on its intersection with the cone.

The subspace also satisfies the required majorization condition. For every
\(f_{\mathrm{split}}\in C([-1,1],\mathbb R)\) and
\(t_{\mathrm{split}}\in\mathbb R\),
\[
(\|f_{\mathrm{split}}\|_\infty\,1,-t_{\mathrm{split}})
 +(f_{\mathrm{split}},t_{\mathrm{split}})
=(\|f_{\mathrm{split}}\|_\infty\,1+f_{\mathrm{split}},0)
\in\mathcal C_{\mathrm{split}}.
\]
Indeed, \(\|f_{\mathrm{split}}\|_\infty\,1+f_{\mathrm{split}}\ge0\), so this is the cone element obtained with that function as \(f_+\) and zero as \(f_-\). Positive extension therefore supplies a linear functional
\[
G_{\mathrm{split}}:C([-1,1],\mathbb R)\times\mathbb R\longrightarrow\mathbb R
\]
that is nonnegative on \(\mathcal C_{\mathrm{split}}\) and satisfies
\[
G_{\mathrm{split}}(a_{\mathrm{cone}}\,1,b_{\mathrm{cone}})
=a_{\mathrm{cone}}/2+b_{\mathrm{cone}}.
\]
Define, for every \(f_{\mathrm{split}}\in C([-1,1],\mathbb R)\),
\[
\Lambda_+(f_{\mathrm{split}})=G_{\mathrm{split}}(f_{\mathrm{split}},0),
\qquad
\Lambda_-(f_{\mathrm{split}})
=\Lambda_+(f_{\mathrm{split}})-\Lambda_{\mathrm{abs}}(f_{\mathrm{split}}).
\]
For \(f_{\mathrm{split}}\ge0\), both
\((f_{\mathrm{split}},0)\) and
\((f_{\mathrm{split}},-\Lambda_{\mathrm{abs}}(f_{\mathrm{split}}))\) belong to the cone. Since
\(G_{\mathrm{split}}(0,t_{\mathrm{split}})=t_{\mathrm{split}}\), this gives
\[
\Lambda_+(f_{\mathrm{split}})\ge0,
\qquad
\Lambda_-(f_{\mathrm{split}})
=G_{\mathrm{split}}(f_{\mathrm{split}},-\Lambda_{\mathrm{abs}}(f_{\mathrm{split}}))\ge0.
\]
Their masses and difference are
\[
\Lambda_+(1)=G_{\mathrm{split}}(1,0)=\frac12,
\qquad
\Lambda_-(1)=\Lambda_+(1)-\Lambda_{\mathrm{abs}}(1)=\frac12,
\qquad
\Lambda_+-\Lambda_-=\Lambda_{\mathrm{abs}}.
\]
Positivity bounds each functional on a continuous function by its mass times the uniform norm, so both have representing finite measures.
Their representing measures have mass \(1/2\). Averaging each measure with its reflection under \(z\mapsto-z\) preserves the absolute-value pairing and the annihilation of all powers through degree \(K_{\mathrm{abs}}\). Doubling these symmetric measures gives symmetric probability measures
\(\eta_0,\eta_1\) on \([-1,1]\) satisfying
\[
\int z^\ell\,d\eta_0=\int z^\ell\,d\eta_1
\quad(0\le\ell\le K_{\mathrm{abs}}),
\]
\[
\int|z|\,d\eta_1-\int|z|\,d\eta_0
=2E_{\mathrm{abs}}(K_{\mathrm{abs}})
\ge\frac1{50K_{\mathrm{abs}}}.
\]
% lean: CausalSmith.Stat.OptvalueVanishingoverlapRate.dense_symmetric_priors_absGap_lower

\item
For \(0<\mathfrak a\le1/2\), define the independent prior draws by
\[
\mathbf T^{\mathrm{de}}
=(T_x^{\mathrm{de}})_{x\in[d]}
\sim\eta_\iota^{\otimes d},
\qquad
\mathbf T^{\mathrm{de}}\in[-1,1]^d,
\]
and define the dense observed law by
\[
p_x=\frac1d,
\qquad
\pi_x=\epsilon,
\qquad
\mu_{0x}=\frac12,
\qquad
\mu_{1x}=\frac12+\mathfrak a T_x^{\mathrm{de}}.
\]
Equivalently, its four atom probabilities are
\[
q_{00,x}=q_{01,x}=\frac{1-\epsilon}{2d},
\qquad
q_{10,x}=\frac{\epsilon}{d}\left(\frac12-\mathfrak a T_x^{\mathrm{de}}\right),
\qquad
q_{11,x}=\frac{\epsilon}{d}\left(\frac12+\mathfrak a T_x^{\mathrm{de}}\right).
\]
These are legal observed laws, including the endpoints \(T_x^{\mathrm{de}}=\pm1\).
Their target and prior centers are
\[
\Psi_{\mathrm{de}}(\mathbf T^{\mathrm{de}})
=\frac12+\frac{\mathfrak a}{2d}
 \sum_{x\in[d]}(|T_x^{\mathrm{de}}|+T_x^{\mathrm{de}}),
\]
\[
\theta_{\mathrm{de},\iota}
=\frac12+\frac{\mathfrak a}{2}\int|z|\,d\eta_\iota,
\qquad
\delta_{\mathrm{de}}
=\frac{\mathfrak a}{2}
 \left(\int|z|\,d\eta_1-\int|z|\,d\eta_0\right)
\ge\frac{\mathfrak a}{100K_{\mathrm{abs}}}.
\]
Symmetry gives \(\int z\,d\eta_\iota=0\). Since
\((|z|+z)/2=z_+\in[0,1]\), independence gives
\[
E_\iota\{\Psi_{\mathrm{de}}-\theta_{\mathrm{de},\iota}\}^2
\le\frac{\mathfrak a^2}{d}.
\]

In the Poisson experiment, define the total treated intensity per cell by
\[
\lambda_{\mathrm{tr}}=\frac{Bn\epsilon}{d}.
\]
The treated counts have means
\[
\lambda_{\mathrm{tr}}\left(\frac12-\mathfrak a z\right),
\qquad
\lambda_{\mathrm{tr}}\left(\frac12+\mathfrak a z\right).
\]
For nonnegative integer treated failure and success counts, respectively denoted by \(N_{\mathrm{de},0}\) and \(N_{\mathrm{de},1}\), define the likelihood ratio relative to \(z=0\) by
\[
\begin{aligned}
L_{\mathrm{de},z}(N_{\mathrm{de},0},N_{\mathrm{de},1})
&=e^{\lambda_{\mathrm{tr}}\mathfrak a z}
 (1-2\mathfrak a z)^{N_{\mathrm{de},0}}
 e^{-\lambda_{\mathrm{tr}}\mathfrak a z}
 (1+2\mathfrak a z)^{N_{\mathrm{de},1}}\\
&=(1-2\mathfrak a z)^{N_{\mathrm{de},0}}
  (1+2\mathfrak a z)^{N_{\mathrm{de},1}},
\qquad z\in[-1,1].
\end{aligned}
\]
At a zero alternative mean, use \(0^0=1\); a zero base raised to a positive count is zero. The exponential factors cancel because the sum of the two treated means is \(\lambda_{\mathrm{tr}}\). Multiplication of these likelihood factors and their expectations gives the Gram identity
\[
E_0[L_{\mathrm{de},z}L_{\mathrm{de},w}]
=\exp\{4\lambda_{\mathrm{tr}}\mathfrak a^2zw\},
\qquad z,w\in[-1,1].
\]
Thus, defining
\[
\tau_{\mathrm{de}}=4\lambda_{\mathrm{tr}}\mathfrak a^2,
\qquad
\eta_{\mathrm{de}}=\frac{\beta_{\mathrm{tail}}}{4},
\]
the same moment-matching and Cauchy--Schwarz argument gives predictive total variation at most
\[
d\sqrt{\mathcal T_{K_{\mathrm{abs}}}(\tau_{\mathrm{de}})}
\le\frac18
\]
whenever
\[
K_{\mathrm{abs}}\ge C_{\mathrm{tv}}L_d,
\qquad
\lambda_{\mathrm{tr}}\mathfrak a^2
\le\eta_{\mathrm{de}}K_{\mathrm{abs}}.
\]
The two untreated counts have common means
\(Bn(1-\epsilon)/(2d)\), independent of the prior parameter. Adding them preserves this total-variation bound, so it applies to all four observed counts.

If
\[
128\mathfrak a^2\le d\delta_{\mathrm{de}}^2,
\]
the target-deviation probability at radius
\(\delta_{\mathrm{de}}/4\) is at most \(1/8\). Applying the midpoint argument above to the two dense centers gives, for every measurable estimator of the full counts,
\[
\max_{\iota=0,1}\mathcal B_{\mathrm{de},\iota}(\widehat W)
\ge\frac{5\delta_{\mathrm{de}}^2}{256}
\ge\frac5{2560000}\frac{\mathfrak a^2}{K_{\mathrm{abs}}^2},
\]
where
\[
\mathcal B_{\mathrm{de},\iota}(\widehat W)
=
E_{\mathbf T^{\mathrm{de}}\sim\eta_\iota^{\otimes d}}
E_{\mathrm{counts}\mid\mathbf T^{\mathrm{de}}}
[\{\widehat W-\Psi_{\mathrm{de}}(\mathbf T^{\mathrm{de}})\}^2].
\]
Choose
\[
C_{\mathrm{de}}=\max\{C_{\mathrm{tv}},2\},
\qquad
K_{\mathrm{de}}=\mathcal K_{\mathrm{cal}}(C_{\mathrm{de}},d).
\]
The same logarithmic cutoff construction, applied with packet comparison constant \(1/100\), supplies a universal
\(D_{\mathrm{de}}\ge2\) such that
\[
1280000K_{\mathrm{de}}^2\le d
\quad(d\ge D_{\mathrm{de}}).
\]
Since
\(\delta_{\mathrm{de}}^2\ge\mathfrak a^2/(10000K_{\mathrm{de}}^2)\),
this implies the required target-deviation condition. Reordering the four counts into the observed atom histogram is a bijection, so both integrated-risk bounds hold on that histogram space.
% lean: CausalSmith.Stat.OptvalueVanishingoverlapRate.dense_full_observed_logDegree_bayesRisk_lower

\item
We transfer these bounds to fixed samples. Let \(\widehat V\) be any measurable estimator based on \(n\) observations. For a histogram
\[
\boldsymbol h_{\mathrm{hist}}:\mathcal J_d\to\mathbb N,
\qquad
n_{\mathrm{hist}}=\sum_{o\in\mathcal J_d}\boldsymbol h_{\mathrm{hist}}(o),
\]
define
\[
\mathcal S(\boldsymbol h_{\mathrm{hist}})
=\{\boldsymbol s\in\mathcal J_d^{n_{\mathrm{hist}}}:
\boldsymbol s\text{ has histogram }\boldsymbol h_{\mathrm{hist}}\},
\qquad
\boldsymbol o^\circ=((0,0,0),\ldots,(0,0,0))\in\mathcal J_d^n,
\]
\[
\boldsymbol o_n(\boldsymbol s)=
\begin{cases}
(s_1,\ldots,s_n),&n_{\mathrm{hist}}\ge n,\\
\boldsymbol o^\circ,&n_{\mathrm{hist}}<n,
\end{cases}
\]
\[
\widehat W_{\widehat V}(\boldsymbol h_{\mathrm{hist}})
=\frac1{|\mathcal S(\boldsymbol h_{\mathrm{hist}})|}
 \sum_{\boldsymbol s\in\mathcal S(\boldsymbol h_{\mathrm{hist}})}
 \Pi_{[0,1]}\{\widehat V(\boldsymbol o_n(\boldsymbol s))\}.
\]
The set \(\mathcal S(\boldsymbol h_{\mathrm{hist}})\) is nonempty, including the zero histogram, for which it contains the empty string.

For independent Poisson atom counts with total mean
\(\lambda_{\mathrm{tot}}\), conditional on their total, the labels are independent with the normalized observed law; conditional on the histogram, their ordering is uniform on \(\mathcal S(h)\). Jensen's inequality, followed by projection onto \([0,1]\), therefore gives
\[
E\{\widehat W_{\widehat V}-\Psi\}^2
\le E\{\widehat V-\Psi\}^2
+\Pr\{\operatorname{Pois}(\lambda_{\mathrm{tot}})<n\}.
\]
For the sparse experiment,
\[
\lambda_{\mathrm{tot}}=BnS_{\mathrm{sp}}\ge Bn/2;
\]
for the dense experiment, \(\lambda_{\mathrm{tot}}=Bn\).

For any \(\lambda_{\mathrm{tot}}\ge Bn/2>n\), put
\[
\gamma_{\mathrm{sh}}
=\frac{\lambda_{\mathrm{tot}}-n}{\lambda_{\mathrm{tot}}}\in(0,1).
\]
Exponential Markov inequality and
\(e^{-\gamma_{\mathrm{sh}}}
\le1-\gamma_{\mathrm{sh}}+\gamma_{\mathrm{sh}}^2/2\) give
\[
\Pr\{\operatorname{Pois}(\lambda_{\mathrm{tot}})<n\}
\le
\exp\left\{-\frac{(\lambda_{\mathrm{tot}}-n)^2}
                       {2\lambda_{\mathrm{tot}}}\right\}
\le
\exp\left\{-\frac{n(B-2)^2}{4B}\right\}.
\]
The last inequality uses that
\((\lambda_{\mathrm{tot}}-n)^2/(2\lambda_{\mathrm{tot}})\)
is increasing for \(\lambda_{\mathrm{tot}}>n\).

Define the shortage penalty by
\[
p_{\mathrm{sh}}(n,B)
=\exp\left\{-\frac{n(B-2)^2}{4B}\right\}.
\]
Integrating the reconstruction bound under either supported prior and then taking the fixed-sample estimator infimum proves
\[
c_{\mathrm{sp}}\frac M{K_{\mathrm{sp}}^2}
-p_{\mathrm{sh}}(n,B)
\le\mathfrak R_{n,d,\epsilon}
\]
under the sparse conditions, and
\[
\frac5{2560000}\frac{\mathfrak a^2}{K_{\mathrm{de}}^2}
-p_{\mathrm{sh}}(n,B)
\le\mathfrak R_{n,d,\epsilon}
\]
under the dense conditions.

For the sparse bound choose
\[
M_{\mathrm{sp}}
=\min\left\{
K_{\mathrm{sp}}^2,
\frac{\eta_{\mathrm{sp}}dK_{\mathrm{sp}}}{Bn\epsilon}
\right\}.
\]
Whenever the second entry is at least two, this is an admissible support endpoint, and
\[
c_{\mathrm{sp}}
\min\left\{1,\frac{\eta_{\mathrm{sp}}d}
                    {Bn\epsilon K_{\mathrm{sp}}}\right\}
-p_{\mathrm{sh}}(n,B)
\le\mathfrak R_{n,d,\epsilon}.
\]
For the dense bound choose
\[
\mathfrak a_{\mathrm{de}}^2
=\min\left\{\frac14,
\frac{\eta_{\mathrm{de}}dK_{\mathrm{de}}}{Bn\epsilon}\right\}.
\]
Then \(0<\mathfrak a_{\mathrm{de}}\le1/2\) and the likelihood condition holds. Dividing by \(K_{\mathrm{de}}^2\) gives
\[
\frac5{2560000}
\min\left\{
\frac1{4K_{\mathrm{de}}^2},
\frac{\eta_{\mathrm{de}}d}{Bn\epsilon K_{\mathrm{de}}}
\right\}
-p_{\mathrm{sh}}(n,B)
\le\mathfrak R_{n,d,\epsilon}.
\]
% lean: CausalSmith.Stat.OptvalueVanishingoverlapRate.sparse_fixedSample_calibrated_lower

\item
\textbf{Fixed-sample transfer.}\quad
We join these two ranges while retaining their degree comparison. Define
\[
D_{\mathrm{lg}}=\max\{D_{\mathrm{sp}},D_{\mathrm{de}}\},
\qquad
q_{\mathrm{deg}}=\frac{C_{\mathrm{de}}+2}{C_{\mathrm{sp}}},
\]
\[
b_{\mathrm{splice}}
=\min\left\{
\frac{\eta_{\mathrm{sp}}}{8q_{\mathrm{deg}}^2},
\frac{\eta_{\mathrm{de}}}{q_{\mathrm{deg}}}
\right\},
\qquad
c_{\mathrm{splice}}
=\min\left\{
c_{\mathrm{sp}}\min\{1,\eta_{\mathrm{sp}}\},
\frac5{2560000}b_{\mathrm{splice}}
\right\}>0,
\]
\[
x_{\mathrm{splice}}=\frac d{Bn\epsilon}.
\]
The degree bounds imply
\[
K_{\mathrm{de}}\le q_{\mathrm{deg}}K_{\mathrm{sp}}.
\]

If \(\eta_{\mathrm{sp}}x_{\mathrm{splice}}K_{\mathrm{sp}}\ge2\), use the sparse bound. The elementary inequality
\[
\min\{1,\eta_{\mathrm{sp}}\}
\min\left\{1,\frac{x_{\mathrm{splice}}}{K_{\mathrm{sp}}}\right\}
\le
\min\left\{1,
\frac{\eta_{\mathrm{sp}}x_{\mathrm{splice}}}{K_{\mathrm{sp}}}\right\}
\]
shows that its leading term is at least
\(c_{\mathrm{splice}}\min\{1,x_{\mathrm{splice}}/K_{\mathrm{sp}}\}\).

Otherwise,
\(\eta_{\mathrm{sp}}x_{\mathrm{splice}}K_{\mathrm{sp}}\le2\).
The degree comparison and the definition of \(b_{\mathrm{splice}}\) give
\[
b_{\mathrm{splice}}
\frac{x_{\mathrm{splice}}}{K_{\mathrm{sp}}}
\le\frac1{4K_{\mathrm{de}}^2},
\qquad
b_{\mathrm{splice}}
\frac{x_{\mathrm{splice}}}{K_{\mathrm{sp}}}
\le\frac{\eta_{\mathrm{de}}x_{\mathrm{splice}}}{K_{\mathrm{de}}}.
\]
For the first inequality, multiply by \(4K_{\mathrm{de}}^2\) and use
\[
\frac{\eta_{\mathrm{sp}}}{8q_{\mathrm{deg}}^2}
\frac{x_{\mathrm{splice}}}{K_{\mathrm{sp}}}
\,4K_{\mathrm{de}}^2
\le\frac{\eta_{\mathrm{sp}}x_{\mathrm{splice}}K_{\mathrm{sp}}}{2}
\le1.
\]
For the second, use
\(K_{\mathrm{de}}\le q_{\mathrm{deg}}K_{\mathrm{sp}}\).
The dense bound therefore supplies the same leading term. In both cases,
\[
c_{\mathrm{splice}}
\min\left\{1,\frac d{Bn\epsilon K_{\mathrm{sp}}}\right\}
-p_{\mathrm{sh}}(n,B)
\le\mathfrak R_{n,d,\epsilon}
\qquad(d\ge D_{\mathrm{lg}}).
\]

Set
\[
c_{\mathrm{sh}}=\frac{c_{\mathrm{splice}}}{C_{\mathrm{sp}}+2}.
\]
Since \(K_{\mathrm{sp}}\le(C_{\mathrm{sp}}+2)L_d\) and
\(1/\{B(C_{\mathrm{sp}}+2)\}\le1\),
\[
\frac1{B(C_{\mathrm{sp}}+2)}\rho_{n,d,\epsilon}
\le
\min\left\{1,\frac d{Bn\epsilon K_{\mathrm{sp}}}\right\}.
\]
Thus
\[
\frac{c_{\mathrm{sh}}}{B}\rho_{n,d,\epsilon}
-p_{\mathrm{sh}}(n,B)
\le\mathfrak R_{n,d,\epsilon}
\qquad(d\ge D_{\mathrm{lg}},\ B>2).
\]
% lean: CausalSmith.Stat.OptvalueVanishingoverlapRate.largeAlphabet_lower_with_shortage

\item
Choose once and for all
\[
B_{\mathrm{lg}}=\max\{4,2048/c_{\mathrm{sh}}\},
\qquad
c_{\mathrm{lg}}=\frac{c_{\mathrm{sh}}}{2B_{\mathrm{lg}}}>0.
\]
For \(B\ge4\), define
\[
a_{\mathrm{sh}}=\frac{(B-2)^2}{8B}\ge\frac B{32}.
\]
The inequality \(e^{-v}\le1/v\) for \(v>0\) gives
\[
e^{-a_{\mathrm{sh}}n}\le\frac{32}{Bn},
\qquad
e^{-a_{\mathrm{sh}}n}\le e^{-a_{\mathrm{sh}}}\le\frac{32}{B}.
\]
Multiplying these bounds proves
\[
p_{\mathrm{sh}}(n,B)
=e^{-a_{\mathrm{sh}}n}e^{-a_{\mathrm{sh}}n}
\le\frac{1024}{B^2n}.
\]
Also \(\log d\le d-1\), so \(L_d\le d\). Since \(n\ge1\) and
\(\epsilon\le1/2\le1\),
\[
\frac1n\le1,
\qquad
\frac1n\le\frac d{n\epsilon L_d},
\qquad
\frac1n\le\rho_{n,d,\epsilon}.
\]
Because \(B_{\mathrm{lg}}c_{\mathrm{sh}}\ge2048\),
\[
p_{\mathrm{sh}}(n,B_{\mathrm{lg}})
\le\frac{1024}{B_{\mathrm{lg}}^2n}
\le\frac{c_{\mathrm{sh}}}{2B_{\mathrm{lg}}}\frac1n
\le c_{\mathrm{lg}}\rho_{n,d,\epsilon}.
\]
The preceding lower bound has leading coefficient
\(c_{\mathrm{sh}}/B_{\mathrm{lg}}=2c_{\mathrm{lg}}\). Subtracting the shortage penalty therefore gives
\[
c_{\mathrm{lg}}\rho_{n,d,\epsilon}
\le\mathfrak R_{n,d,\epsilon}
\qquad(d\ge D_{\mathrm{lg}}).
\]
% lean: CausalSmith.Stat.OptvalueVanishingoverlapRate.largeAlphabet_minimaxRisk_lower

\item
For the bounded-alphabet comparison, define
\[
r_{\mathrm{par}}=\min\left\{1,\frac1{n\epsilon}\right\},
\qquad
a_{\mathrm{par}}=\sqrt{\frac{r_{\mathrm{par}}}{64}}.
\]
Then
\[
0<a_{\mathrm{par}}\le\frac12,
\qquad
n\epsilon a_{\mathrm{par}}^2\le\frac1{16}.
\]
Use two dense observed laws \(\mathbb P_{\mathrm{par},1}\) and
\(\mathbb P_{\mathrm{par},0}\) with uniform covariate probabilities, propensity
\(\epsilon\), control mean \(1/2\), and treated means respectively
\(1/2+a_{\mathrm{par}}\) and \(1/2\). Define their scalar targets by
\[
\theta_{\mathrm{par},1}
=\Psi(\mathbb P_{\mathrm{par},1})=\frac12+a_{\mathrm{par}},
\qquad
\theta_{\mathrm{par},0}
=\Psi(\mathbb P_{\mathrm{par},0})=\frac12.
\]
For every \(x\in[d]\), the atom probabilities of these laws are
\[
\begin{aligned}
\mathbb P_{\mathrm{par},1}(x,0,0)
&=\mathbb P_{\mathrm{par},1}(x,0,1)
=\mathbb P_{\mathrm{par},0}(x,0,0)
=\mathbb P_{\mathrm{par},0}(x,0,1)
=\frac{1-\epsilon}{2d},\\
\mathbb P_{\mathrm{par},0}(x,1,0)
&=\mathbb P_{\mathrm{par},0}(x,1,1)=\frac{\epsilon}{2d},\\
\mathbb P_{\mathrm{par},1}(x,1,0)
&=\frac{\epsilon}{d}\left(\frac12-a_{\mathrm{par}}\right),
\qquad
\mathbb P_{\mathrm{par},1}(x,1,1)
=\frac{\epsilon}{d}\left(\frac12+a_{\mathrm{par}}\right).
\end{aligned}
\]
All reference atom probabilities are positive. Each untreated atom has zero probability difference and hence zero chi-square contribution. Each treated atom contributes
\[
\frac{(\epsilon a_{\mathrm{par}}/d)^2}{\epsilon/(2d)}
=\frac{2\epsilon a_{\mathrm{par}}^2}{d}.
\]
Summing over both treated outcomes and all \(d\) cells gives the finite-atom calculation
\[
\begin{aligned}
\chi^2(\mathbb P_{\mathrm{par},1},\mathbb P_{\mathrm{par},0})
&=\sum_{x\in[d]}\sum_{a,y\in\{0,1\}}
\frac{\{\mathbb P_{\mathrm{par},1}(x,a,y)
-\mathbb P_{\mathrm{par},0}(x,a,y)\}^2}
 {\mathbb P_{\mathrm{par},0}(x,a,y)}\\
&=d\cdot2\cdot\frac{2\epsilon a_{\mathrm{par}}^2}{d}
=4\epsilon a_{\mathrm{par}}^2.
\end{aligned}
\]
Define the two \(n\)-observation laws by
\[
\mathbb Q_{\mathrm{par},\iota}
=\mathbb P_{\mathrm{par},\iota}^{\otimes n},
\qquad \iota\in\{0,1\}.
\]
Product likelihoods give
\[
1+\chi^2(\mathbb Q_{\mathrm{par},1},\mathbb Q_{\mathrm{par},0})
=(1+4\epsilon a_{\mathrm{par}}^2)^n
\le e^{4n\epsilon a_{\mathrm{par}}^2}
\le e^{1/4}\le\frac43.
\]
In particular their chi-square divergence is at most one, so their total variation is at most \(1/2\). Thresholding an arbitrary estimator at
\(1/2+a_{\mathrm{par}}/2\) shows that one of the two laws has probability at least \(1/4\) of absolute error at least \(a_{\mathrm{par}}/2\). Hence
\[
\max_{\iota=0,1}
E_{\mathbb Q_{\mathrm{par},\iota}}
\{\widehat V-\theta_{\mathrm{par},\iota}\}^2
\ge\frac{a_{\mathrm{par}}^2}{16}
=\frac1{1024}r_{\mathrm{par}}.
\]
Taking the estimator infimum gives
\[
\frac1{1024}\min\left\{1,\frac1{n\epsilon}\right\}
\le\mathfrak R_{n,d,\epsilon}
\qquad(d\ge2).
\]

Now define
\[
c_{\mathrm{all}}
=\min\left\{c_{\mathrm{lg}},\frac1{1024D_{\mathrm{lg}}}\right\}>0.
\]
For \(d\ge D_{\mathrm{lg}}\), the large-alphabet bound applies.
For \(2\le d<D_{\mathrm{lg}}\), using \(L_d\ge1\) gives
\[
\rho_{n,d,\epsilon}
\le
D_{\mathrm{lg}}\min\left\{1,\frac1{n\epsilon}\right\}.
\]
Indeed, \(\rho_{n,d,\epsilon}\le1\le D_{\mathrm{lg}}\), and
\[
\rho_{n,d,\epsilon}
\le\frac d{n\epsilon L_d}
\le\frac{D_{\mathrm{lg}}}{n\epsilon}.
\]
Thus in this second case,
\[
c_{\mathrm{all}}\rho_{n,d,\epsilon}
\le
\frac1{1024}\min\left\{1,\frac1{n\epsilon}\right\}
\le\mathfrak R_{n,d,\epsilon}.
\]
Together the two cases establish
\[
c_{\mathrm{all}}\rho_{n,d,\epsilon}
\le\mathfrak R_{n,d,\epsilon}
\qquad(n\ge1,\ d\ge2,\ 0<\epsilon\le1/2).
\]
% lean: CausalSmith.Stat.OptvalueVanishingoverlapRate.allAlphabet_minimaxRisk_lower

\item
\textbf{Final consequences.}\quad
By \cref{obj:thm:observable-upper}, there are universal calibrations
\(H_0>0,\kappa>0,D_0\ge2\) and a universal \(C_{\mathrm u}>0\) for the measurable estimator in
\cref{obj:def:armwise-estimator}. Define its worst-case observed risk by
\[
\mathcal R^{\mathrm{AF}}_{n,d,\epsilon}
=
\sup_{\mathbb P\in\mathcal D_{d,\epsilon}^{\mathrm{obs}}}
E_{\mathbb P}
\left[
\{\widehat V_{n,d,\epsilon}^{\mathrm{AF}}-\Psi(\mathbb P)\}^2
\right].
\]
Apply the cited bound to the product sampling law
\(\mathbb P^{\otimes n}\), which satisfies
\cref{obj:ass:iid-sampling}. This gives
\[
\mathcal R^{\mathrm{AF}}_{n,d,\epsilon}
\le C_{\mathrm u}\rho_{n,d,\epsilon}.
\]
Set
\[
c_\star=\min\{c_{\mathrm{all}},C_{\mathrm u}\},
\qquad
C_\star=C_{\mathrm u}.
\]
Then \(0<c_\star\le C_\star<\infty\). Nonnegativity of
\(\rho_{n,d,\epsilon}\) permits reducing the lower constant, and the estimator infimum in
\cref{obj:def:minimax-risk} is bounded by the risk of this particular estimator. Consequently,
\[
c_\star\rho_{n,d,\epsilon}
\le\mathfrak R_{n,d,\epsilon}
\le\mathcal R^{\mathrm{AF}}_{n,d,\epsilon}
\le C_\star\rho_{n,d,\epsilon}.
\]

For each measurable \(\widehat V\), identification and completion in
\cref{obj:prop:identification} give the equality
\[
\begin{aligned}
&\sup_{P\in\mathcal P_{d,\epsilon}}
E_{\operatorname{Obs}(P)}
[\{\widehat V-V^\star(P)\}^2]\\
&\hspace{2cm}=
\sup_{\mathbb P\in\mathcal D_{d,\epsilon}^{\mathrm{obs}}}
E_{\mathbb P}
[\{\widehat V-\Psi(\mathbb P)\}^2].
\end{aligned}
\]
For one inequality use
\(V^\star(P)=\Psi\{\operatorname{Obs}(P)\}\) and the observed marginal's membership in the observed class. For the other, complete each observed law to a causal law with that same marginal and value. Taking the estimator infimum preserves equality, so the observed bounds also prove the two stated causal minimax bounds.
% lean: CausalSmith.Stat.OptvalueVanishingoverlapRate.weighted_separation_and_frontier

\item
Finally, let the sequences in the statement be given. Define
\[
r_k=\frac{d_k}{n_k\epsilon_k\log(ed_k)},
\qquad
\rho_k=\min\{1,r_k\}.
\]
For the finitely many possible indices with \(n_k=0\), use the totalized division convention \(r_k=0\). Since \(n_k\to\infty\), one has \(n_k\ge1\) eventually. For all \(k\), \(r_k,\rho_k\ge0\), and
\[
\rho_k\longrightarrow0
\quad\Longleftrightarrow\quad
r_k\longrightarrow0.
\]
The reverse implication follows from \(0\le\rho_k\le r_k\).
For the forward implication, eventually \(\rho_k<1\); then
\(r_k\ge1\) would force \(\rho_k=1\), so eventually
\(r_k<1\) and \(\rho_k=r_k\).

Suppose a measurable estimator sequence \((\widehat V_k)\) has worst-case risks
\[
\mathcal W_k
=
\sup_{\mathbb P\in\mathcal D_{d_k,\epsilon_k}^{\mathrm{obs}}}
E_{\mathbb P}
[\{\widehat V_k-\Psi(\mathbb P)\}^2]
\longrightarrow0.
\]
The sample alphabet is finite. Thus each \(\widehat V_k\) is bounded on it, and these squared-error suprema are finite and nonnegative. The lower bound and the estimator infimum give, eventually,
\[
0\le\rho_k
\le\frac{\mathfrak R_{n_k,d_k,\epsilon_k}}{c_\star}
\le\frac{\mathcal W_k}{c_\star}.
\]
The squeeze theorem implies \(\rho_k\to0\), hence \(r_k\to0\).

Conversely, suppose \(r_k\to0\). Define a measurable estimator for every index by
\[
\widehat V_k=
\begin{cases}
\widehat V_{n_k,d_k,\epsilon_k}^{\mathrm{AF}},&n_k\ge1,\\
0,&n_k=0\text{ and }d_k<D_0,\\
1/2,&n_k=0\text{ and }d_k\ge D_0.
\end{cases}
\]
The zero-sample branches define constants on the singleton empty-sample space. For the eventual positive-sample indices, the established upper bound gives
\[
0\le
\sup_{\mathbb P\in\mathcal D_{d_k,\epsilon_k}^{\mathrm{obs}}}
E_{\mathbb P}
[\{\widehat V_k-\Psi(\mathbb P)\}^2]
\le C_\star\rho_k\longrightarrow0.
\]
A final application of the squeeze theorem proves uniform consistency and completes the equivalence.
% lean: CausalSmith.Stat.OptvalueVanishingoverlapRate.uniformConsistency_iff_of_frontier
\end{enumerate}
\end{proof}

\begin{proof}[Proof of \cref{obj:thm:all-estimator-lower}]
Choose the universal lower comparison constant \(c_\star>0\) supplied by the observed-risk conclusion of \cref{obj:thm:weighted-separation-and-frontier}. For arbitrary \(n,d\in\mathbb N\) and \(\epsilon\in\mathbb R\) satisfying \(n\ge1\), \(d\ge2\), and \(0<\epsilon\le1/2\), that conclusion gives
\[
c_\star\min\left\{1,\frac{d}{n\epsilon\log(ed)}\right\}
\le \mathfrak R_{n,d,\epsilon}.
\]
% lean: CausalSmith.Stat.OptvalueVanishingoverlapRate.weighted_separation_and_frontier
Since \(L_d=\log(ed)\), this is the required inequality. The constant \(c_\star\) is positive and independent of \(n,d,\epsilon\), as required.
% lean: CausalSmith.Stat.OptvalueVanishingoverlapRate.all_estimator_lower
\end{proof}

\bibliographystyle{plainnat}

\bibliography{references}

\end{document}
